% ============================================================================
%  Reporte de seguimiento - Trabajo Terminal II
%
%  Un solo archivo genera los reportes de cada mes. El mes se elige con \corte:
%      1 = Agosto 2026, 2 = Septiembre 2026
%
%  Compilación (desde esta carpeta):
%      pdflatex -jobname=1ReporteSeguimiento "\def\corte{1}% ============================================================================
%  Reporte de seguimiento - Trabajo Terminal II
%
%  Un solo archivo genera los reportes de cada mes. El mes se elige con \corte:
%      1 = Agosto 2026, 2 = Septiembre 2026
%
%  Compilación (desde esta carpeta):
%      pdflatex -jobname=1ReporteSeguimiento "\def\corte{1}% ============================================================================
%  Reporte de seguimiento - Trabajo Terminal II
%
%  Un solo archivo genera los reportes de cada mes. El mes se elige con \corte:
%      1 = Agosto 2026, 2 = Septiembre 2026
%
%  Compilación (desde esta carpeta):
%      pdflatex -jobname=1ReporteSeguimiento "\def\corte{1}% ============================================================================
%  Reporte de seguimiento - Trabajo Terminal II
%
%  Un solo archivo genera los reportes de cada mes. El mes se elige con \corte:
%      1 = Agosto 2026, 2 = Septiembre 2026
%
%  Compilación (desde esta carpeta):
%      pdflatex -jobname=1ReporteSeguimiento "\def\corte{1}\input{reporteSeguimiento}"
%      pdflatex -jobname=2ReporteSeguimiento "\def\corte{2}\input{reporteSeguimiento}"
%  o simplemente:  ./compilar.sh
%
%  Las páginas y secciones citadas corresponden a Documentation/main.pdf
%  (reporte técnico de TT), que es el PDF que se adjunta como evidencia.
% ============================================================================
\documentclass[letterpaper,10pt,oneside]{article}
\usepackage[margin=2.54cm]{geometry}
\usepackage[spanish]{babel}
\newcommand\authormark[1]{\textsuperscript{#1}}
\usepackage[colorlinks=true,bookmarks=true,linkcolor=black,citecolor=black,urlcolor=blue]{hyperref}
\usepackage{times}
\usepackage[table]{xcolor}
\usepackage{booktabs}
\usepackage{graphicx}
\usepackage{titlesec}
\usepackage{tabularx}
\usepackage{array}
\usepackage{enumitem}
\usepackage{float}

\titleformat{\section}
{\normalfont\fontsize{12pt}{14pt}\bfseries}
{\thesection}{1em}{}

% ----------------------------------------------------------------------------
%  Mes reportado
% ----------------------------------------------------------------------------
\providecommand\corte{2}
\ifnum\corte=1
	\def\numReporte{1}
	\def\mes{Agosto}
	\def\fechaCorte{31 de agosto de 2026}
\else
	\def\numReporte{2}
	\def\mes{Septiembre}
	\def\fechaCorte{30 de septiembre de 2026}
\fi
% \segun{texto agosto}{texto septiembre}
\newcommand\segun[2]{\ifnum\corte=1 #1\else #2\fi}

% ----------------------------------------------------------------------------
%  Porcentajes (editar aquí). Formato: \segun{agosto}{septiembre}
%  "esp" = avance esperado según cronograma al corte, "real" = avance logrado.
% ----------------------------------------------------------------------------
\def\espEntrenamiento{\segun{33}{67}}   \def\realEntrenamiento{\segun{33}{67}}
\def\espExperimentos{\segun{25}{50}}    \def\realExperimentos{\segun{25}{50}}
\def\espAnalisis{\segun{25}{50}}        \def\realAnalisis{\segun{25}{50}}

% ----------------------------------------------------------------------------
%  Enlaces a video (dejar vacío si no aplica)
% ----------------------------------------------------------------------------
\def\videoSimulacion{}
\def\videoEntrenamiento{}

% ----------------------------------------------------------------------------
%  Bloque de evidencia por actividad
%  \actividad{nombre}{fase}{periodo}{tipo de evidencia}{ubicación en reporte TT}
%            {descripción}{porcentaje}{enlace de video}
% ----------------------------------------------------------------------------
\definecolor{encabezado}{gray}{0.85}
\newcommand\video[1]{\edef\tmp{#1}\ifx\tmp\empty No aplica.\else\url{#1}\fi}
\newcommand\actividad[8]{%
	\par\noindent
	\begin{minipage}{\linewidth}
		\renewcommand{\arraystretch}{1.25}
		\begin{tabularx}{\linewidth}{|>{\raggedright\bfseries\arraybackslash}p{0.24\linewidth}|X|}
			\hline
			\rowcolor{encabezado}
			Nombre de la actividad & \textbf{#1} \\ \hline
			Fase y periodo & #2 \,(#3) \\ \hline
			Evidencia de trabajo & #4 \\ \hline
			Ubicación en el reporte de TT & #5 \\ \hline
			Descripción del avance & #6 \\ \hline
			Porcentaje de cumplimiento & \textbf{#7} \\ \hline
			Enlace de video & \video{#8} \\ \hline
		\end{tabularx}
	\end{minipage}
	\par\vspace{1.2em}%
}

% Porcentaje de cumplimiento: real respecto a lo esperado al corte
\newcommand\cumplimiento[2]{\fpeval{min(100, round(100*(#2)/(#1), 0))}\%}

% Columna del mes reportado en el cronograma
\newcommand\mesActual[2]{\ifnum\corte=#1\cellcolor{black!70}\textcolor{white}{\textbf{#2}}\else\textbf{#2}\fi}

\begin{document}

	\begin{center}
		\fontsize{14}{16}\selectfont\bfseries
		\numReporte° Reporte de seguimiento (\mes{} 2026): \\
		Simulación de sistemas complejos mediante autómatas celulares neuronales
	\end{center}

	\begin{center}
		\fontsize{14}{16}\selectfont\bfseries
		\textit{Trabajo Terminal No. 2026-A114}
	\end{center}

	\begin{center}
		\fontsize{12}{16}\selectfont
		\textit{Alumno: \authormark{*}Angeles López Erick Jesse\authormark{1}} \\
		\textit{Director: Juárez Martínez Genaro\authormark{1}} \\
		\textit{\authormark{*}e-mail: eangelesl1700@alumno.ipn.mx}
	\end{center}

	\section{Información general}

	\begin{itemize}[leftmargin=*]
		\item \textbf{Periodo reportado:} \mes{} 2026 (corte al \fechaCorte).
		\item \textbf{Documento de evidencia adjunto:} reporte técnico de Trabajo Terminal
		\textit{``Simulación de sistemas complejos mediante autómatas celulares neuronales''} (PDF, 204 páginas).
		El avance es incremental: cada evidencia listada en este reporte se encuentra integrada en el capítulo
		y sección correspondiente de dicho documento, cuyas páginas se indican en cada actividad.
		\item \textbf{Actividades programadas en el mes:} \segun{Módulos de entrenamiento (cierre), Entrenamiento y ajuste,
		Experimentación y análisis, Análisis de resultados (global).}{Entrenamiento y ajuste, Experimentación y análisis,
		Análisis de resultados (global).}
	\end{itemize}

	\section{Resumen de cumplimiento}

	\begin{table}[H]
		\centering
		\renewcommand{\arraystretch}{1.2}
		\begin{tabular}{l c c c c}
			\toprule
			\textbf{Actividad} & \textbf{Periodo} & \textbf{Esperado al corte} & \textbf{Avance real} & \textbf{Cumplimiento} \\
			\midrule
			\multicolumn{5}{l}{\textit{Fase 1: Planificación y Revisión Bibliográfica}} \\
			Estado del arte                  & Ago/25 -- Oct/25 & 100\% & 100\% & 100\% \\
			Marco teórico y metodológico     & Sep/25 -- Oct/25 & 100\% & 100\% & 100\% \\
			\multicolumn{5}{l}{\textit{Fase 2: Diseño e implementación del simulador}} \\
			Arquitectura del simulador       & Oct/25 -- Nov/25 & 100\% & 100\% & 100\% \\
			Módulos de simulación            & Nov/25 -- Dic/25 & 100\% & 100\% & 100\% \\
			Módulos entrenamiento            & Dic/25 -- Ago/26 & 100\% & 100\% & 100\% \\
			\multicolumn{5}{l}{\textit{Fase 3: Entrenamiento y Validación}} \\
			Entrenamiento y ajuste           & Ago/26 -- Oct/26 & \espEntrenamiento\% & \realEntrenamiento\% & \cumplimiento{\espEntrenamiento}{\realEntrenamiento} \\
			Experimentación y análisis       & Ago/26 -- Nov/26 & \espExperimentos\%  & \realExperimentos\%  & \cumplimiento{\espExperimentos}{\realExperimentos} \\
			Análisis de resultados (global)  & Ago/26 -- Nov/26 & \espAnalisis\%      & \realAnalisis\%      & \cumplimiento{\espAnalisis}{\realAnalisis} \\
			\multicolumn{5}{l}{\textit{Fase 4: Documentación y Difusión}} \\
			Documentación técnica            & Oct/26 -- Dic/26 & 0\% & -- & No iniciada \\
			Artículo académico               & Oct/26 -- Dic/26 & 0\% & -- & No iniciada \\
			\bottomrule
		\end{tabular}
		\caption{Avance acumulado por actividad al \fechaCorte. ``Esperado al corte'' es la proporción del periodo
		programado en el cronograma que ha transcurrido; ``Cumplimiento'' compara el avance real con dicho valor.}
	\end{table}

	\section{Evidencias por actividad}

	\subsection*{Fase 1: Planificación y Revisión Bibliográfica}

	\actividad{Estado del arte}
	{Fase 1}{Ago/25 -- Oct/25}
	{Documento escrito.}
	{Capítulo 5, sección 5.1 \textit{Estado del arte, aplicaciones y líneas actuales de investigación} (pp.~61--77);
	Capítulo 1 \textit{Introducción} (pp.~1--6).}
	{Revisión de trabajos sobre autómatas celulares neuronales: morfogénesis, NCA isotrópicos, autoclasificación,
	representaciones tridimensionales, NCA en malla, modelos generativos y aplicaciones adicionales.
	Actividad concluida en Trabajo Terminal I.}
	{100\%}{}

	\actividad{Marco teórico y metodológico}
	{Fase 1}{Sep/25 -- Oct/25}
	{Documento escrito.}
	{Capítulos 2 a 5 (marco teórico, pp.~7--77) y Capítulo 6 \textit{Marco metodológico} (pp.~78--90).}
	{Marco teórico de sistemas complejos y teoría del caos, autómatas celulares, inteligencia artificial y redes
	neuronales, y autómatas celulares neuronales. Marco metodológico con diseño cuasiexperimental, población,
	muestra, técnicas de recolección y procesamiento de datos, y contextualización tecnológica.
	Actividad concluida en Trabajo Terminal I.}
	{100\%}{}

	\subsection*{Fase 2: Diseño e implementación del simulador}

	\actividad{Arquitectura del simulador}
	{Fase 2}{Oct/25 -- Nov/25}
	{Documento escrito (análisis y diseño de software, diagramas UML).}
	{Capítulo 7 \textit{Análisis de software} (pp.~91--135); Capítulo 8, secciones 8.1 \textit{Arquitectura},
	8.2 \textit{Diagrama de clases}, 8.3 \textit{Diseño de datos} y 8.5 \textit{Interfaces de usuarios} (pp.~136--143, 153--156).}
	{Metodología de desarrollo, reglas de negocio, requerimientos, análisis de riesgos, estudio de factibilidad,
	casos de uso con sus especificaciones, diagramas de secuencia, arquitectura MVC, diagrama de clases,
	diseño de datos e interfaces gráficas.}
	{100\%}{}

	\actividad{Módulos de simulación}
	{Fase 2}{Nov/25 -- Dic/25}
	{Documento escrito, tablas de resultados y pruebas de operatividad.}
	{Sección 8.4 \textit{Diseño de algoritmos} (8.4.1--8.4.4 y 8.4.6, pp.~144--152); secciones 9.1 \textit{Entorno de desarrollo}
	y 9.2 \textit{Estructura del proyecto} (pp.~157--162); subsecciones 9.4.1, 9.4.3, 9.4.4 y 9.4.5 de
	\textit{Validación y pruebas} (pp.~166--169), tablas 9.6 y 9.7.}
	{Núcleo de simulación en C++/CUDA con enlaces a Python e interfaz en Qt: indexación de la regla semitotalística,
	paso de evolución, condiciones de frontera, \textit{forward pass} del NCA y \textit{prefetch} asíncrono.
	Validación funcional contra reglas conocidas y pruebas de rendimiento (hasta 345\,000 pasos/s en
	$100\times100$ y más de 100 fps de renderizado en $2000\times2000$).}
	{100\%}{\videoSimulacion}

	\actividad{Módulos entrenamiento}
	{Fase 2}{Dic/25 -- Ago/26}
	{Documento escrito, pruebas de operatividad.}
	{Subsección 7.8.4 \textit{Especificaciones de casos de uso: Entrenamiento} (pp.~125--130); 8.3.5
	\textit{Configuración del NCA} (p.~143); 8.4.5 \textit{Entrenamiento del autómata celular neuronal} (pp.~150--151);
	sección 9.3 \textit{Entrenamiento del modelo} (pp.~163--166, tabla 9.5); subsecciones 9.4.2 y 9.4.6 (pp.~167, 169).}
	{Módulo de entrenamiento en PyTorch con construcción dinámica del perceptrón multicapa, ciclo de optimización
	con Adam y detención temprana, persistencia de pesos en \texttt{.npz} compatible con el \textit{kernel} CUDA e
	integración con la pestaña de entrenamiento de la GUI (subproceso con \textit{log} en tiempo real).
	Se agregó el parámetro de semilla aleatoria para reproducir experimentos.
	\segun{Actividad concluida en el mes reportado.}{Actividad concluida en agosto de 2026.}}
	{100\%}{\videoEntrenamiento}

	\subsection*{Fase 3: Entrenamiento y Validación}

	\actividad{Entrenamiento y ajuste}
	{Fase 3}{Ago/26 -- Oct/26}
	{Documento escrito y tablas de resultados.}
	{Capítulo 10, sección 10.1 \textit{Duplicación de reglas del juego de la vida}:
	\segun{tablas 10.1 y 10.2 y figura 10.3 (pp.~173--174).}
	{tablas 10.1 y 10.2 y figura 10.3 (pp.~173--174); figuras 10.5 a 10.9 (pp.~176--178).}}
	{\segun{Barrido de hiperparámetros para la regla B3/S23: 10 combinaciones de funciones de activación con
	una capa oculta de 4 neuronas (una sola converge) y con 4 capas ocultas de 8 neuronas (cinco convergen a una
	pérdida menor a $10^{-4}$). Se identifica el estancamiento en las primeras 12 épocas.}
	{Además del barrido de hiperparámetros de agosto, se entrenaron las reglas B3/S23, B3/S238, B38/S23, B3/S236
	y B38/S238 con distintas pérdidas objetivo y semillas, midiendo la distancia Manhattan entre el vector de regla
	aprendido y el original. Se determinó que la pérdida objetivo debe ser $\leq 0.1$ para B3/S23.}}
	{\realEntrenamiento\% (esperado: \espEntrenamiento\%)}{}

	\actividad{Experimentación y análisis}
	{Fase 3}{Ago/26 -- Nov/26}
	{Documento escrito, tablas y gráficas de resultados.}
	{Capítulo 10, sección 10.1:
	\segun{figuras 10.1 y 10.2 (pp.~170--172).}
	{figuras 10.1 a 10.13 (pp.~170--181).}}
	{\segun{Generación del conjunto de datos (100 condiciones iniciales en $100\times100$ durante 100 pasos) y
	análisis de su balance: el estado 1 aparece un $16.5\% \pm 4.2\%$ y los patrones $C0N8$ y $C1N8$ representan
	menos del $0.02\%$ de las entradas.}
	{Primer experimento (duplicación de reglas del juego de la vida) completo. Se diseñó un espacio con las
	512 permutaciones de la vecindad para extraer el vector de regla aprendido y verificar si el NCA conserva la
	propiedad semitotalística. Se identificaron reglas aprendidas alternativas (B3/S23 en lugar de B3/S238,
	B3/S23678 en lugar de B3/S236) y reglas no semitotalísticas.}}
	{\realExperimentos\% (esperado: \espExperimentos\%)}{}

	\actividad{Análisis de resultados (global)}
	{Fase 3}{Ago/26 -- Nov/26}
	{Documento escrito.}
	{Capítulo 10, sección 10.1\segun{ (pp.~170--172).}{, conclusiones del experimento (p.~181).}}
	{\segun{Análisis estadístico de la distribución de estados y de patrones de vecindad para la regla B3/S23,
	identificando el sesgo de muestra y proponiendo sobremuestreo o ponderación de la función de costo.}
	{Conclusiones del primer experimento: la generalización depende del sesgo de muestra de patrones poco frecuentes,
	el entrenamiento de reglas semitotalísticas no garantiza una regla aprendida semitotalística, y no se requieren
	arquitecturas complejas para replicar estas dinámicas. Se valida la metodología de extracción del vector de
	regla para los siguientes experimentos.}}
	{\realAnalisis\% (esperado: \espAnalisis\%)}{}

	\subsection*{Fase 4: Documentación y Difusión}

	\actividad{Documentación técnica}
	{Fase 4}{Oct/26 -- Dic/26}
	{No aplica en el periodo.}
	{Reporte técnico de TT completo (204 páginas).}
	{Actividad programada a partir de octubre de 2026. El reporte técnico se ha redactado de forma incremental
	conforme avanza cada actividad, por lo que su integración final partirá del documento adjunto.}
	{No iniciada (0\% esperado al corte)}{}

	\actividad{Artículo académico}
	{Fase 4}{Oct/26 -- Dic/26}
	{No aplica en el periodo.}
	{--}
	{Actividad programada a partir de octubre de 2026.}
	{No iniciada (0\% esperado al corte)}{}

	\section{Cronograma}

	\resizebox{\linewidth}{!}{%
		\begin{tabular}{p{5cm}*{10}{>{\centering\arraybackslash}p{1cm}}}
			\toprule
			\textbf{Actividad} & \textbf{Ago/25} & \textbf{Sep/25} & \textbf{Oct/25} & \textbf{Nov/25} & \textbf{Dic/25} & \mesActual{1}{Ago/26} & \mesActual{2}{Sep/26} & \textbf{Oct/26} & \textbf{Nov/26} & \textbf{Dic/26} \\
			\midrule
			\multicolumn{11}{l}{\textbf{Fase 1: Planificación y Revisión Bibliográfica}} \\

			Estado del arte       & \cellcolor{gray!30} & \cellcolor{gray!30} & \cellcolor{gray!30} &  &  &  &  &  &  &   \\

			Marco teórico y metodológico &  & \cellcolor{gray!30} & \cellcolor{gray!30} &  &  &  &  &  &  &   \\

			\midrule

			\multicolumn{11}{l}{\textbf{Fase 2: Diseño e implementación del simulador}} \\
			Arquitectura del simulador &  &  & \cellcolor{gray!30} & \cellcolor{gray!30} &  &  &  &  &  &   \\

			Módulos de simulación  &  &  &  & \cellcolor{gray!30} & \cellcolor{gray!30} &  &  &  &  &   \\

			Módulos entrenamiento   &  &  &  &  & \cellcolor{gray!30} & \cellcolor{gray!30} &  &  &  &   \\

			\midrule

			\multicolumn{11}{l}{\textbf{Fase 3: Entrenamiento y Validación}} \\
			Entrenamiento y ajuste   &  &  &  &  &  & \cellcolor{gray!30} & \cellcolor{gray!30} & \cellcolor{gray!30} &   &   \\

			Experimentación y análisis  &  &  &  &  &  & \cellcolor{gray!30}  & \cellcolor{gray!30} & \cellcolor{gray!30} & \cellcolor{gray!30} &   \\

			Análisis de resultados (global) &  &  &  &  &  & \cellcolor{gray!30} & \cellcolor{gray!30} & \cellcolor{gray!30} & \cellcolor{gray!30} &   \\

			\midrule

			\multicolumn{11}{l}{\textbf{Fase 4: Documentación y Difusión}} \\
			Documentación técnica  &  &  &  &  &  &  &  & \cellcolor{gray!30}  & \cellcolor{gray!30} & \cellcolor{gray!30}  \\

			Artículo académico  &  &  &  &  &  &  &  & \cellcolor{gray!30}  & \cellcolor{gray!30}  & \cellcolor{gray!30}  \\

			\bottomrule
		\end{tabular}
	}

	\vspace{0.5em}
	{\small La columna resaltada corresponde al mes reportado.}

	\section{Alumnos y directores}

	\vspace{1em}
	\begin{tabular}{@{} p{0.45\textwidth} p{0.45\textwidth} @{}}
		\textit{Erick Jesse Angeles López} & \textit{Dr. Genaro Juárez Martínez} \\[5em]
		Firma: \rule{6cm}{0.4pt} & Firma: \rule{6cm}{0.4pt} \\
	\end{tabular}

\end{document}
"
%      pdflatex -jobname=2ReporteSeguimiento "\def\corte{2}% ============================================================================
%  Reporte de seguimiento - Trabajo Terminal II
%
%  Un solo archivo genera los reportes de cada mes. El mes se elige con \corte:
%      1 = Agosto 2026, 2 = Septiembre 2026
%
%  Compilación (desde esta carpeta):
%      pdflatex -jobname=1ReporteSeguimiento "\def\corte{1}\input{reporteSeguimiento}"
%      pdflatex -jobname=2ReporteSeguimiento "\def\corte{2}\input{reporteSeguimiento}"
%  o simplemente:  ./compilar.sh
%
%  Las páginas y secciones citadas corresponden a Documentation/main.pdf
%  (reporte técnico de TT), que es el PDF que se adjunta como evidencia.
% ============================================================================
\documentclass[letterpaper,10pt,oneside]{article}
\usepackage[margin=2.54cm]{geometry}
\usepackage[spanish]{babel}
\newcommand\authormark[1]{\textsuperscript{#1}}
\usepackage[colorlinks=true,bookmarks=true,linkcolor=black,citecolor=black,urlcolor=blue]{hyperref}
\usepackage{times}
\usepackage[table]{xcolor}
\usepackage{booktabs}
\usepackage{graphicx}
\usepackage{titlesec}
\usepackage{tabularx}
\usepackage{array}
\usepackage{enumitem}
\usepackage{float}

\titleformat{\section}
{\normalfont\fontsize{12pt}{14pt}\bfseries}
{\thesection}{1em}{}

% ----------------------------------------------------------------------------
%  Mes reportado
% ----------------------------------------------------------------------------
\providecommand\corte{2}
\ifnum\corte=1
	\def\numReporte{1}
	\def\mes{Agosto}
	\def\fechaCorte{31 de agosto de 2026}
\else
	\def\numReporte{2}
	\def\mes{Septiembre}
	\def\fechaCorte{30 de septiembre de 2026}
\fi
% \segun{texto agosto}{texto septiembre}
\newcommand\segun[2]{\ifnum\corte=1 #1\else #2\fi}

% ----------------------------------------------------------------------------
%  Porcentajes (editar aquí). Formato: \segun{agosto}{septiembre}
%  "esp" = avance esperado según cronograma al corte, "real" = avance logrado.
% ----------------------------------------------------------------------------
\def\espEntrenamiento{\segun{33}{67}}   \def\realEntrenamiento{\segun{33}{67}}
\def\espExperimentos{\segun{25}{50}}    \def\realExperimentos{\segun{25}{50}}
\def\espAnalisis{\segun{25}{50}}        \def\realAnalisis{\segun{25}{50}}

% ----------------------------------------------------------------------------
%  Enlaces a video (dejar vacío si no aplica)
% ----------------------------------------------------------------------------
\def\videoSimulacion{}
\def\videoEntrenamiento{}

% ----------------------------------------------------------------------------
%  Bloque de evidencia por actividad
%  \actividad{nombre}{fase}{periodo}{tipo de evidencia}{ubicación en reporte TT}
%            {descripción}{porcentaje}{enlace de video}
% ----------------------------------------------------------------------------
\definecolor{encabezado}{gray}{0.85}
\newcommand\video[1]{\edef\tmp{#1}\ifx\tmp\empty No aplica.\else\url{#1}\fi}
\newcommand\actividad[8]{%
	\par\noindent
	\begin{minipage}{\linewidth}
		\renewcommand{\arraystretch}{1.25}
		\begin{tabularx}{\linewidth}{|>{\raggedright\bfseries\arraybackslash}p{0.24\linewidth}|X|}
			\hline
			\rowcolor{encabezado}
			Nombre de la actividad & \textbf{#1} \\ \hline
			Fase y periodo & #2 \,(#3) \\ \hline
			Evidencia de trabajo & #4 \\ \hline
			Ubicación en el reporte de TT & #5 \\ \hline
			Descripción del avance & #6 \\ \hline
			Porcentaje de cumplimiento & \textbf{#7} \\ \hline
			Enlace de video & \video{#8} \\ \hline
		\end{tabularx}
	\end{minipage}
	\par\vspace{1.2em}%
}

% Porcentaje de cumplimiento: real respecto a lo esperado al corte
\newcommand\cumplimiento[2]{\fpeval{min(100, round(100*(#2)/(#1), 0))}\%}

% Columna del mes reportado en el cronograma
\newcommand\mesActual[2]{\ifnum\corte=#1\cellcolor{black!70}\textcolor{white}{\textbf{#2}}\else\textbf{#2}\fi}

\begin{document}

	\begin{center}
		\fontsize{14}{16}\selectfont\bfseries
		\numReporte° Reporte de seguimiento (\mes{} 2026): \\
		Simulación de sistemas complejos mediante autómatas celulares neuronales
	\end{center}

	\begin{center}
		\fontsize{14}{16}\selectfont\bfseries
		\textit{Trabajo Terminal No. 2026-A114}
	\end{center}

	\begin{center}
		\fontsize{12}{16}\selectfont
		\textit{Alumno: \authormark{*}Angeles López Erick Jesse\authormark{1}} \\
		\textit{Director: Juárez Martínez Genaro\authormark{1}} \\
		\textit{\authormark{*}e-mail: eangelesl1700@alumno.ipn.mx}
	\end{center}

	\section{Información general}

	\begin{itemize}[leftmargin=*]
		\item \textbf{Periodo reportado:} \mes{} 2026 (corte al \fechaCorte).
		\item \textbf{Documento de evidencia adjunto:} reporte técnico de Trabajo Terminal
		\textit{``Simulación de sistemas complejos mediante autómatas celulares neuronales''} (PDF, 204 páginas).
		El avance es incremental: cada evidencia listada en este reporte se encuentra integrada en el capítulo
		y sección correspondiente de dicho documento, cuyas páginas se indican en cada actividad.
		\item \textbf{Actividades programadas en el mes:} \segun{Módulos de entrenamiento (cierre), Entrenamiento y ajuste,
		Experimentación y análisis, Análisis de resultados (global).}{Entrenamiento y ajuste, Experimentación y análisis,
		Análisis de resultados (global).}
	\end{itemize}

	\section{Resumen de cumplimiento}

	\begin{table}[H]
		\centering
		\renewcommand{\arraystretch}{1.2}
		\begin{tabular}{l c c c c}
			\toprule
			\textbf{Actividad} & \textbf{Periodo} & \textbf{Esperado al corte} & \textbf{Avance real} & \textbf{Cumplimiento} \\
			\midrule
			\multicolumn{5}{l}{\textit{Fase 1: Planificación y Revisión Bibliográfica}} \\
			Estado del arte                  & Ago/25 -- Oct/25 & 100\% & 100\% & 100\% \\
			Marco teórico y metodológico     & Sep/25 -- Oct/25 & 100\% & 100\% & 100\% \\
			\multicolumn{5}{l}{\textit{Fase 2: Diseño e implementación del simulador}} \\
			Arquitectura del simulador       & Oct/25 -- Nov/25 & 100\% & 100\% & 100\% \\
			Módulos de simulación            & Nov/25 -- Dic/25 & 100\% & 100\% & 100\% \\
			Módulos entrenamiento            & Dic/25 -- Ago/26 & 100\% & 100\% & 100\% \\
			\multicolumn{5}{l}{\textit{Fase 3: Entrenamiento y Validación}} \\
			Entrenamiento y ajuste           & Ago/26 -- Oct/26 & \espEntrenamiento\% & \realEntrenamiento\% & \cumplimiento{\espEntrenamiento}{\realEntrenamiento} \\
			Experimentación y análisis       & Ago/26 -- Nov/26 & \espExperimentos\%  & \realExperimentos\%  & \cumplimiento{\espExperimentos}{\realExperimentos} \\
			Análisis de resultados (global)  & Ago/26 -- Nov/26 & \espAnalisis\%      & \realAnalisis\%      & \cumplimiento{\espAnalisis}{\realAnalisis} \\
			\multicolumn{5}{l}{\textit{Fase 4: Documentación y Difusión}} \\
			Documentación técnica            & Oct/26 -- Dic/26 & 0\% & -- & No iniciada \\
			Artículo académico               & Oct/26 -- Dic/26 & 0\% & -- & No iniciada \\
			\bottomrule
		\end{tabular}
		\caption{Avance acumulado por actividad al \fechaCorte. ``Esperado al corte'' es la proporción del periodo
		programado en el cronograma que ha transcurrido; ``Cumplimiento'' compara el avance real con dicho valor.}
	\end{table}

	\section{Evidencias por actividad}

	\subsection*{Fase 1: Planificación y Revisión Bibliográfica}

	\actividad{Estado del arte}
	{Fase 1}{Ago/25 -- Oct/25}
	{Documento escrito.}
	{Capítulo 5, sección 5.1 \textit{Estado del arte, aplicaciones y líneas actuales de investigación} (pp.~61--77);
	Capítulo 1 \textit{Introducción} (pp.~1--6).}
	{Revisión de trabajos sobre autómatas celulares neuronales: morfogénesis, NCA isotrópicos, autoclasificación,
	representaciones tridimensionales, NCA en malla, modelos generativos y aplicaciones adicionales.
	Actividad concluida en Trabajo Terminal I.}
	{100\%}{}

	\actividad{Marco teórico y metodológico}
	{Fase 1}{Sep/25 -- Oct/25}
	{Documento escrito.}
	{Capítulos 2 a 5 (marco teórico, pp.~7--77) y Capítulo 6 \textit{Marco metodológico} (pp.~78--90).}
	{Marco teórico de sistemas complejos y teoría del caos, autómatas celulares, inteligencia artificial y redes
	neuronales, y autómatas celulares neuronales. Marco metodológico con diseño cuasiexperimental, población,
	muestra, técnicas de recolección y procesamiento de datos, y contextualización tecnológica.
	Actividad concluida en Trabajo Terminal I.}
	{100\%}{}

	\subsection*{Fase 2: Diseño e implementación del simulador}

	\actividad{Arquitectura del simulador}
	{Fase 2}{Oct/25 -- Nov/25}
	{Documento escrito (análisis y diseño de software, diagramas UML).}
	{Capítulo 7 \textit{Análisis de software} (pp.~91--135); Capítulo 8, secciones 8.1 \textit{Arquitectura},
	8.2 \textit{Diagrama de clases}, 8.3 \textit{Diseño de datos} y 8.5 \textit{Interfaces de usuarios} (pp.~136--143, 153--156).}
	{Metodología de desarrollo, reglas de negocio, requerimientos, análisis de riesgos, estudio de factibilidad,
	casos de uso con sus especificaciones, diagramas de secuencia, arquitectura MVC, diagrama de clases,
	diseño de datos e interfaces gráficas.}
	{100\%}{}

	\actividad{Módulos de simulación}
	{Fase 2}{Nov/25 -- Dic/25}
	{Documento escrito, tablas de resultados y pruebas de operatividad.}
	{Sección 8.4 \textit{Diseño de algoritmos} (8.4.1--8.4.4 y 8.4.6, pp.~144--152); secciones 9.1 \textit{Entorno de desarrollo}
	y 9.2 \textit{Estructura del proyecto} (pp.~157--162); subsecciones 9.4.1, 9.4.3, 9.4.4 y 9.4.5 de
	\textit{Validación y pruebas} (pp.~166--169), tablas 9.6 y 9.7.}
	{Núcleo de simulación en C++/CUDA con enlaces a Python e interfaz en Qt: indexación de la regla semitotalística,
	paso de evolución, condiciones de frontera, \textit{forward pass} del NCA y \textit{prefetch} asíncrono.
	Validación funcional contra reglas conocidas y pruebas de rendimiento (hasta 345\,000 pasos/s en
	$100\times100$ y más de 100 fps de renderizado en $2000\times2000$).}
	{100\%}{\videoSimulacion}

	\actividad{Módulos entrenamiento}
	{Fase 2}{Dic/25 -- Ago/26}
	{Documento escrito, pruebas de operatividad.}
	{Subsección 7.8.4 \textit{Especificaciones de casos de uso: Entrenamiento} (pp.~125--130); 8.3.5
	\textit{Configuración del NCA} (p.~143); 8.4.5 \textit{Entrenamiento del autómata celular neuronal} (pp.~150--151);
	sección 9.3 \textit{Entrenamiento del modelo} (pp.~163--166, tabla 9.5); subsecciones 9.4.2 y 9.4.6 (pp.~167, 169).}
	{Módulo de entrenamiento en PyTorch con construcción dinámica del perceptrón multicapa, ciclo de optimización
	con Adam y detención temprana, persistencia de pesos en \texttt{.npz} compatible con el \textit{kernel} CUDA e
	integración con la pestaña de entrenamiento de la GUI (subproceso con \textit{log} en tiempo real).
	Se agregó el parámetro de semilla aleatoria para reproducir experimentos.
	\segun{Actividad concluida en el mes reportado.}{Actividad concluida en agosto de 2026.}}
	{100\%}{\videoEntrenamiento}

	\subsection*{Fase 3: Entrenamiento y Validación}

	\actividad{Entrenamiento y ajuste}
	{Fase 3}{Ago/26 -- Oct/26}
	{Documento escrito y tablas de resultados.}
	{Capítulo 10, sección 10.1 \textit{Duplicación de reglas del juego de la vida}:
	\segun{tablas 10.1 y 10.2 y figura 10.3 (pp.~173--174).}
	{tablas 10.1 y 10.2 y figura 10.3 (pp.~173--174); figuras 10.5 a 10.9 (pp.~176--178).}}
	{\segun{Barrido de hiperparámetros para la regla B3/S23: 10 combinaciones de funciones de activación con
	una capa oculta de 4 neuronas (una sola converge) y con 4 capas ocultas de 8 neuronas (cinco convergen a una
	pérdida menor a $10^{-4}$). Se identifica el estancamiento en las primeras 12 épocas.}
	{Además del barrido de hiperparámetros de agosto, se entrenaron las reglas B3/S23, B3/S238, B38/S23, B3/S236
	y B38/S238 con distintas pérdidas objetivo y semillas, midiendo la distancia Manhattan entre el vector de regla
	aprendido y el original. Se determinó que la pérdida objetivo debe ser $\leq 0.1$ para B3/S23.}}
	{\realEntrenamiento\% (esperado: \espEntrenamiento\%)}{}

	\actividad{Experimentación y análisis}
	{Fase 3}{Ago/26 -- Nov/26}
	{Documento escrito, tablas y gráficas de resultados.}
	{Capítulo 10, sección 10.1:
	\segun{figuras 10.1 y 10.2 (pp.~170--172).}
	{figuras 10.1 a 10.13 (pp.~170--181).}}
	{\segun{Generación del conjunto de datos (100 condiciones iniciales en $100\times100$ durante 100 pasos) y
	análisis de su balance: el estado 1 aparece un $16.5\% \pm 4.2\%$ y los patrones $C0N8$ y $C1N8$ representan
	menos del $0.02\%$ de las entradas.}
	{Primer experimento (duplicación de reglas del juego de la vida) completo. Se diseñó un espacio con las
	512 permutaciones de la vecindad para extraer el vector de regla aprendido y verificar si el NCA conserva la
	propiedad semitotalística. Se identificaron reglas aprendidas alternativas (B3/S23 en lugar de B3/S238,
	B3/S23678 en lugar de B3/S236) y reglas no semitotalísticas.}}
	{\realExperimentos\% (esperado: \espExperimentos\%)}{}

	\actividad{Análisis de resultados (global)}
	{Fase 3}{Ago/26 -- Nov/26}
	{Documento escrito.}
	{Capítulo 10, sección 10.1\segun{ (pp.~170--172).}{, conclusiones del experimento (p.~181).}}
	{\segun{Análisis estadístico de la distribución de estados y de patrones de vecindad para la regla B3/S23,
	identificando el sesgo de muestra y proponiendo sobremuestreo o ponderación de la función de costo.}
	{Conclusiones del primer experimento: la generalización depende del sesgo de muestra de patrones poco frecuentes,
	el entrenamiento de reglas semitotalísticas no garantiza una regla aprendida semitotalística, y no se requieren
	arquitecturas complejas para replicar estas dinámicas. Se valida la metodología de extracción del vector de
	regla para los siguientes experimentos.}}
	{\realAnalisis\% (esperado: \espAnalisis\%)}{}

	\subsection*{Fase 4: Documentación y Difusión}

	\actividad{Documentación técnica}
	{Fase 4}{Oct/26 -- Dic/26}
	{No aplica en el periodo.}
	{Reporte técnico de TT completo (204 páginas).}
	{Actividad programada a partir de octubre de 2026. El reporte técnico se ha redactado de forma incremental
	conforme avanza cada actividad, por lo que su integración final partirá del documento adjunto.}
	{No iniciada (0\% esperado al corte)}{}

	\actividad{Artículo académico}
	{Fase 4}{Oct/26 -- Dic/26}
	{No aplica en el periodo.}
	{--}
	{Actividad programada a partir de octubre de 2026.}
	{No iniciada (0\% esperado al corte)}{}

	\section{Cronograma}

	\resizebox{\linewidth}{!}{%
		\begin{tabular}{p{5cm}*{10}{>{\centering\arraybackslash}p{1cm}}}
			\toprule
			\textbf{Actividad} & \textbf{Ago/25} & \textbf{Sep/25} & \textbf{Oct/25} & \textbf{Nov/25} & \textbf{Dic/25} & \mesActual{1}{Ago/26} & \mesActual{2}{Sep/26} & \textbf{Oct/26} & \textbf{Nov/26} & \textbf{Dic/26} \\
			\midrule
			\multicolumn{11}{l}{\textbf{Fase 1: Planificación y Revisión Bibliográfica}} \\

			Estado del arte       & \cellcolor{gray!30} & \cellcolor{gray!30} & \cellcolor{gray!30} &  &  &  &  &  &  &   \\

			Marco teórico y metodológico &  & \cellcolor{gray!30} & \cellcolor{gray!30} &  &  &  &  &  &  &   \\

			\midrule

			\multicolumn{11}{l}{\textbf{Fase 2: Diseño e implementación del simulador}} \\
			Arquitectura del simulador &  &  & \cellcolor{gray!30} & \cellcolor{gray!30} &  &  &  &  &  &   \\

			Módulos de simulación  &  &  &  & \cellcolor{gray!30} & \cellcolor{gray!30} &  &  &  &  &   \\

			Módulos entrenamiento   &  &  &  &  & \cellcolor{gray!30} & \cellcolor{gray!30} &  &  &  &   \\

			\midrule

			\multicolumn{11}{l}{\textbf{Fase 3: Entrenamiento y Validación}} \\
			Entrenamiento y ajuste   &  &  &  &  &  & \cellcolor{gray!30} & \cellcolor{gray!30} & \cellcolor{gray!30} &   &   \\

			Experimentación y análisis  &  &  &  &  &  & \cellcolor{gray!30}  & \cellcolor{gray!30} & \cellcolor{gray!30} & \cellcolor{gray!30} &   \\

			Análisis de resultados (global) &  &  &  &  &  & \cellcolor{gray!30} & \cellcolor{gray!30} & \cellcolor{gray!30} & \cellcolor{gray!30} &   \\

			\midrule

			\multicolumn{11}{l}{\textbf{Fase 4: Documentación y Difusión}} \\
			Documentación técnica  &  &  &  &  &  &  &  & \cellcolor{gray!30}  & \cellcolor{gray!30} & \cellcolor{gray!30}  \\

			Artículo académico  &  &  &  &  &  &  &  & \cellcolor{gray!30}  & \cellcolor{gray!30}  & \cellcolor{gray!30}  \\

			\bottomrule
		\end{tabular}
	}

	\vspace{0.5em}
	{\small La columna resaltada corresponde al mes reportado.}

	\section{Alumnos y directores}

	\vspace{1em}
	\begin{tabular}{@{} p{0.45\textwidth} p{0.45\textwidth} @{}}
		\textit{Erick Jesse Angeles López} & \textit{Dr. Genaro Juárez Martínez} \\[5em]
		Firma: \rule{6cm}{0.4pt} & Firma: \rule{6cm}{0.4pt} \\
	\end{tabular}

\end{document}
"
%  o simplemente:  ./compilar.sh
%
%  Las páginas y secciones citadas corresponden a Documentation/main.pdf
%  (reporte técnico de TT), que es el PDF que se adjunta como evidencia.
% ============================================================================
\documentclass[letterpaper,10pt,oneside]{article}
\usepackage[margin=2.54cm]{geometry}
\usepackage[spanish]{babel}
\newcommand\authormark[1]{\textsuperscript{#1}}
\usepackage[colorlinks=true,bookmarks=true,linkcolor=black,citecolor=black,urlcolor=blue]{hyperref}
\usepackage{times}
\usepackage[table]{xcolor}
\usepackage{booktabs}
\usepackage{graphicx}
\usepackage{titlesec}
\usepackage{tabularx}
\usepackage{array}
\usepackage{enumitem}
\usepackage{float}

\titleformat{\section}
{\normalfont\fontsize{12pt}{14pt}\bfseries}
{\thesection}{1em}{}

% ----------------------------------------------------------------------------
%  Mes reportado
% ----------------------------------------------------------------------------
\providecommand\corte{2}
\ifnum\corte=1
	\def\numReporte{1}
	\def\mes{Agosto}
	\def\fechaCorte{31 de agosto de 2026}
\else
	\def\numReporte{2}
	\def\mes{Septiembre}
	\def\fechaCorte{30 de septiembre de 2026}
\fi
% \segun{texto agosto}{texto septiembre}
\newcommand\segun[2]{\ifnum\corte=1 #1\else #2\fi}

% ----------------------------------------------------------------------------
%  Porcentajes (editar aquí). Formato: \segun{agosto}{septiembre}
%  "esp" = avance esperado según cronograma al corte, "real" = avance logrado.
% ----------------------------------------------------------------------------
\def\espEntrenamiento{\segun{33}{67}}   \def\realEntrenamiento{\segun{33}{67}}
\def\espExperimentos{\segun{25}{50}}    \def\realExperimentos{\segun{25}{50}}
\def\espAnalisis{\segun{25}{50}}        \def\realAnalisis{\segun{25}{50}}

% ----------------------------------------------------------------------------
%  Enlaces a video (dejar vacío si no aplica)
% ----------------------------------------------------------------------------
\def\videoSimulacion{}
\def\videoEntrenamiento{}

% ----------------------------------------------------------------------------
%  Bloque de evidencia por actividad
%  \actividad{nombre}{fase}{periodo}{tipo de evidencia}{ubicación en reporte TT}
%            {descripción}{porcentaje}{enlace de video}
% ----------------------------------------------------------------------------
\definecolor{encabezado}{gray}{0.85}
\newcommand\video[1]{\edef\tmp{#1}\ifx\tmp\empty No aplica.\else\url{#1}\fi}
\newcommand\actividad[8]{%
	\par\noindent
	\begin{minipage}{\linewidth}
		\renewcommand{\arraystretch}{1.25}
		\begin{tabularx}{\linewidth}{|>{\raggedright\bfseries\arraybackslash}p{0.24\linewidth}|X|}
			\hline
			\rowcolor{encabezado}
			Nombre de la actividad & \textbf{#1} \\ \hline
			Fase y periodo & #2 \,(#3) \\ \hline
			Evidencia de trabajo & #4 \\ \hline
			Ubicación en el reporte de TT & #5 \\ \hline
			Descripción del avance & #6 \\ \hline
			Porcentaje de cumplimiento & \textbf{#7} \\ \hline
			Enlace de video & \video{#8} \\ \hline
		\end{tabularx}
	\end{minipage}
	\par\vspace{1.2em}%
}

% Porcentaje de cumplimiento: real respecto a lo esperado al corte
\newcommand\cumplimiento[2]{\fpeval{min(100, round(100*(#2)/(#1), 0))}\%}

% Columna del mes reportado en el cronograma
\newcommand\mesActual[2]{\ifnum\corte=#1\cellcolor{black!70}\textcolor{white}{\textbf{#2}}\else\textbf{#2}\fi}

\begin{document}

	\begin{center}
		\fontsize{14}{16}\selectfont\bfseries
		\numReporte° Reporte de seguimiento (\mes{} 2026): \\
		Simulación de sistemas complejos mediante autómatas celulares neuronales
	\end{center}

	\begin{center}
		\fontsize{14}{16}\selectfont\bfseries
		\textit{Trabajo Terminal No. 2026-A114}
	\end{center}

	\begin{center}
		\fontsize{12}{16}\selectfont
		\textit{Alumno: \authormark{*}Angeles López Erick Jesse\authormark{1}} \\
		\textit{Director: Juárez Martínez Genaro\authormark{1}} \\
		\textit{\authormark{*}e-mail: eangelesl1700@alumno.ipn.mx}
	\end{center}

	\section{Información general}

	\begin{itemize}[leftmargin=*]
		\item \textbf{Periodo reportado:} \mes{} 2026 (corte al \fechaCorte).
		\item \textbf{Documento de evidencia adjunto:} reporte técnico de Trabajo Terminal
		\textit{``Simulación de sistemas complejos mediante autómatas celulares neuronales''} (PDF, 204 páginas).
		El avance es incremental: cada evidencia listada en este reporte se encuentra integrada en el capítulo
		y sección correspondiente de dicho documento, cuyas páginas se indican en cada actividad.
		\item \textbf{Actividades programadas en el mes:} \segun{Módulos de entrenamiento (cierre), Entrenamiento y ajuste,
		Experimentación y análisis, Análisis de resultados (global).}{Entrenamiento y ajuste, Experimentación y análisis,
		Análisis de resultados (global).}
	\end{itemize}

	\section{Resumen de cumplimiento}

	\begin{table}[H]
		\centering
		\renewcommand{\arraystretch}{1.2}
		\begin{tabular}{l c c c c}
			\toprule
			\textbf{Actividad} & \textbf{Periodo} & \textbf{Esperado al corte} & \textbf{Avance real} & \textbf{Cumplimiento} \\
			\midrule
			\multicolumn{5}{l}{\textit{Fase 1: Planificación y Revisión Bibliográfica}} \\
			Estado del arte                  & Ago/25 -- Oct/25 & 100\% & 100\% & 100\% \\
			Marco teórico y metodológico     & Sep/25 -- Oct/25 & 100\% & 100\% & 100\% \\
			\multicolumn{5}{l}{\textit{Fase 2: Diseño e implementación del simulador}} \\
			Arquitectura del simulador       & Oct/25 -- Nov/25 & 100\% & 100\% & 100\% \\
			Módulos de simulación            & Nov/25 -- Dic/25 & 100\% & 100\% & 100\% \\
			Módulos entrenamiento            & Dic/25 -- Ago/26 & 100\% & 100\% & 100\% \\
			\multicolumn{5}{l}{\textit{Fase 3: Entrenamiento y Validación}} \\
			Entrenamiento y ajuste           & Ago/26 -- Oct/26 & \espEntrenamiento\% & \realEntrenamiento\% & \cumplimiento{\espEntrenamiento}{\realEntrenamiento} \\
			Experimentación y análisis       & Ago/26 -- Nov/26 & \espExperimentos\%  & \realExperimentos\%  & \cumplimiento{\espExperimentos}{\realExperimentos} \\
			Análisis de resultados (global)  & Ago/26 -- Nov/26 & \espAnalisis\%      & \realAnalisis\%      & \cumplimiento{\espAnalisis}{\realAnalisis} \\
			\multicolumn{5}{l}{\textit{Fase 4: Documentación y Difusión}} \\
			Documentación técnica            & Oct/26 -- Dic/26 & 0\% & -- & No iniciada \\
			Artículo académico               & Oct/26 -- Dic/26 & 0\% & -- & No iniciada \\
			\bottomrule
		\end{tabular}
		\caption{Avance acumulado por actividad al \fechaCorte. ``Esperado al corte'' es la proporción del periodo
		programado en el cronograma que ha transcurrido; ``Cumplimiento'' compara el avance real con dicho valor.}
	\end{table}

	\section{Evidencias por actividad}

	\subsection*{Fase 1: Planificación y Revisión Bibliográfica}

	\actividad{Estado del arte}
	{Fase 1}{Ago/25 -- Oct/25}
	{Documento escrito.}
	{Capítulo 5, sección 5.1 \textit{Estado del arte, aplicaciones y líneas actuales de investigación} (pp.~61--77);
	Capítulo 1 \textit{Introducción} (pp.~1--6).}
	{Revisión de trabajos sobre autómatas celulares neuronales: morfogénesis, NCA isotrópicos, autoclasificación,
	representaciones tridimensionales, NCA en malla, modelos generativos y aplicaciones adicionales.
	Actividad concluida en Trabajo Terminal I.}
	{100\%}{}

	\actividad{Marco teórico y metodológico}
	{Fase 1}{Sep/25 -- Oct/25}
	{Documento escrito.}
	{Capítulos 2 a 5 (marco teórico, pp.~7--77) y Capítulo 6 \textit{Marco metodológico} (pp.~78--90).}
	{Marco teórico de sistemas complejos y teoría del caos, autómatas celulares, inteligencia artificial y redes
	neuronales, y autómatas celulares neuronales. Marco metodológico con diseño cuasiexperimental, población,
	muestra, técnicas de recolección y procesamiento de datos, y contextualización tecnológica.
	Actividad concluida en Trabajo Terminal I.}
	{100\%}{}

	\subsection*{Fase 2: Diseño e implementación del simulador}

	\actividad{Arquitectura del simulador}
	{Fase 2}{Oct/25 -- Nov/25}
	{Documento escrito (análisis y diseño de software, diagramas UML).}
	{Capítulo 7 \textit{Análisis de software} (pp.~91--135); Capítulo 8, secciones 8.1 \textit{Arquitectura},
	8.2 \textit{Diagrama de clases}, 8.3 \textit{Diseño de datos} y 8.5 \textit{Interfaces de usuarios} (pp.~136--143, 153--156).}
	{Metodología de desarrollo, reglas de negocio, requerimientos, análisis de riesgos, estudio de factibilidad,
	casos de uso con sus especificaciones, diagramas de secuencia, arquitectura MVC, diagrama de clases,
	diseño de datos e interfaces gráficas.}
	{100\%}{}

	\actividad{Módulos de simulación}
	{Fase 2}{Nov/25 -- Dic/25}
	{Documento escrito, tablas de resultados y pruebas de operatividad.}
	{Sección 8.4 \textit{Diseño de algoritmos} (8.4.1--8.4.4 y 8.4.6, pp.~144--152); secciones 9.1 \textit{Entorno de desarrollo}
	y 9.2 \textit{Estructura del proyecto} (pp.~157--162); subsecciones 9.4.1, 9.4.3, 9.4.4 y 9.4.5 de
	\textit{Validación y pruebas} (pp.~166--169), tablas 9.6 y 9.7.}
	{Núcleo de simulación en C++/CUDA con enlaces a Python e interfaz en Qt: indexación de la regla semitotalística,
	paso de evolución, condiciones de frontera, \textit{forward pass} del NCA y \textit{prefetch} asíncrono.
	Validación funcional contra reglas conocidas y pruebas de rendimiento (hasta 345\,000 pasos/s en
	$100\times100$ y más de 100 fps de renderizado en $2000\times2000$).}
	{100\%}{\videoSimulacion}

	\actividad{Módulos entrenamiento}
	{Fase 2}{Dic/25 -- Ago/26}
	{Documento escrito, pruebas de operatividad.}
	{Subsección 7.8.4 \textit{Especificaciones de casos de uso: Entrenamiento} (pp.~125--130); 8.3.5
	\textit{Configuración del NCA} (p.~143); 8.4.5 \textit{Entrenamiento del autómata celular neuronal} (pp.~150--151);
	sección 9.3 \textit{Entrenamiento del modelo} (pp.~163--166, tabla 9.5); subsecciones 9.4.2 y 9.4.6 (pp.~167, 169).}
	{Módulo de entrenamiento en PyTorch con construcción dinámica del perceptrón multicapa, ciclo de optimización
	con Adam y detención temprana, persistencia de pesos en \texttt{.npz} compatible con el \textit{kernel} CUDA e
	integración con la pestaña de entrenamiento de la GUI (subproceso con \textit{log} en tiempo real).
	Se agregó el parámetro de semilla aleatoria para reproducir experimentos.
	\segun{Actividad concluida en el mes reportado.}{Actividad concluida en agosto de 2026.}}
	{100\%}{\videoEntrenamiento}

	\subsection*{Fase 3: Entrenamiento y Validación}

	\actividad{Entrenamiento y ajuste}
	{Fase 3}{Ago/26 -- Oct/26}
	{Documento escrito y tablas de resultados.}
	{Capítulo 10, sección 10.1 \textit{Duplicación de reglas del juego de la vida}:
	\segun{tablas 10.1 y 10.2 y figura 10.3 (pp.~173--174).}
	{tablas 10.1 y 10.2 y figura 10.3 (pp.~173--174); figuras 10.5 a 10.9 (pp.~176--178).}}
	{\segun{Barrido de hiperparámetros para la regla B3/S23: 10 combinaciones de funciones de activación con
	una capa oculta de 4 neuronas (una sola converge) y con 4 capas ocultas de 8 neuronas (cinco convergen a una
	pérdida menor a $10^{-4}$). Se identifica el estancamiento en las primeras 12 épocas.}
	{Además del barrido de hiperparámetros de agosto, se entrenaron las reglas B3/S23, B3/S238, B38/S23, B3/S236
	y B38/S238 con distintas pérdidas objetivo y semillas, midiendo la distancia Manhattan entre el vector de regla
	aprendido y el original. Se determinó que la pérdida objetivo debe ser $\leq 0.1$ para B3/S23.}}
	{\realEntrenamiento\% (esperado: \espEntrenamiento\%)}{}

	\actividad{Experimentación y análisis}
	{Fase 3}{Ago/26 -- Nov/26}
	{Documento escrito, tablas y gráficas de resultados.}
	{Capítulo 10, sección 10.1:
	\segun{figuras 10.1 y 10.2 (pp.~170--172).}
	{figuras 10.1 a 10.13 (pp.~170--181).}}
	{\segun{Generación del conjunto de datos (100 condiciones iniciales en $100\times100$ durante 100 pasos) y
	análisis de su balance: el estado 1 aparece un $16.5\% \pm 4.2\%$ y los patrones $C0N8$ y $C1N8$ representan
	menos del $0.02\%$ de las entradas.}
	{Primer experimento (duplicación de reglas del juego de la vida) completo. Se diseñó un espacio con las
	512 permutaciones de la vecindad para extraer el vector de regla aprendido y verificar si el NCA conserva la
	propiedad semitotalística. Se identificaron reglas aprendidas alternativas (B3/S23 en lugar de B3/S238,
	B3/S23678 en lugar de B3/S236) y reglas no semitotalísticas.}}
	{\realExperimentos\% (esperado: \espExperimentos\%)}{}

	\actividad{Análisis de resultados (global)}
	{Fase 3}{Ago/26 -- Nov/26}
	{Documento escrito.}
	{Capítulo 10, sección 10.1\segun{ (pp.~170--172).}{, conclusiones del experimento (p.~181).}}
	{\segun{Análisis estadístico de la distribución de estados y de patrones de vecindad para la regla B3/S23,
	identificando el sesgo de muestra y proponiendo sobremuestreo o ponderación de la función de costo.}
	{Conclusiones del primer experimento: la generalización depende del sesgo de muestra de patrones poco frecuentes,
	el entrenamiento de reglas semitotalísticas no garantiza una regla aprendida semitotalística, y no se requieren
	arquitecturas complejas para replicar estas dinámicas. Se valida la metodología de extracción del vector de
	regla para los siguientes experimentos.}}
	{\realAnalisis\% (esperado: \espAnalisis\%)}{}

	\subsection*{Fase 4: Documentación y Difusión}

	\actividad{Documentación técnica}
	{Fase 4}{Oct/26 -- Dic/26}
	{No aplica en el periodo.}
	{Reporte técnico de TT completo (204 páginas).}
	{Actividad programada a partir de octubre de 2026. El reporte técnico se ha redactado de forma incremental
	conforme avanza cada actividad, por lo que su integración final partirá del documento adjunto.}
	{No iniciada (0\% esperado al corte)}{}

	\actividad{Artículo académico}
	{Fase 4}{Oct/26 -- Dic/26}
	{No aplica en el periodo.}
	{--}
	{Actividad programada a partir de octubre de 2026.}
	{No iniciada (0\% esperado al corte)}{}

	\section{Cronograma}

	\resizebox{\linewidth}{!}{%
		\begin{tabular}{p{5cm}*{10}{>{\centering\arraybackslash}p{1cm}}}
			\toprule
			\textbf{Actividad} & \textbf{Ago/25} & \textbf{Sep/25} & \textbf{Oct/25} & \textbf{Nov/25} & \textbf{Dic/25} & \mesActual{1}{Ago/26} & \mesActual{2}{Sep/26} & \textbf{Oct/26} & \textbf{Nov/26} & \textbf{Dic/26} \\
			\midrule
			\multicolumn{11}{l}{\textbf{Fase 1: Planificación y Revisión Bibliográfica}} \\

			Estado del arte       & \cellcolor{gray!30} & \cellcolor{gray!30} & \cellcolor{gray!30} &  &  &  &  &  &  &   \\

			Marco teórico y metodológico &  & \cellcolor{gray!30} & \cellcolor{gray!30} &  &  &  &  &  &  &   \\

			\midrule

			\multicolumn{11}{l}{\textbf{Fase 2: Diseño e implementación del simulador}} \\
			Arquitectura del simulador &  &  & \cellcolor{gray!30} & \cellcolor{gray!30} &  &  &  &  &  &   \\

			Módulos de simulación  &  &  &  & \cellcolor{gray!30} & \cellcolor{gray!30} &  &  &  &  &   \\

			Módulos entrenamiento   &  &  &  &  & \cellcolor{gray!30} & \cellcolor{gray!30} &  &  &  &   \\

			\midrule

			\multicolumn{11}{l}{\textbf{Fase 3: Entrenamiento y Validación}} \\
			Entrenamiento y ajuste   &  &  &  &  &  & \cellcolor{gray!30} & \cellcolor{gray!30} & \cellcolor{gray!30} &   &   \\

			Experimentación y análisis  &  &  &  &  &  & \cellcolor{gray!30}  & \cellcolor{gray!30} & \cellcolor{gray!30} & \cellcolor{gray!30} &   \\

			Análisis de resultados (global) &  &  &  &  &  & \cellcolor{gray!30} & \cellcolor{gray!30} & \cellcolor{gray!30} & \cellcolor{gray!30} &   \\

			\midrule

			\multicolumn{11}{l}{\textbf{Fase 4: Documentación y Difusión}} \\
			Documentación técnica  &  &  &  &  &  &  &  & \cellcolor{gray!30}  & \cellcolor{gray!30} & \cellcolor{gray!30}  \\

			Artículo académico  &  &  &  &  &  &  &  & \cellcolor{gray!30}  & \cellcolor{gray!30}  & \cellcolor{gray!30}  \\

			\bottomrule
		\end{tabular}
	}

	\vspace{0.5em}
	{\small La columna resaltada corresponde al mes reportado.}

	\section{Alumnos y directores}

	\vspace{1em}
	\begin{tabular}{@{} p{0.45\textwidth} p{0.45\textwidth} @{}}
		\textit{Erick Jesse Angeles López} & \textit{Dr. Genaro Juárez Martínez} \\[5em]
		Firma: \rule{6cm}{0.4pt} & Firma: \rule{6cm}{0.4pt} \\
	\end{tabular}

\end{document}
"
%      pdflatex -jobname=2ReporteSeguimiento "\def\corte{2}% ============================================================================
%  Reporte de seguimiento - Trabajo Terminal II
%
%  Un solo archivo genera los reportes de cada mes. El mes se elige con \corte:
%      1 = Agosto 2026, 2 = Septiembre 2026
%
%  Compilación (desde esta carpeta):
%      pdflatex -jobname=1ReporteSeguimiento "\def\corte{1}% ============================================================================
%  Reporte de seguimiento - Trabajo Terminal II
%
%  Un solo archivo genera los reportes de cada mes. El mes se elige con \corte:
%      1 = Agosto 2026, 2 = Septiembre 2026
%
%  Compilación (desde esta carpeta):
%      pdflatex -jobname=1ReporteSeguimiento "\def\corte{1}\input{reporteSeguimiento}"
%      pdflatex -jobname=2ReporteSeguimiento "\def\corte{2}\input{reporteSeguimiento}"
%  o simplemente:  ./compilar.sh
%
%  Las páginas y secciones citadas corresponden a Documentation/main.pdf
%  (reporte técnico de TT), que es el PDF que se adjunta como evidencia.
% ============================================================================
\documentclass[letterpaper,10pt,oneside]{article}
\usepackage[margin=2.54cm]{geometry}
\usepackage[spanish]{babel}
\newcommand\authormark[1]{\textsuperscript{#1}}
\usepackage[colorlinks=true,bookmarks=true,linkcolor=black,citecolor=black,urlcolor=blue]{hyperref}
\usepackage{times}
\usepackage[table]{xcolor}
\usepackage{booktabs}
\usepackage{graphicx}
\usepackage{titlesec}
\usepackage{tabularx}
\usepackage{array}
\usepackage{enumitem}
\usepackage{float}

\titleformat{\section}
{\normalfont\fontsize{12pt}{14pt}\bfseries}
{\thesection}{1em}{}

% ----------------------------------------------------------------------------
%  Mes reportado
% ----------------------------------------------------------------------------
\providecommand\corte{2}
\ifnum\corte=1
	\def\numReporte{1}
	\def\mes{Agosto}
	\def\fechaCorte{31 de agosto de 2026}
\else
	\def\numReporte{2}
	\def\mes{Septiembre}
	\def\fechaCorte{30 de septiembre de 2026}
\fi
% \segun{texto agosto}{texto septiembre}
\newcommand\segun[2]{\ifnum\corte=1 #1\else #2\fi}

% ----------------------------------------------------------------------------
%  Porcentajes (editar aquí). Formato: \segun{agosto}{septiembre}
%  "esp" = avance esperado según cronograma al corte, "real" = avance logrado.
% ----------------------------------------------------------------------------
\def\espEntrenamiento{\segun{33}{67}}   \def\realEntrenamiento{\segun{33}{67}}
\def\espExperimentos{\segun{25}{50}}    \def\realExperimentos{\segun{25}{50}}
\def\espAnalisis{\segun{25}{50}}        \def\realAnalisis{\segun{25}{50}}

% ----------------------------------------------------------------------------
%  Enlaces a video (dejar vacío si no aplica)
% ----------------------------------------------------------------------------
\def\videoSimulacion{}
\def\videoEntrenamiento{}

% ----------------------------------------------------------------------------
%  Bloque de evidencia por actividad
%  \actividad{nombre}{fase}{periodo}{tipo de evidencia}{ubicación en reporte TT}
%            {descripción}{porcentaje}{enlace de video}
% ----------------------------------------------------------------------------
\definecolor{encabezado}{gray}{0.85}
\newcommand\video[1]{\edef\tmp{#1}\ifx\tmp\empty No aplica.\else\url{#1}\fi}
\newcommand\actividad[8]{%
	\par\noindent
	\begin{minipage}{\linewidth}
		\renewcommand{\arraystretch}{1.25}
		\begin{tabularx}{\linewidth}{|>{\raggedright\bfseries\arraybackslash}p{0.24\linewidth}|X|}
			\hline
			\rowcolor{encabezado}
			Nombre de la actividad & \textbf{#1} \\ \hline
			Fase y periodo & #2 \,(#3) \\ \hline
			Evidencia de trabajo & #4 \\ \hline
			Ubicación en el reporte de TT & #5 \\ \hline
			Descripción del avance & #6 \\ \hline
			Porcentaje de cumplimiento & \textbf{#7} \\ \hline
			Enlace de video & \video{#8} \\ \hline
		\end{tabularx}
	\end{minipage}
	\par\vspace{1.2em}%
}

% Porcentaje de cumplimiento: real respecto a lo esperado al corte
\newcommand\cumplimiento[2]{\fpeval{min(100, round(100*(#2)/(#1), 0))}\%}

% Columna del mes reportado en el cronograma
\newcommand\mesActual[2]{\ifnum\corte=#1\cellcolor{black!70}\textcolor{white}{\textbf{#2}}\else\textbf{#2}\fi}

\begin{document}

	\begin{center}
		\fontsize{14}{16}\selectfont\bfseries
		\numReporte° Reporte de seguimiento (\mes{} 2026): \\
		Simulación de sistemas complejos mediante autómatas celulares neuronales
	\end{center}

	\begin{center}
		\fontsize{14}{16}\selectfont\bfseries
		\textit{Trabajo Terminal No. 2026-A114}
	\end{center}

	\begin{center}
		\fontsize{12}{16}\selectfont
		\textit{Alumno: \authormark{*}Angeles López Erick Jesse\authormark{1}} \\
		\textit{Director: Juárez Martínez Genaro\authormark{1}} \\
		\textit{\authormark{*}e-mail: eangelesl1700@alumno.ipn.mx}
	\end{center}

	\section{Información general}

	\begin{itemize}[leftmargin=*]
		\item \textbf{Periodo reportado:} \mes{} 2026 (corte al \fechaCorte).
		\item \textbf{Documento de evidencia adjunto:} reporte técnico de Trabajo Terminal
		\textit{``Simulación de sistemas complejos mediante autómatas celulares neuronales''} (PDF, 204 páginas).
		El avance es incremental: cada evidencia listada en este reporte se encuentra integrada en el capítulo
		y sección correspondiente de dicho documento, cuyas páginas se indican en cada actividad.
		\item \textbf{Actividades programadas en el mes:} \segun{Módulos de entrenamiento (cierre), Entrenamiento y ajuste,
		Experimentación y análisis, Análisis de resultados (global).}{Entrenamiento y ajuste, Experimentación y análisis,
		Análisis de resultados (global).}
	\end{itemize}

	\section{Resumen de cumplimiento}

	\begin{table}[H]
		\centering
		\renewcommand{\arraystretch}{1.2}
		\begin{tabular}{l c c c c}
			\toprule
			\textbf{Actividad} & \textbf{Periodo} & \textbf{Esperado al corte} & \textbf{Avance real} & \textbf{Cumplimiento} \\
			\midrule
			\multicolumn{5}{l}{\textit{Fase 1: Planificación y Revisión Bibliográfica}} \\
			Estado del arte                  & Ago/25 -- Oct/25 & 100\% & 100\% & 100\% \\
			Marco teórico y metodológico     & Sep/25 -- Oct/25 & 100\% & 100\% & 100\% \\
			\multicolumn{5}{l}{\textit{Fase 2: Diseño e implementación del simulador}} \\
			Arquitectura del simulador       & Oct/25 -- Nov/25 & 100\% & 100\% & 100\% \\
			Módulos de simulación            & Nov/25 -- Dic/25 & 100\% & 100\% & 100\% \\
			Módulos entrenamiento            & Dic/25 -- Ago/26 & 100\% & 100\% & 100\% \\
			\multicolumn{5}{l}{\textit{Fase 3: Entrenamiento y Validación}} \\
			Entrenamiento y ajuste           & Ago/26 -- Oct/26 & \espEntrenamiento\% & \realEntrenamiento\% & \cumplimiento{\espEntrenamiento}{\realEntrenamiento} \\
			Experimentación y análisis       & Ago/26 -- Nov/26 & \espExperimentos\%  & \realExperimentos\%  & \cumplimiento{\espExperimentos}{\realExperimentos} \\
			Análisis de resultados (global)  & Ago/26 -- Nov/26 & \espAnalisis\%      & \realAnalisis\%      & \cumplimiento{\espAnalisis}{\realAnalisis} \\
			\multicolumn{5}{l}{\textit{Fase 4: Documentación y Difusión}} \\
			Documentación técnica            & Oct/26 -- Dic/26 & 0\% & -- & No iniciada \\
			Artículo académico               & Oct/26 -- Dic/26 & 0\% & -- & No iniciada \\
			\bottomrule
		\end{tabular}
		\caption{Avance acumulado por actividad al \fechaCorte. ``Esperado al corte'' es la proporción del periodo
		programado en el cronograma que ha transcurrido; ``Cumplimiento'' compara el avance real con dicho valor.}
	\end{table}

	\section{Evidencias por actividad}

	\subsection*{Fase 1: Planificación y Revisión Bibliográfica}

	\actividad{Estado del arte}
	{Fase 1}{Ago/25 -- Oct/25}
	{Documento escrito.}
	{Capítulo 5, sección 5.1 \textit{Estado del arte, aplicaciones y líneas actuales de investigación} (pp.~61--77);
	Capítulo 1 \textit{Introducción} (pp.~1--6).}
	{Revisión de trabajos sobre autómatas celulares neuronales: morfogénesis, NCA isotrópicos, autoclasificación,
	representaciones tridimensionales, NCA en malla, modelos generativos y aplicaciones adicionales.
	Actividad concluida en Trabajo Terminal I.}
	{100\%}{}

	\actividad{Marco teórico y metodológico}
	{Fase 1}{Sep/25 -- Oct/25}
	{Documento escrito.}
	{Capítulos 2 a 5 (marco teórico, pp.~7--77) y Capítulo 6 \textit{Marco metodológico} (pp.~78--90).}
	{Marco teórico de sistemas complejos y teoría del caos, autómatas celulares, inteligencia artificial y redes
	neuronales, y autómatas celulares neuronales. Marco metodológico con diseño cuasiexperimental, población,
	muestra, técnicas de recolección y procesamiento de datos, y contextualización tecnológica.
	Actividad concluida en Trabajo Terminal I.}
	{100\%}{}

	\subsection*{Fase 2: Diseño e implementación del simulador}

	\actividad{Arquitectura del simulador}
	{Fase 2}{Oct/25 -- Nov/25}
	{Documento escrito (análisis y diseño de software, diagramas UML).}
	{Capítulo 7 \textit{Análisis de software} (pp.~91--135); Capítulo 8, secciones 8.1 \textit{Arquitectura},
	8.2 \textit{Diagrama de clases}, 8.3 \textit{Diseño de datos} y 8.5 \textit{Interfaces de usuarios} (pp.~136--143, 153--156).}
	{Metodología de desarrollo, reglas de negocio, requerimientos, análisis de riesgos, estudio de factibilidad,
	casos de uso con sus especificaciones, diagramas de secuencia, arquitectura MVC, diagrama de clases,
	diseño de datos e interfaces gráficas.}
	{100\%}{}

	\actividad{Módulos de simulación}
	{Fase 2}{Nov/25 -- Dic/25}
	{Documento escrito, tablas de resultados y pruebas de operatividad.}
	{Sección 8.4 \textit{Diseño de algoritmos} (8.4.1--8.4.4 y 8.4.6, pp.~144--152); secciones 9.1 \textit{Entorno de desarrollo}
	y 9.2 \textit{Estructura del proyecto} (pp.~157--162); subsecciones 9.4.1, 9.4.3, 9.4.4 y 9.4.5 de
	\textit{Validación y pruebas} (pp.~166--169), tablas 9.6 y 9.7.}
	{Núcleo de simulación en C++/CUDA con enlaces a Python e interfaz en Qt: indexación de la regla semitotalística,
	paso de evolución, condiciones de frontera, \textit{forward pass} del NCA y \textit{prefetch} asíncrono.
	Validación funcional contra reglas conocidas y pruebas de rendimiento (hasta 345\,000 pasos/s en
	$100\times100$ y más de 100 fps de renderizado en $2000\times2000$).}
	{100\%}{\videoSimulacion}

	\actividad{Módulos entrenamiento}
	{Fase 2}{Dic/25 -- Ago/26}
	{Documento escrito, pruebas de operatividad.}
	{Subsección 7.8.4 \textit{Especificaciones de casos de uso: Entrenamiento} (pp.~125--130); 8.3.5
	\textit{Configuración del NCA} (p.~143); 8.4.5 \textit{Entrenamiento del autómata celular neuronal} (pp.~150--151);
	sección 9.3 \textit{Entrenamiento del modelo} (pp.~163--166, tabla 9.5); subsecciones 9.4.2 y 9.4.6 (pp.~167, 169).}
	{Módulo de entrenamiento en PyTorch con construcción dinámica del perceptrón multicapa, ciclo de optimización
	con Adam y detención temprana, persistencia de pesos en \texttt{.npz} compatible con el \textit{kernel} CUDA e
	integración con la pestaña de entrenamiento de la GUI (subproceso con \textit{log} en tiempo real).
	Se agregó el parámetro de semilla aleatoria para reproducir experimentos.
	\segun{Actividad concluida en el mes reportado.}{Actividad concluida en agosto de 2026.}}
	{100\%}{\videoEntrenamiento}

	\subsection*{Fase 3: Entrenamiento y Validación}

	\actividad{Entrenamiento y ajuste}
	{Fase 3}{Ago/26 -- Oct/26}
	{Documento escrito y tablas de resultados.}
	{Capítulo 10, sección 10.1 \textit{Duplicación de reglas del juego de la vida}:
	\segun{tablas 10.1 y 10.2 y figura 10.3 (pp.~173--174).}
	{tablas 10.1 y 10.2 y figura 10.3 (pp.~173--174); figuras 10.5 a 10.9 (pp.~176--178).}}
	{\segun{Barrido de hiperparámetros para la regla B3/S23: 10 combinaciones de funciones de activación con
	una capa oculta de 4 neuronas (una sola converge) y con 4 capas ocultas de 8 neuronas (cinco convergen a una
	pérdida menor a $10^{-4}$). Se identifica el estancamiento en las primeras 12 épocas.}
	{Además del barrido de hiperparámetros de agosto, se entrenaron las reglas B3/S23, B3/S238, B38/S23, B3/S236
	y B38/S238 con distintas pérdidas objetivo y semillas, midiendo la distancia Manhattan entre el vector de regla
	aprendido y el original. Se determinó que la pérdida objetivo debe ser $\leq 0.1$ para B3/S23.}}
	{\realEntrenamiento\% (esperado: \espEntrenamiento\%)}{}

	\actividad{Experimentación y análisis}
	{Fase 3}{Ago/26 -- Nov/26}
	{Documento escrito, tablas y gráficas de resultados.}
	{Capítulo 10, sección 10.1:
	\segun{figuras 10.1 y 10.2 (pp.~170--172).}
	{figuras 10.1 a 10.13 (pp.~170--181).}}
	{\segun{Generación del conjunto de datos (100 condiciones iniciales en $100\times100$ durante 100 pasos) y
	análisis de su balance: el estado 1 aparece un $16.5\% \pm 4.2\%$ y los patrones $C0N8$ y $C1N8$ representan
	menos del $0.02\%$ de las entradas.}
	{Primer experimento (duplicación de reglas del juego de la vida) completo. Se diseñó un espacio con las
	512 permutaciones de la vecindad para extraer el vector de regla aprendido y verificar si el NCA conserva la
	propiedad semitotalística. Se identificaron reglas aprendidas alternativas (B3/S23 en lugar de B3/S238,
	B3/S23678 en lugar de B3/S236) y reglas no semitotalísticas.}}
	{\realExperimentos\% (esperado: \espExperimentos\%)}{}

	\actividad{Análisis de resultados (global)}
	{Fase 3}{Ago/26 -- Nov/26}
	{Documento escrito.}
	{Capítulo 10, sección 10.1\segun{ (pp.~170--172).}{, conclusiones del experimento (p.~181).}}
	{\segun{Análisis estadístico de la distribución de estados y de patrones de vecindad para la regla B3/S23,
	identificando el sesgo de muestra y proponiendo sobremuestreo o ponderación de la función de costo.}
	{Conclusiones del primer experimento: la generalización depende del sesgo de muestra de patrones poco frecuentes,
	el entrenamiento de reglas semitotalísticas no garantiza una regla aprendida semitotalística, y no se requieren
	arquitecturas complejas para replicar estas dinámicas. Se valida la metodología de extracción del vector de
	regla para los siguientes experimentos.}}
	{\realAnalisis\% (esperado: \espAnalisis\%)}{}

	\subsection*{Fase 4: Documentación y Difusión}

	\actividad{Documentación técnica}
	{Fase 4}{Oct/26 -- Dic/26}
	{No aplica en el periodo.}
	{Reporte técnico de TT completo (204 páginas).}
	{Actividad programada a partir de octubre de 2026. El reporte técnico se ha redactado de forma incremental
	conforme avanza cada actividad, por lo que su integración final partirá del documento adjunto.}
	{No iniciada (0\% esperado al corte)}{}

	\actividad{Artículo académico}
	{Fase 4}{Oct/26 -- Dic/26}
	{No aplica en el periodo.}
	{--}
	{Actividad programada a partir de octubre de 2026.}
	{No iniciada (0\% esperado al corte)}{}

	\section{Cronograma}

	\resizebox{\linewidth}{!}{%
		\begin{tabular}{p{5cm}*{10}{>{\centering\arraybackslash}p{1cm}}}
			\toprule
			\textbf{Actividad} & \textbf{Ago/25} & \textbf{Sep/25} & \textbf{Oct/25} & \textbf{Nov/25} & \textbf{Dic/25} & \mesActual{1}{Ago/26} & \mesActual{2}{Sep/26} & \textbf{Oct/26} & \textbf{Nov/26} & \textbf{Dic/26} \\
			\midrule
			\multicolumn{11}{l}{\textbf{Fase 1: Planificación y Revisión Bibliográfica}} \\

			Estado del arte       & \cellcolor{gray!30} & \cellcolor{gray!30} & \cellcolor{gray!30} &  &  &  &  &  &  &   \\

			Marco teórico y metodológico &  & \cellcolor{gray!30} & \cellcolor{gray!30} &  &  &  &  &  &  &   \\

			\midrule

			\multicolumn{11}{l}{\textbf{Fase 2: Diseño e implementación del simulador}} \\
			Arquitectura del simulador &  &  & \cellcolor{gray!30} & \cellcolor{gray!30} &  &  &  &  &  &   \\

			Módulos de simulación  &  &  &  & \cellcolor{gray!30} & \cellcolor{gray!30} &  &  &  &  &   \\

			Módulos entrenamiento   &  &  &  &  & \cellcolor{gray!30} & \cellcolor{gray!30} &  &  &  &   \\

			\midrule

			\multicolumn{11}{l}{\textbf{Fase 3: Entrenamiento y Validación}} \\
			Entrenamiento y ajuste   &  &  &  &  &  & \cellcolor{gray!30} & \cellcolor{gray!30} & \cellcolor{gray!30} &   &   \\

			Experimentación y análisis  &  &  &  &  &  & \cellcolor{gray!30}  & \cellcolor{gray!30} & \cellcolor{gray!30} & \cellcolor{gray!30} &   \\

			Análisis de resultados (global) &  &  &  &  &  & \cellcolor{gray!30} & \cellcolor{gray!30} & \cellcolor{gray!30} & \cellcolor{gray!30} &   \\

			\midrule

			\multicolumn{11}{l}{\textbf{Fase 4: Documentación y Difusión}} \\
			Documentación técnica  &  &  &  &  &  &  &  & \cellcolor{gray!30}  & \cellcolor{gray!30} & \cellcolor{gray!30}  \\

			Artículo académico  &  &  &  &  &  &  &  & \cellcolor{gray!30}  & \cellcolor{gray!30}  & \cellcolor{gray!30}  \\

			\bottomrule
		\end{tabular}
	}

	\vspace{0.5em}
	{\small La columna resaltada corresponde al mes reportado.}

	\section{Alumnos y directores}

	\vspace{1em}
	\begin{tabular}{@{} p{0.45\textwidth} p{0.45\textwidth} @{}}
		\textit{Erick Jesse Angeles López} & \textit{Dr. Genaro Juárez Martínez} \\[5em]
		Firma: \rule{6cm}{0.4pt} & Firma: \rule{6cm}{0.4pt} \\
	\end{tabular}

\end{document}
"
%      pdflatex -jobname=2ReporteSeguimiento "\def\corte{2}% ============================================================================
%  Reporte de seguimiento - Trabajo Terminal II
%
%  Un solo archivo genera los reportes de cada mes. El mes se elige con \corte:
%      1 = Agosto 2026, 2 = Septiembre 2026
%
%  Compilación (desde esta carpeta):
%      pdflatex -jobname=1ReporteSeguimiento "\def\corte{1}\input{reporteSeguimiento}"
%      pdflatex -jobname=2ReporteSeguimiento "\def\corte{2}\input{reporteSeguimiento}"
%  o simplemente:  ./compilar.sh
%
%  Las páginas y secciones citadas corresponden a Documentation/main.pdf
%  (reporte técnico de TT), que es el PDF que se adjunta como evidencia.
% ============================================================================
\documentclass[letterpaper,10pt,oneside]{article}
\usepackage[margin=2.54cm]{geometry}
\usepackage[spanish]{babel}
\newcommand\authormark[1]{\textsuperscript{#1}}
\usepackage[colorlinks=true,bookmarks=true,linkcolor=black,citecolor=black,urlcolor=blue]{hyperref}
\usepackage{times}
\usepackage[table]{xcolor}
\usepackage{booktabs}
\usepackage{graphicx}
\usepackage{titlesec}
\usepackage{tabularx}
\usepackage{array}
\usepackage{enumitem}
\usepackage{float}

\titleformat{\section}
{\normalfont\fontsize{12pt}{14pt}\bfseries}
{\thesection}{1em}{}

% ----------------------------------------------------------------------------
%  Mes reportado
% ----------------------------------------------------------------------------
\providecommand\corte{2}
\ifnum\corte=1
	\def\numReporte{1}
	\def\mes{Agosto}
	\def\fechaCorte{31 de agosto de 2026}
\else
	\def\numReporte{2}
	\def\mes{Septiembre}
	\def\fechaCorte{30 de septiembre de 2026}
\fi
% \segun{texto agosto}{texto septiembre}
\newcommand\segun[2]{\ifnum\corte=1 #1\else #2\fi}

% ----------------------------------------------------------------------------
%  Porcentajes (editar aquí). Formato: \segun{agosto}{septiembre}
%  "esp" = avance esperado según cronograma al corte, "real" = avance logrado.
% ----------------------------------------------------------------------------
\def\espEntrenamiento{\segun{33}{67}}   \def\realEntrenamiento{\segun{33}{67}}
\def\espExperimentos{\segun{25}{50}}    \def\realExperimentos{\segun{25}{50}}
\def\espAnalisis{\segun{25}{50}}        \def\realAnalisis{\segun{25}{50}}

% ----------------------------------------------------------------------------
%  Enlaces a video (dejar vacío si no aplica)
% ----------------------------------------------------------------------------
\def\videoSimulacion{}
\def\videoEntrenamiento{}

% ----------------------------------------------------------------------------
%  Bloque de evidencia por actividad
%  \actividad{nombre}{fase}{periodo}{tipo de evidencia}{ubicación en reporte TT}
%            {descripción}{porcentaje}{enlace de video}
% ----------------------------------------------------------------------------
\definecolor{encabezado}{gray}{0.85}
\newcommand\video[1]{\edef\tmp{#1}\ifx\tmp\empty No aplica.\else\url{#1}\fi}
\newcommand\actividad[8]{%
	\par\noindent
	\begin{minipage}{\linewidth}
		\renewcommand{\arraystretch}{1.25}
		\begin{tabularx}{\linewidth}{|>{\raggedright\bfseries\arraybackslash}p{0.24\linewidth}|X|}
			\hline
			\rowcolor{encabezado}
			Nombre de la actividad & \textbf{#1} \\ \hline
			Fase y periodo & #2 \,(#3) \\ \hline
			Evidencia de trabajo & #4 \\ \hline
			Ubicación en el reporte de TT & #5 \\ \hline
			Descripción del avance & #6 \\ \hline
			Porcentaje de cumplimiento & \textbf{#7} \\ \hline
			Enlace de video & \video{#8} \\ \hline
		\end{tabularx}
	\end{minipage}
	\par\vspace{1.2em}%
}

% Porcentaje de cumplimiento: real respecto a lo esperado al corte
\newcommand\cumplimiento[2]{\fpeval{min(100, round(100*(#2)/(#1), 0))}\%}

% Columna del mes reportado en el cronograma
\newcommand\mesActual[2]{\ifnum\corte=#1\cellcolor{black!70}\textcolor{white}{\textbf{#2}}\else\textbf{#2}\fi}

\begin{document}

	\begin{center}
		\fontsize{14}{16}\selectfont\bfseries
		\numReporte° Reporte de seguimiento (\mes{} 2026): \\
		Simulación de sistemas complejos mediante autómatas celulares neuronales
	\end{center}

	\begin{center}
		\fontsize{14}{16}\selectfont\bfseries
		\textit{Trabajo Terminal No. 2026-A114}
	\end{center}

	\begin{center}
		\fontsize{12}{16}\selectfont
		\textit{Alumno: \authormark{*}Angeles López Erick Jesse\authormark{1}} \\
		\textit{Director: Juárez Martínez Genaro\authormark{1}} \\
		\textit{\authormark{*}e-mail: eangelesl1700@alumno.ipn.mx}
	\end{center}

	\section{Información general}

	\begin{itemize}[leftmargin=*]
		\item \textbf{Periodo reportado:} \mes{} 2026 (corte al \fechaCorte).
		\item \textbf{Documento de evidencia adjunto:} reporte técnico de Trabajo Terminal
		\textit{``Simulación de sistemas complejos mediante autómatas celulares neuronales''} (PDF, 204 páginas).
		El avance es incremental: cada evidencia listada en este reporte se encuentra integrada en el capítulo
		y sección correspondiente de dicho documento, cuyas páginas se indican en cada actividad.
		\item \textbf{Actividades programadas en el mes:} \segun{Módulos de entrenamiento (cierre), Entrenamiento y ajuste,
		Experimentación y análisis, Análisis de resultados (global).}{Entrenamiento y ajuste, Experimentación y análisis,
		Análisis de resultados (global).}
	\end{itemize}

	\section{Resumen de cumplimiento}

	\begin{table}[H]
		\centering
		\renewcommand{\arraystretch}{1.2}
		\begin{tabular}{l c c c c}
			\toprule
			\textbf{Actividad} & \textbf{Periodo} & \textbf{Esperado al corte} & \textbf{Avance real} & \textbf{Cumplimiento} \\
			\midrule
			\multicolumn{5}{l}{\textit{Fase 1: Planificación y Revisión Bibliográfica}} \\
			Estado del arte                  & Ago/25 -- Oct/25 & 100\% & 100\% & 100\% \\
			Marco teórico y metodológico     & Sep/25 -- Oct/25 & 100\% & 100\% & 100\% \\
			\multicolumn{5}{l}{\textit{Fase 2: Diseño e implementación del simulador}} \\
			Arquitectura del simulador       & Oct/25 -- Nov/25 & 100\% & 100\% & 100\% \\
			Módulos de simulación            & Nov/25 -- Dic/25 & 100\% & 100\% & 100\% \\
			Módulos entrenamiento            & Dic/25 -- Ago/26 & 100\% & 100\% & 100\% \\
			\multicolumn{5}{l}{\textit{Fase 3: Entrenamiento y Validación}} \\
			Entrenamiento y ajuste           & Ago/26 -- Oct/26 & \espEntrenamiento\% & \realEntrenamiento\% & \cumplimiento{\espEntrenamiento}{\realEntrenamiento} \\
			Experimentación y análisis       & Ago/26 -- Nov/26 & \espExperimentos\%  & \realExperimentos\%  & \cumplimiento{\espExperimentos}{\realExperimentos} \\
			Análisis de resultados (global)  & Ago/26 -- Nov/26 & \espAnalisis\%      & \realAnalisis\%      & \cumplimiento{\espAnalisis}{\realAnalisis} \\
			\multicolumn{5}{l}{\textit{Fase 4: Documentación y Difusión}} \\
			Documentación técnica            & Oct/26 -- Dic/26 & 0\% & -- & No iniciada \\
			Artículo académico               & Oct/26 -- Dic/26 & 0\% & -- & No iniciada \\
			\bottomrule
		\end{tabular}
		\caption{Avance acumulado por actividad al \fechaCorte. ``Esperado al corte'' es la proporción del periodo
		programado en el cronograma que ha transcurrido; ``Cumplimiento'' compara el avance real con dicho valor.}
	\end{table}

	\section{Evidencias por actividad}

	\subsection*{Fase 1: Planificación y Revisión Bibliográfica}

	\actividad{Estado del arte}
	{Fase 1}{Ago/25 -- Oct/25}
	{Documento escrito.}
	{Capítulo 5, sección 5.1 \textit{Estado del arte, aplicaciones y líneas actuales de investigación} (pp.~61--77);
	Capítulo 1 \textit{Introducción} (pp.~1--6).}
	{Revisión de trabajos sobre autómatas celulares neuronales: morfogénesis, NCA isotrópicos, autoclasificación,
	representaciones tridimensionales, NCA en malla, modelos generativos y aplicaciones adicionales.
	Actividad concluida en Trabajo Terminal I.}
	{100\%}{}

	\actividad{Marco teórico y metodológico}
	{Fase 1}{Sep/25 -- Oct/25}
	{Documento escrito.}
	{Capítulos 2 a 5 (marco teórico, pp.~7--77) y Capítulo 6 \textit{Marco metodológico} (pp.~78--90).}
	{Marco teórico de sistemas complejos y teoría del caos, autómatas celulares, inteligencia artificial y redes
	neuronales, y autómatas celulares neuronales. Marco metodológico con diseño cuasiexperimental, población,
	muestra, técnicas de recolección y procesamiento de datos, y contextualización tecnológica.
	Actividad concluida en Trabajo Terminal I.}
	{100\%}{}

	\subsection*{Fase 2: Diseño e implementación del simulador}

	\actividad{Arquitectura del simulador}
	{Fase 2}{Oct/25 -- Nov/25}
	{Documento escrito (análisis y diseño de software, diagramas UML).}
	{Capítulo 7 \textit{Análisis de software} (pp.~91--135); Capítulo 8, secciones 8.1 \textit{Arquitectura},
	8.2 \textit{Diagrama de clases}, 8.3 \textit{Diseño de datos} y 8.5 \textit{Interfaces de usuarios} (pp.~136--143, 153--156).}
	{Metodología de desarrollo, reglas de negocio, requerimientos, análisis de riesgos, estudio de factibilidad,
	casos de uso con sus especificaciones, diagramas de secuencia, arquitectura MVC, diagrama de clases,
	diseño de datos e interfaces gráficas.}
	{100\%}{}

	\actividad{Módulos de simulación}
	{Fase 2}{Nov/25 -- Dic/25}
	{Documento escrito, tablas de resultados y pruebas de operatividad.}
	{Sección 8.4 \textit{Diseño de algoritmos} (8.4.1--8.4.4 y 8.4.6, pp.~144--152); secciones 9.1 \textit{Entorno de desarrollo}
	y 9.2 \textit{Estructura del proyecto} (pp.~157--162); subsecciones 9.4.1, 9.4.3, 9.4.4 y 9.4.5 de
	\textit{Validación y pruebas} (pp.~166--169), tablas 9.6 y 9.7.}
	{Núcleo de simulación en C++/CUDA con enlaces a Python e interfaz en Qt: indexación de la regla semitotalística,
	paso de evolución, condiciones de frontera, \textit{forward pass} del NCA y \textit{prefetch} asíncrono.
	Validación funcional contra reglas conocidas y pruebas de rendimiento (hasta 345\,000 pasos/s en
	$100\times100$ y más de 100 fps de renderizado en $2000\times2000$).}
	{100\%}{\videoSimulacion}

	\actividad{Módulos entrenamiento}
	{Fase 2}{Dic/25 -- Ago/26}
	{Documento escrito, pruebas de operatividad.}
	{Subsección 7.8.4 \textit{Especificaciones de casos de uso: Entrenamiento} (pp.~125--130); 8.3.5
	\textit{Configuración del NCA} (p.~143); 8.4.5 \textit{Entrenamiento del autómata celular neuronal} (pp.~150--151);
	sección 9.3 \textit{Entrenamiento del modelo} (pp.~163--166, tabla 9.5); subsecciones 9.4.2 y 9.4.6 (pp.~167, 169).}
	{Módulo de entrenamiento en PyTorch con construcción dinámica del perceptrón multicapa, ciclo de optimización
	con Adam y detención temprana, persistencia de pesos en \texttt{.npz} compatible con el \textit{kernel} CUDA e
	integración con la pestaña de entrenamiento de la GUI (subproceso con \textit{log} en tiempo real).
	Se agregó el parámetro de semilla aleatoria para reproducir experimentos.
	\segun{Actividad concluida en el mes reportado.}{Actividad concluida en agosto de 2026.}}
	{100\%}{\videoEntrenamiento}

	\subsection*{Fase 3: Entrenamiento y Validación}

	\actividad{Entrenamiento y ajuste}
	{Fase 3}{Ago/26 -- Oct/26}
	{Documento escrito y tablas de resultados.}
	{Capítulo 10, sección 10.1 \textit{Duplicación de reglas del juego de la vida}:
	\segun{tablas 10.1 y 10.2 y figura 10.3 (pp.~173--174).}
	{tablas 10.1 y 10.2 y figura 10.3 (pp.~173--174); figuras 10.5 a 10.9 (pp.~176--178).}}
	{\segun{Barrido de hiperparámetros para la regla B3/S23: 10 combinaciones de funciones de activación con
	una capa oculta de 4 neuronas (una sola converge) y con 4 capas ocultas de 8 neuronas (cinco convergen a una
	pérdida menor a $10^{-4}$). Se identifica el estancamiento en las primeras 12 épocas.}
	{Además del barrido de hiperparámetros de agosto, se entrenaron las reglas B3/S23, B3/S238, B38/S23, B3/S236
	y B38/S238 con distintas pérdidas objetivo y semillas, midiendo la distancia Manhattan entre el vector de regla
	aprendido y el original. Se determinó que la pérdida objetivo debe ser $\leq 0.1$ para B3/S23.}}
	{\realEntrenamiento\% (esperado: \espEntrenamiento\%)}{}

	\actividad{Experimentación y análisis}
	{Fase 3}{Ago/26 -- Nov/26}
	{Documento escrito, tablas y gráficas de resultados.}
	{Capítulo 10, sección 10.1:
	\segun{figuras 10.1 y 10.2 (pp.~170--172).}
	{figuras 10.1 a 10.13 (pp.~170--181).}}
	{\segun{Generación del conjunto de datos (100 condiciones iniciales en $100\times100$ durante 100 pasos) y
	análisis de su balance: el estado 1 aparece un $16.5\% \pm 4.2\%$ y los patrones $C0N8$ y $C1N8$ representan
	menos del $0.02\%$ de las entradas.}
	{Primer experimento (duplicación de reglas del juego de la vida) completo. Se diseñó un espacio con las
	512 permutaciones de la vecindad para extraer el vector de regla aprendido y verificar si el NCA conserva la
	propiedad semitotalística. Se identificaron reglas aprendidas alternativas (B3/S23 en lugar de B3/S238,
	B3/S23678 en lugar de B3/S236) y reglas no semitotalísticas.}}
	{\realExperimentos\% (esperado: \espExperimentos\%)}{}

	\actividad{Análisis de resultados (global)}
	{Fase 3}{Ago/26 -- Nov/26}
	{Documento escrito.}
	{Capítulo 10, sección 10.1\segun{ (pp.~170--172).}{, conclusiones del experimento (p.~181).}}
	{\segun{Análisis estadístico de la distribución de estados y de patrones de vecindad para la regla B3/S23,
	identificando el sesgo de muestra y proponiendo sobremuestreo o ponderación de la función de costo.}
	{Conclusiones del primer experimento: la generalización depende del sesgo de muestra de patrones poco frecuentes,
	el entrenamiento de reglas semitotalísticas no garantiza una regla aprendida semitotalística, y no se requieren
	arquitecturas complejas para replicar estas dinámicas. Se valida la metodología de extracción del vector de
	regla para los siguientes experimentos.}}
	{\realAnalisis\% (esperado: \espAnalisis\%)}{}

	\subsection*{Fase 4: Documentación y Difusión}

	\actividad{Documentación técnica}
	{Fase 4}{Oct/26 -- Dic/26}
	{No aplica en el periodo.}
	{Reporte técnico de TT completo (204 páginas).}
	{Actividad programada a partir de octubre de 2026. El reporte técnico se ha redactado de forma incremental
	conforme avanza cada actividad, por lo que su integración final partirá del documento adjunto.}
	{No iniciada (0\% esperado al corte)}{}

	\actividad{Artículo académico}
	{Fase 4}{Oct/26 -- Dic/26}
	{No aplica en el periodo.}
	{--}
	{Actividad programada a partir de octubre de 2026.}
	{No iniciada (0\% esperado al corte)}{}

	\section{Cronograma}

	\resizebox{\linewidth}{!}{%
		\begin{tabular}{p{5cm}*{10}{>{\centering\arraybackslash}p{1cm}}}
			\toprule
			\textbf{Actividad} & \textbf{Ago/25} & \textbf{Sep/25} & \textbf{Oct/25} & \textbf{Nov/25} & \textbf{Dic/25} & \mesActual{1}{Ago/26} & \mesActual{2}{Sep/26} & \textbf{Oct/26} & \textbf{Nov/26} & \textbf{Dic/26} \\
			\midrule
			\multicolumn{11}{l}{\textbf{Fase 1: Planificación y Revisión Bibliográfica}} \\

			Estado del arte       & \cellcolor{gray!30} & \cellcolor{gray!30} & \cellcolor{gray!30} &  &  &  &  &  &  &   \\

			Marco teórico y metodológico &  & \cellcolor{gray!30} & \cellcolor{gray!30} &  &  &  &  &  &  &   \\

			\midrule

			\multicolumn{11}{l}{\textbf{Fase 2: Diseño e implementación del simulador}} \\
			Arquitectura del simulador &  &  & \cellcolor{gray!30} & \cellcolor{gray!30} &  &  &  &  &  &   \\

			Módulos de simulación  &  &  &  & \cellcolor{gray!30} & \cellcolor{gray!30} &  &  &  &  &   \\

			Módulos entrenamiento   &  &  &  &  & \cellcolor{gray!30} & \cellcolor{gray!30} &  &  &  &   \\

			\midrule

			\multicolumn{11}{l}{\textbf{Fase 3: Entrenamiento y Validación}} \\
			Entrenamiento y ajuste   &  &  &  &  &  & \cellcolor{gray!30} & \cellcolor{gray!30} & \cellcolor{gray!30} &   &   \\

			Experimentación y análisis  &  &  &  &  &  & \cellcolor{gray!30}  & \cellcolor{gray!30} & \cellcolor{gray!30} & \cellcolor{gray!30} &   \\

			Análisis de resultados (global) &  &  &  &  &  & \cellcolor{gray!30} & \cellcolor{gray!30} & \cellcolor{gray!30} & \cellcolor{gray!30} &   \\

			\midrule

			\multicolumn{11}{l}{\textbf{Fase 4: Documentación y Difusión}} \\
			Documentación técnica  &  &  &  &  &  &  &  & \cellcolor{gray!30}  & \cellcolor{gray!30} & \cellcolor{gray!30}  \\

			Artículo académico  &  &  &  &  &  &  &  & \cellcolor{gray!30}  & \cellcolor{gray!30}  & \cellcolor{gray!30}  \\

			\bottomrule
		\end{tabular}
	}

	\vspace{0.5em}
	{\small La columna resaltada corresponde al mes reportado.}

	\section{Alumnos y directores}

	\vspace{1em}
	\begin{tabular}{@{} p{0.45\textwidth} p{0.45\textwidth} @{}}
		\textit{Erick Jesse Angeles López} & \textit{Dr. Genaro Juárez Martínez} \\[5em]
		Firma: \rule{6cm}{0.4pt} & Firma: \rule{6cm}{0.4pt} \\
	\end{tabular}

\end{document}
"
%  o simplemente:  ./compilar.sh
%
%  Las páginas y secciones citadas corresponden a Documentation/main.pdf
%  (reporte técnico de TT), que es el PDF que se adjunta como evidencia.
% ============================================================================
\documentclass[letterpaper,10pt,oneside]{article}
\usepackage[margin=2.54cm]{geometry}
\usepackage[spanish]{babel}
\newcommand\authormark[1]{\textsuperscript{#1}}
\usepackage[colorlinks=true,bookmarks=true,linkcolor=black,citecolor=black,urlcolor=blue]{hyperref}
\usepackage{times}
\usepackage[table]{xcolor}
\usepackage{booktabs}
\usepackage{graphicx}
\usepackage{titlesec}
\usepackage{tabularx}
\usepackage{array}
\usepackage{enumitem}
\usepackage{float}

\titleformat{\section}
{\normalfont\fontsize{12pt}{14pt}\bfseries}
{\thesection}{1em}{}

% ----------------------------------------------------------------------------
%  Mes reportado
% ----------------------------------------------------------------------------
\providecommand\corte{2}
\ifnum\corte=1
	\def\numReporte{1}
	\def\mes{Agosto}
	\def\fechaCorte{31 de agosto de 2026}
\else
	\def\numReporte{2}
	\def\mes{Septiembre}
	\def\fechaCorte{30 de septiembre de 2026}
\fi
% \segun{texto agosto}{texto septiembre}
\newcommand\segun[2]{\ifnum\corte=1 #1\else #2\fi}

% ----------------------------------------------------------------------------
%  Porcentajes (editar aquí). Formato: \segun{agosto}{septiembre}
%  "esp" = avance esperado según cronograma al corte, "real" = avance logrado.
% ----------------------------------------------------------------------------
\def\espEntrenamiento{\segun{33}{67}}   \def\realEntrenamiento{\segun{33}{67}}
\def\espExperimentos{\segun{25}{50}}    \def\realExperimentos{\segun{25}{50}}
\def\espAnalisis{\segun{25}{50}}        \def\realAnalisis{\segun{25}{50}}

% ----------------------------------------------------------------------------
%  Enlaces a video (dejar vacío si no aplica)
% ----------------------------------------------------------------------------
\def\videoSimulacion{}
\def\videoEntrenamiento{}

% ----------------------------------------------------------------------------
%  Bloque de evidencia por actividad
%  \actividad{nombre}{fase}{periodo}{tipo de evidencia}{ubicación en reporte TT}
%            {descripción}{porcentaje}{enlace de video}
% ----------------------------------------------------------------------------
\definecolor{encabezado}{gray}{0.85}
\newcommand\video[1]{\edef\tmp{#1}\ifx\tmp\empty No aplica.\else\url{#1}\fi}
\newcommand\actividad[8]{%
	\par\noindent
	\begin{minipage}{\linewidth}
		\renewcommand{\arraystretch}{1.25}
		\begin{tabularx}{\linewidth}{|>{\raggedright\bfseries\arraybackslash}p{0.24\linewidth}|X|}
			\hline
			\rowcolor{encabezado}
			Nombre de la actividad & \textbf{#1} \\ \hline
			Fase y periodo & #2 \,(#3) \\ \hline
			Evidencia de trabajo & #4 \\ \hline
			Ubicación en el reporte de TT & #5 \\ \hline
			Descripción del avance & #6 \\ \hline
			Porcentaje de cumplimiento & \textbf{#7} \\ \hline
			Enlace de video & \video{#8} \\ \hline
		\end{tabularx}
	\end{minipage}
	\par\vspace{1.2em}%
}

% Porcentaje de cumplimiento: real respecto a lo esperado al corte
\newcommand\cumplimiento[2]{\fpeval{min(100, round(100*(#2)/(#1), 0))}\%}

% Columna del mes reportado en el cronograma
\newcommand\mesActual[2]{\ifnum\corte=#1\cellcolor{black!70}\textcolor{white}{\textbf{#2}}\else\textbf{#2}\fi}

\begin{document}

	\begin{center}
		\fontsize{14}{16}\selectfont\bfseries
		\numReporte° Reporte de seguimiento (\mes{} 2026): \\
		Simulación de sistemas complejos mediante autómatas celulares neuronales
	\end{center}

	\begin{center}
		\fontsize{14}{16}\selectfont\bfseries
		\textit{Trabajo Terminal No. 2026-A114}
	\end{center}

	\begin{center}
		\fontsize{12}{16}\selectfont
		\textit{Alumno: \authormark{*}Angeles López Erick Jesse\authormark{1}} \\
		\textit{Director: Juárez Martínez Genaro\authormark{1}} \\
		\textit{\authormark{*}e-mail: eangelesl1700@alumno.ipn.mx}
	\end{center}

	\section{Información general}

	\begin{itemize}[leftmargin=*]
		\item \textbf{Periodo reportado:} \mes{} 2026 (corte al \fechaCorte).
		\item \textbf{Documento de evidencia adjunto:} reporte técnico de Trabajo Terminal
		\textit{``Simulación de sistemas complejos mediante autómatas celulares neuronales''} (PDF, 204 páginas).
		El avance es incremental: cada evidencia listada en este reporte se encuentra integrada en el capítulo
		y sección correspondiente de dicho documento, cuyas páginas se indican en cada actividad.
		\item \textbf{Actividades programadas en el mes:} \segun{Módulos de entrenamiento (cierre), Entrenamiento y ajuste,
		Experimentación y análisis, Análisis de resultados (global).}{Entrenamiento y ajuste, Experimentación y análisis,
		Análisis de resultados (global).}
	\end{itemize}

	\section{Resumen de cumplimiento}

	\begin{table}[H]
		\centering
		\renewcommand{\arraystretch}{1.2}
		\begin{tabular}{l c c c c}
			\toprule
			\textbf{Actividad} & \textbf{Periodo} & \textbf{Esperado al corte} & \textbf{Avance real} & \textbf{Cumplimiento} \\
			\midrule
			\multicolumn{5}{l}{\textit{Fase 1: Planificación y Revisión Bibliográfica}} \\
			Estado del arte                  & Ago/25 -- Oct/25 & 100\% & 100\% & 100\% \\
			Marco teórico y metodológico     & Sep/25 -- Oct/25 & 100\% & 100\% & 100\% \\
			\multicolumn{5}{l}{\textit{Fase 2: Diseño e implementación del simulador}} \\
			Arquitectura del simulador       & Oct/25 -- Nov/25 & 100\% & 100\% & 100\% \\
			Módulos de simulación            & Nov/25 -- Dic/25 & 100\% & 100\% & 100\% \\
			Módulos entrenamiento            & Dic/25 -- Ago/26 & 100\% & 100\% & 100\% \\
			\multicolumn{5}{l}{\textit{Fase 3: Entrenamiento y Validación}} \\
			Entrenamiento y ajuste           & Ago/26 -- Oct/26 & \espEntrenamiento\% & \realEntrenamiento\% & \cumplimiento{\espEntrenamiento}{\realEntrenamiento} \\
			Experimentación y análisis       & Ago/26 -- Nov/26 & \espExperimentos\%  & \realExperimentos\%  & \cumplimiento{\espExperimentos}{\realExperimentos} \\
			Análisis de resultados (global)  & Ago/26 -- Nov/26 & \espAnalisis\%      & \realAnalisis\%      & \cumplimiento{\espAnalisis}{\realAnalisis} \\
			\multicolumn{5}{l}{\textit{Fase 4: Documentación y Difusión}} \\
			Documentación técnica            & Oct/26 -- Dic/26 & 0\% & -- & No iniciada \\
			Artículo académico               & Oct/26 -- Dic/26 & 0\% & -- & No iniciada \\
			\bottomrule
		\end{tabular}
		\caption{Avance acumulado por actividad al \fechaCorte. ``Esperado al corte'' es la proporción del periodo
		programado en el cronograma que ha transcurrido; ``Cumplimiento'' compara el avance real con dicho valor.}
	\end{table}

	\section{Evidencias por actividad}

	\subsection*{Fase 1: Planificación y Revisión Bibliográfica}

	\actividad{Estado del arte}
	{Fase 1}{Ago/25 -- Oct/25}
	{Documento escrito.}
	{Capítulo 5, sección 5.1 \textit{Estado del arte, aplicaciones y líneas actuales de investigación} (pp.~61--77);
	Capítulo 1 \textit{Introducción} (pp.~1--6).}
	{Revisión de trabajos sobre autómatas celulares neuronales: morfogénesis, NCA isotrópicos, autoclasificación,
	representaciones tridimensionales, NCA en malla, modelos generativos y aplicaciones adicionales.
	Actividad concluida en Trabajo Terminal I.}
	{100\%}{}

	\actividad{Marco teórico y metodológico}
	{Fase 1}{Sep/25 -- Oct/25}
	{Documento escrito.}
	{Capítulos 2 a 5 (marco teórico, pp.~7--77) y Capítulo 6 \textit{Marco metodológico} (pp.~78--90).}
	{Marco teórico de sistemas complejos y teoría del caos, autómatas celulares, inteligencia artificial y redes
	neuronales, y autómatas celulares neuronales. Marco metodológico con diseño cuasiexperimental, población,
	muestra, técnicas de recolección y procesamiento de datos, y contextualización tecnológica.
	Actividad concluida en Trabajo Terminal I.}
	{100\%}{}

	\subsection*{Fase 2: Diseño e implementación del simulador}

	\actividad{Arquitectura del simulador}
	{Fase 2}{Oct/25 -- Nov/25}
	{Documento escrito (análisis y diseño de software, diagramas UML).}
	{Capítulo 7 \textit{Análisis de software} (pp.~91--135); Capítulo 8, secciones 8.1 \textit{Arquitectura},
	8.2 \textit{Diagrama de clases}, 8.3 \textit{Diseño de datos} y 8.5 \textit{Interfaces de usuarios} (pp.~136--143, 153--156).}
	{Metodología de desarrollo, reglas de negocio, requerimientos, análisis de riesgos, estudio de factibilidad,
	casos de uso con sus especificaciones, diagramas de secuencia, arquitectura MVC, diagrama de clases,
	diseño de datos e interfaces gráficas.}
	{100\%}{}

	\actividad{Módulos de simulación}
	{Fase 2}{Nov/25 -- Dic/25}
	{Documento escrito, tablas de resultados y pruebas de operatividad.}
	{Sección 8.4 \textit{Diseño de algoritmos} (8.4.1--8.4.4 y 8.4.6, pp.~144--152); secciones 9.1 \textit{Entorno de desarrollo}
	y 9.2 \textit{Estructura del proyecto} (pp.~157--162); subsecciones 9.4.1, 9.4.3, 9.4.4 y 9.4.5 de
	\textit{Validación y pruebas} (pp.~166--169), tablas 9.6 y 9.7.}
	{Núcleo de simulación en C++/CUDA con enlaces a Python e interfaz en Qt: indexación de la regla semitotalística,
	paso de evolución, condiciones de frontera, \textit{forward pass} del NCA y \textit{prefetch} asíncrono.
	Validación funcional contra reglas conocidas y pruebas de rendimiento (hasta 345\,000 pasos/s en
	$100\times100$ y más de 100 fps de renderizado en $2000\times2000$).}
	{100\%}{\videoSimulacion}

	\actividad{Módulos entrenamiento}
	{Fase 2}{Dic/25 -- Ago/26}
	{Documento escrito, pruebas de operatividad.}
	{Subsección 7.8.4 \textit{Especificaciones de casos de uso: Entrenamiento} (pp.~125--130); 8.3.5
	\textit{Configuración del NCA} (p.~143); 8.4.5 \textit{Entrenamiento del autómata celular neuronal} (pp.~150--151);
	sección 9.3 \textit{Entrenamiento del modelo} (pp.~163--166, tabla 9.5); subsecciones 9.4.2 y 9.4.6 (pp.~167, 169).}
	{Módulo de entrenamiento en PyTorch con construcción dinámica del perceptrón multicapa, ciclo de optimización
	con Adam y detención temprana, persistencia de pesos en \texttt{.npz} compatible con el \textit{kernel} CUDA e
	integración con la pestaña de entrenamiento de la GUI (subproceso con \textit{log} en tiempo real).
	Se agregó el parámetro de semilla aleatoria para reproducir experimentos.
	\segun{Actividad concluida en el mes reportado.}{Actividad concluida en agosto de 2026.}}
	{100\%}{\videoEntrenamiento}

	\subsection*{Fase 3: Entrenamiento y Validación}

	\actividad{Entrenamiento y ajuste}
	{Fase 3}{Ago/26 -- Oct/26}
	{Documento escrito y tablas de resultados.}
	{Capítulo 10, sección 10.1 \textit{Duplicación de reglas del juego de la vida}:
	\segun{tablas 10.1 y 10.2 y figura 10.3 (pp.~173--174).}
	{tablas 10.1 y 10.2 y figura 10.3 (pp.~173--174); figuras 10.5 a 10.9 (pp.~176--178).}}
	{\segun{Barrido de hiperparámetros para la regla B3/S23: 10 combinaciones de funciones de activación con
	una capa oculta de 4 neuronas (una sola converge) y con 4 capas ocultas de 8 neuronas (cinco convergen a una
	pérdida menor a $10^{-4}$). Se identifica el estancamiento en las primeras 12 épocas.}
	{Además del barrido de hiperparámetros de agosto, se entrenaron las reglas B3/S23, B3/S238, B38/S23, B3/S236
	y B38/S238 con distintas pérdidas objetivo y semillas, midiendo la distancia Manhattan entre el vector de regla
	aprendido y el original. Se determinó que la pérdida objetivo debe ser $\leq 0.1$ para B3/S23.}}
	{\realEntrenamiento\% (esperado: \espEntrenamiento\%)}{}

	\actividad{Experimentación y análisis}
	{Fase 3}{Ago/26 -- Nov/26}
	{Documento escrito, tablas y gráficas de resultados.}
	{Capítulo 10, sección 10.1:
	\segun{figuras 10.1 y 10.2 (pp.~170--172).}
	{figuras 10.1 a 10.13 (pp.~170--181).}}
	{\segun{Generación del conjunto de datos (100 condiciones iniciales en $100\times100$ durante 100 pasos) y
	análisis de su balance: el estado 1 aparece un $16.5\% \pm 4.2\%$ y los patrones $C0N8$ y $C1N8$ representan
	menos del $0.02\%$ de las entradas.}
	{Primer experimento (duplicación de reglas del juego de la vida) completo. Se diseñó un espacio con las
	512 permutaciones de la vecindad para extraer el vector de regla aprendido y verificar si el NCA conserva la
	propiedad semitotalística. Se identificaron reglas aprendidas alternativas (B3/S23 en lugar de B3/S238,
	B3/S23678 en lugar de B3/S236) y reglas no semitotalísticas.}}
	{\realExperimentos\% (esperado: \espExperimentos\%)}{}

	\actividad{Análisis de resultados (global)}
	{Fase 3}{Ago/26 -- Nov/26}
	{Documento escrito.}
	{Capítulo 10, sección 10.1\segun{ (pp.~170--172).}{, conclusiones del experimento (p.~181).}}
	{\segun{Análisis estadístico de la distribución de estados y de patrones de vecindad para la regla B3/S23,
	identificando el sesgo de muestra y proponiendo sobremuestreo o ponderación de la función de costo.}
	{Conclusiones del primer experimento: la generalización depende del sesgo de muestra de patrones poco frecuentes,
	el entrenamiento de reglas semitotalísticas no garantiza una regla aprendida semitotalística, y no se requieren
	arquitecturas complejas para replicar estas dinámicas. Se valida la metodología de extracción del vector de
	regla para los siguientes experimentos.}}
	{\realAnalisis\% (esperado: \espAnalisis\%)}{}

	\subsection*{Fase 4: Documentación y Difusión}

	\actividad{Documentación técnica}
	{Fase 4}{Oct/26 -- Dic/26}
	{No aplica en el periodo.}
	{Reporte técnico de TT completo (204 páginas).}
	{Actividad programada a partir de octubre de 2026. El reporte técnico se ha redactado de forma incremental
	conforme avanza cada actividad, por lo que su integración final partirá del documento adjunto.}
	{No iniciada (0\% esperado al corte)}{}

	\actividad{Artículo académico}
	{Fase 4}{Oct/26 -- Dic/26}
	{No aplica en el periodo.}
	{--}
	{Actividad programada a partir de octubre de 2026.}
	{No iniciada (0\% esperado al corte)}{}

	\section{Cronograma}

	\resizebox{\linewidth}{!}{%
		\begin{tabular}{p{5cm}*{10}{>{\centering\arraybackslash}p{1cm}}}
			\toprule
			\textbf{Actividad} & \textbf{Ago/25} & \textbf{Sep/25} & \textbf{Oct/25} & \textbf{Nov/25} & \textbf{Dic/25} & \mesActual{1}{Ago/26} & \mesActual{2}{Sep/26} & \textbf{Oct/26} & \textbf{Nov/26} & \textbf{Dic/26} \\
			\midrule
			\multicolumn{11}{l}{\textbf{Fase 1: Planificación y Revisión Bibliográfica}} \\

			Estado del arte       & \cellcolor{gray!30} & \cellcolor{gray!30} & \cellcolor{gray!30} &  &  &  &  &  &  &   \\

			Marco teórico y metodológico &  & \cellcolor{gray!30} & \cellcolor{gray!30} &  &  &  &  &  &  &   \\

			\midrule

			\multicolumn{11}{l}{\textbf{Fase 2: Diseño e implementación del simulador}} \\
			Arquitectura del simulador &  &  & \cellcolor{gray!30} & \cellcolor{gray!30} &  &  &  &  &  &   \\

			Módulos de simulación  &  &  &  & \cellcolor{gray!30} & \cellcolor{gray!30} &  &  &  &  &   \\

			Módulos entrenamiento   &  &  &  &  & \cellcolor{gray!30} & \cellcolor{gray!30} &  &  &  &   \\

			\midrule

			\multicolumn{11}{l}{\textbf{Fase 3: Entrenamiento y Validación}} \\
			Entrenamiento y ajuste   &  &  &  &  &  & \cellcolor{gray!30} & \cellcolor{gray!30} & \cellcolor{gray!30} &   &   \\

			Experimentación y análisis  &  &  &  &  &  & \cellcolor{gray!30}  & \cellcolor{gray!30} & \cellcolor{gray!30} & \cellcolor{gray!30} &   \\

			Análisis de resultados (global) &  &  &  &  &  & \cellcolor{gray!30} & \cellcolor{gray!30} & \cellcolor{gray!30} & \cellcolor{gray!30} &   \\

			\midrule

			\multicolumn{11}{l}{\textbf{Fase 4: Documentación y Difusión}} \\
			Documentación técnica  &  &  &  &  &  &  &  & \cellcolor{gray!30}  & \cellcolor{gray!30} & \cellcolor{gray!30}  \\

			Artículo académico  &  &  &  &  &  &  &  & \cellcolor{gray!30}  & \cellcolor{gray!30}  & \cellcolor{gray!30}  \\

			\bottomrule
		\end{tabular}
	}

	\vspace{0.5em}
	{\small La columna resaltada corresponde al mes reportado.}

	\section{Alumnos y directores}

	\vspace{1em}
	\begin{tabular}{@{} p{0.45\textwidth} p{0.45\textwidth} @{}}
		\textit{Erick Jesse Angeles López} & \textit{Dr. Genaro Juárez Martínez} \\[5em]
		Firma: \rule{6cm}{0.4pt} & Firma: \rule{6cm}{0.4pt} \\
	\end{tabular}

\end{document}
"
%  o simplemente:  ./compilar.sh
%
%  Las páginas y secciones citadas corresponden a Documentation/main.pdf
%  (reporte técnico de TT), que es el PDF que se adjunta como evidencia.
% ============================================================================
\documentclass[letterpaper,10pt,oneside]{article}
\usepackage[margin=2.54cm]{geometry}
\usepackage[spanish]{babel}
\newcommand\authormark[1]{\textsuperscript{#1}}
\usepackage[colorlinks=true,bookmarks=true,linkcolor=black,citecolor=black,urlcolor=blue]{hyperref}
\usepackage{times}
\usepackage[table]{xcolor}
\usepackage{booktabs}
\usepackage{graphicx}
\usepackage{titlesec}
\usepackage{tabularx}
\usepackage{array}
\usepackage{enumitem}
\usepackage{float}

\titleformat{\section}
{\normalfont\fontsize{12pt}{14pt}\bfseries}
{\thesection}{1em}{}

% ----------------------------------------------------------------------------
%  Mes reportado
% ----------------------------------------------------------------------------
\providecommand\corte{2}
\ifnum\corte=1
	\def\numReporte{1}
	\def\mes{Agosto}
	\def\fechaCorte{31 de agosto de 2026}
\else
	\def\numReporte{2}
	\def\mes{Septiembre}
	\def\fechaCorte{30 de septiembre de 2026}
\fi
% \segun{texto agosto}{texto septiembre}
\newcommand\segun[2]{\ifnum\corte=1 #1\else #2\fi}

% ----------------------------------------------------------------------------
%  Porcentajes (editar aquí). Formato: \segun{agosto}{septiembre}
%  "esp" = avance esperado según cronograma al corte, "real" = avance logrado.
% ----------------------------------------------------------------------------
\def\espEntrenamiento{\segun{33}{67}}   \def\realEntrenamiento{\segun{33}{67}}
\def\espExperimentos{\segun{25}{50}}    \def\realExperimentos{\segun{25}{50}}
\def\espAnalisis{\segun{25}{50}}        \def\realAnalisis{\segun{25}{50}}

% ----------------------------------------------------------------------------
%  Enlaces a video (dejar vacío si no aplica)
% ----------------------------------------------------------------------------
\def\videoSimulacion{}
\def\videoEntrenamiento{}

% ----------------------------------------------------------------------------
%  Bloque de evidencia por actividad
%  \actividad{nombre}{fase}{periodo}{tipo de evidencia}{ubicación en reporte TT}
%            {descripción}{porcentaje}{enlace de video}
% ----------------------------------------------------------------------------
\definecolor{encabezado}{gray}{0.85}
\newcommand\video[1]{\edef\tmp{#1}\ifx\tmp\empty No aplica.\else\url{#1}\fi}
\newcommand\actividad[8]{%
	\par\noindent
	\begin{minipage}{\linewidth}
		\renewcommand{\arraystretch}{1.25}
		\begin{tabularx}{\linewidth}{|>{\raggedright\bfseries\arraybackslash}p{0.24\linewidth}|X|}
			\hline
			\rowcolor{encabezado}
			Nombre de la actividad & \textbf{#1} \\ \hline
			Fase y periodo & #2 \,(#3) \\ \hline
			Evidencia de trabajo & #4 \\ \hline
			Ubicación en el reporte de TT & #5 \\ \hline
			Descripción del avance & #6 \\ \hline
			Porcentaje de cumplimiento & \textbf{#7} \\ \hline
			Enlace de video & \video{#8} \\ \hline
		\end{tabularx}
	\end{minipage}
	\par\vspace{1.2em}%
}

% Porcentaje de cumplimiento: real respecto a lo esperado al corte
\newcommand\cumplimiento[2]{\fpeval{min(100, round(100*(#2)/(#1), 0))}\%}

% Columna del mes reportado en el cronograma
\newcommand\mesActual[2]{\ifnum\corte=#1\cellcolor{black!70}\textcolor{white}{\textbf{#2}}\else\textbf{#2}\fi}

\begin{document}

	\begin{center}
		\fontsize{14}{16}\selectfont\bfseries
		\numReporte° Reporte de seguimiento (\mes{} 2026): \\
		Simulación de sistemas complejos mediante autómatas celulares neuronales
	\end{center}

	\begin{center}
		\fontsize{14}{16}\selectfont\bfseries
		\textit{Trabajo Terminal No. 2026-A114}
	\end{center}

	\begin{center}
		\fontsize{12}{16}\selectfont
		\textit{Alumno: \authormark{*}Angeles López Erick Jesse\authormark{1}} \\
		\textit{Director: Juárez Martínez Genaro\authormark{1}} \\
		\textit{\authormark{*}e-mail: eangelesl1700@alumno.ipn.mx}
	\end{center}

	\section{Información general}

	\begin{itemize}[leftmargin=*]
		\item \textbf{Periodo reportado:} \mes{} 2026 (corte al \fechaCorte).
		\item \textbf{Documento de evidencia adjunto:} reporte técnico de Trabajo Terminal
		\textit{``Simulación de sistemas complejos mediante autómatas celulares neuronales''} (PDF, 204 páginas).
		El avance es incremental: cada evidencia listada en este reporte se encuentra integrada en el capítulo
		y sección correspondiente de dicho documento, cuyas páginas se indican en cada actividad.
		\item \textbf{Actividades programadas en el mes:} \segun{Módulos de entrenamiento (cierre), Entrenamiento y ajuste,
		Experimentación y análisis, Análisis de resultados (global).}{Entrenamiento y ajuste, Experimentación y análisis,
		Análisis de resultados (global).}
	\end{itemize}

	\section{Resumen de cumplimiento}

	\begin{table}[H]
		\centering
		\renewcommand{\arraystretch}{1.2}
		\begin{tabular}{l c c c c}
			\toprule
			\textbf{Actividad} & \textbf{Periodo} & \textbf{Esperado al corte} & \textbf{Avance real} & \textbf{Cumplimiento} \\
			\midrule
			\multicolumn{5}{l}{\textit{Fase 1: Planificación y Revisión Bibliográfica}} \\
			Estado del arte                  & Ago/25 -- Oct/25 & 100\% & 100\% & 100\% \\
			Marco teórico y metodológico     & Sep/25 -- Oct/25 & 100\% & 100\% & 100\% \\
			\multicolumn{5}{l}{\textit{Fase 2: Diseño e implementación del simulador}} \\
			Arquitectura del simulador       & Oct/25 -- Nov/25 & 100\% & 100\% & 100\% \\
			Módulos de simulación            & Nov/25 -- Dic/25 & 100\% & 100\% & 100\% \\
			Módulos entrenamiento            & Dic/25 -- Ago/26 & 100\% & 100\% & 100\% \\
			\multicolumn{5}{l}{\textit{Fase 3: Entrenamiento y Validación}} \\
			Entrenamiento y ajuste           & Ago/26 -- Oct/26 & \espEntrenamiento\% & \realEntrenamiento\% & \cumplimiento{\espEntrenamiento}{\realEntrenamiento} \\
			Experimentación y análisis       & Ago/26 -- Nov/26 & \espExperimentos\%  & \realExperimentos\%  & \cumplimiento{\espExperimentos}{\realExperimentos} \\
			Análisis de resultados (global)  & Ago/26 -- Nov/26 & \espAnalisis\%      & \realAnalisis\%      & \cumplimiento{\espAnalisis}{\realAnalisis} \\
			\multicolumn{5}{l}{\textit{Fase 4: Documentación y Difusión}} \\
			Documentación técnica            & Oct/26 -- Dic/26 & 0\% & -- & No iniciada \\
			Artículo académico               & Oct/26 -- Dic/26 & 0\% & -- & No iniciada \\
			\bottomrule
		\end{tabular}
		\caption{Avance acumulado por actividad al \fechaCorte. ``Esperado al corte'' es la proporción del periodo
		programado en el cronograma que ha transcurrido; ``Cumplimiento'' compara el avance real con dicho valor.}
	\end{table}

	\section{Evidencias por actividad}

	\subsection*{Fase 1: Planificación y Revisión Bibliográfica}

	\actividad{Estado del arte}
	{Fase 1}{Ago/25 -- Oct/25}
	{Documento escrito.}
	{Capítulo 5, sección 5.1 \textit{Estado del arte, aplicaciones y líneas actuales de investigación} (pp.~61--77);
	Capítulo 1 \textit{Introducción} (pp.~1--6).}
	{Revisión de trabajos sobre autómatas celulares neuronales: morfogénesis, NCA isotrópicos, autoclasificación,
	representaciones tridimensionales, NCA en malla, modelos generativos y aplicaciones adicionales.
	Actividad concluida en Trabajo Terminal I.}
	{100\%}{}

	\actividad{Marco teórico y metodológico}
	{Fase 1}{Sep/25 -- Oct/25}
	{Documento escrito.}
	{Capítulos 2 a 5 (marco teórico, pp.~7--77) y Capítulo 6 \textit{Marco metodológico} (pp.~78--90).}
	{Marco teórico de sistemas complejos y teoría del caos, autómatas celulares, inteligencia artificial y redes
	neuronales, y autómatas celulares neuronales. Marco metodológico con diseño cuasiexperimental, población,
	muestra, técnicas de recolección y procesamiento de datos, y contextualización tecnológica.
	Actividad concluida en Trabajo Terminal I.}
	{100\%}{}

	\subsection*{Fase 2: Diseño e implementación del simulador}

	\actividad{Arquitectura del simulador}
	{Fase 2}{Oct/25 -- Nov/25}
	{Documento escrito (análisis y diseño de software, diagramas UML).}
	{Capítulo 7 \textit{Análisis de software} (pp.~91--135); Capítulo 8, secciones 8.1 \textit{Arquitectura},
	8.2 \textit{Diagrama de clases}, 8.3 \textit{Diseño de datos} y 8.5 \textit{Interfaces de usuarios} (pp.~136--143, 153--156).}
	{Metodología de desarrollo, reglas de negocio, requerimientos, análisis de riesgos, estudio de factibilidad,
	casos de uso con sus especificaciones, diagramas de secuencia, arquitectura MVC, diagrama de clases,
	diseño de datos e interfaces gráficas.}
	{100\%}{}

	\actividad{Módulos de simulación}
	{Fase 2}{Nov/25 -- Dic/25}
	{Documento escrito, tablas de resultados y pruebas de operatividad.}
	{Sección 8.4 \textit{Diseño de algoritmos} (8.4.1--8.4.4 y 8.4.6, pp.~144--152); secciones 9.1 \textit{Entorno de desarrollo}
	y 9.2 \textit{Estructura del proyecto} (pp.~157--162); subsecciones 9.4.1, 9.4.3, 9.4.4 y 9.4.5 de
	\textit{Validación y pruebas} (pp.~166--169), tablas 9.6 y 9.7.}
	{Núcleo de simulación en C++/CUDA con enlaces a Python e interfaz en Qt: indexación de la regla semitotalística,
	paso de evolución, condiciones de frontera, \textit{forward pass} del NCA y \textit{prefetch} asíncrono.
	Validación funcional contra reglas conocidas y pruebas de rendimiento (hasta 345\,000 pasos/s en
	$100\times100$ y más de 100 fps de renderizado en $2000\times2000$).}
	{100\%}{\videoSimulacion}

	\actividad{Módulos entrenamiento}
	{Fase 2}{Dic/25 -- Ago/26}
	{Documento escrito, pruebas de operatividad.}
	{Subsección 7.8.4 \textit{Especificaciones de casos de uso: Entrenamiento} (pp.~125--130); 8.3.5
	\textit{Configuración del NCA} (p.~143); 8.4.5 \textit{Entrenamiento del autómata celular neuronal} (pp.~150--151);
	sección 9.3 \textit{Entrenamiento del modelo} (pp.~163--166, tabla 9.5); subsecciones 9.4.2 y 9.4.6 (pp.~167, 169).}
	{Módulo de entrenamiento en PyTorch con construcción dinámica del perceptrón multicapa, ciclo de optimización
	con Adam y detención temprana, persistencia de pesos en \texttt{.npz} compatible con el \textit{kernel} CUDA e
	integración con la pestaña de entrenamiento de la GUI (subproceso con \textit{log} en tiempo real).
	Se agregó el parámetro de semilla aleatoria para reproducir experimentos.
	\segun{Actividad concluida en el mes reportado.}{Actividad concluida en agosto de 2026.}}
	{100\%}{\videoEntrenamiento}

	\subsection*{Fase 3: Entrenamiento y Validación}

	\actividad{Entrenamiento y ajuste}
	{Fase 3}{Ago/26 -- Oct/26}
	{Documento escrito y tablas de resultados.}
	{Capítulo 10, sección 10.1 \textit{Duplicación de reglas del juego de la vida}:
	\segun{tablas 10.1 y 10.2 y figura 10.3 (pp.~173--174).}
	{tablas 10.1 y 10.2 y figura 10.3 (pp.~173--174); figuras 10.5 a 10.9 (pp.~176--178).}}
	{\segun{Barrido de hiperparámetros para la regla B3/S23: 10 combinaciones de funciones de activación con
	una capa oculta de 4 neuronas (una sola converge) y con 4 capas ocultas de 8 neuronas (cinco convergen a una
	pérdida menor a $10^{-4}$). Se identifica el estancamiento en las primeras 12 épocas.}
	{Además del barrido de hiperparámetros de agosto, se entrenaron las reglas B3/S23, B3/S238, B38/S23, B3/S236
	y B38/S238 con distintas pérdidas objetivo y semillas, midiendo la distancia Manhattan entre el vector de regla
	aprendido y el original. Se determinó que la pérdida objetivo debe ser $\leq 0.1$ para B3/S23.}}
	{\realEntrenamiento\% (esperado: \espEntrenamiento\%)}{}

	\actividad{Experimentación y análisis}
	{Fase 3}{Ago/26 -- Nov/26}
	{Documento escrito, tablas y gráficas de resultados.}
	{Capítulo 10, sección 10.1:
	\segun{figuras 10.1 y 10.2 (pp.~170--172).}
	{figuras 10.1 a 10.13 (pp.~170--181).}}
	{\segun{Generación del conjunto de datos (100 condiciones iniciales en $100\times100$ durante 100 pasos) y
	análisis de su balance: el estado 1 aparece un $16.5\% \pm 4.2\%$ y los patrones $C0N8$ y $C1N8$ representan
	menos del $0.02\%$ de las entradas.}
	{Primer experimento (duplicación de reglas del juego de la vida) completo. Se diseñó un espacio con las
	512 permutaciones de la vecindad para extraer el vector de regla aprendido y verificar si el NCA conserva la
	propiedad semitotalística. Se identificaron reglas aprendidas alternativas (B3/S23 en lugar de B3/S238,
	B3/S23678 en lugar de B3/S236) y reglas no semitotalísticas.}}
	{\realExperimentos\% (esperado: \espExperimentos\%)}{}

	\actividad{Análisis de resultados (global)}
	{Fase 3}{Ago/26 -- Nov/26}
	{Documento escrito.}
	{Capítulo 10, sección 10.1\segun{ (pp.~170--172).}{, conclusiones del experimento (p.~181).}}
	{\segun{Análisis estadístico de la distribución de estados y de patrones de vecindad para la regla B3/S23,
	identificando el sesgo de muestra y proponiendo sobremuestreo o ponderación de la función de costo.}
	{Conclusiones del primer experimento: la generalización depende del sesgo de muestra de patrones poco frecuentes,
	el entrenamiento de reglas semitotalísticas no garantiza una regla aprendida semitotalística, y no se requieren
	arquitecturas complejas para replicar estas dinámicas. Se valida la metodología de extracción del vector de
	regla para los siguientes experimentos.}}
	{\realAnalisis\% (esperado: \espAnalisis\%)}{}

	\subsection*{Fase 4: Documentación y Difusión}

	\actividad{Documentación técnica}
	{Fase 4}{Oct/26 -- Dic/26}
	{No aplica en el periodo.}
	{Reporte técnico de TT completo (204 páginas).}
	{Actividad programada a partir de octubre de 2026. El reporte técnico se ha redactado de forma incremental
	conforme avanza cada actividad, por lo que su integración final partirá del documento adjunto.}
	{No iniciada (0\% esperado al corte)}{}

	\actividad{Artículo académico}
	{Fase 4}{Oct/26 -- Dic/26}
	{No aplica en el periodo.}
	{--}
	{Actividad programada a partir de octubre de 2026.}
	{No iniciada (0\% esperado al corte)}{}

	\section{Cronograma}

	\resizebox{\linewidth}{!}{%
		\begin{tabular}{p{5cm}*{10}{>{\centering\arraybackslash}p{1cm}}}
			\toprule
			\textbf{Actividad} & \textbf{Ago/25} & \textbf{Sep/25} & \textbf{Oct/25} & \textbf{Nov/25} & \textbf{Dic/25} & \mesActual{1}{Ago/26} & \mesActual{2}{Sep/26} & \textbf{Oct/26} & \textbf{Nov/26} & \textbf{Dic/26} \\
			\midrule
			\multicolumn{11}{l}{\textbf{Fase 1: Planificación y Revisión Bibliográfica}} \\

			Estado del arte       & \cellcolor{gray!30} & \cellcolor{gray!30} & \cellcolor{gray!30} &  &  &  &  &  &  &   \\

			Marco teórico y metodológico &  & \cellcolor{gray!30} & \cellcolor{gray!30} &  &  &  &  &  &  &   \\

			\midrule

			\multicolumn{11}{l}{\textbf{Fase 2: Diseño e implementación del simulador}} \\
			Arquitectura del simulador &  &  & \cellcolor{gray!30} & \cellcolor{gray!30} &  &  &  &  &  &   \\

			Módulos de simulación  &  &  &  & \cellcolor{gray!30} & \cellcolor{gray!30} &  &  &  &  &   \\

			Módulos entrenamiento   &  &  &  &  & \cellcolor{gray!30} & \cellcolor{gray!30} &  &  &  &   \\

			\midrule

			\multicolumn{11}{l}{\textbf{Fase 3: Entrenamiento y Validación}} \\
			Entrenamiento y ajuste   &  &  &  &  &  & \cellcolor{gray!30} & \cellcolor{gray!30} & \cellcolor{gray!30} &   &   \\

			Experimentación y análisis  &  &  &  &  &  & \cellcolor{gray!30}  & \cellcolor{gray!30} & \cellcolor{gray!30} & \cellcolor{gray!30} &   \\

			Análisis de resultados (global) &  &  &  &  &  & \cellcolor{gray!30} & \cellcolor{gray!30} & \cellcolor{gray!30} & \cellcolor{gray!30} &   \\

			\midrule

			\multicolumn{11}{l}{\textbf{Fase 4: Documentación y Difusión}} \\
			Documentación técnica  &  &  &  &  &  &  &  & \cellcolor{gray!30}  & \cellcolor{gray!30} & \cellcolor{gray!30}  \\

			Artículo académico  &  &  &  &  &  &  &  & \cellcolor{gray!30}  & \cellcolor{gray!30}  & \cellcolor{gray!30}  \\

			\bottomrule
		\end{tabular}
	}

	\vspace{0.5em}
	{\small La columna resaltada corresponde al mes reportado.}

	\section{Alumnos y directores}

	\vspace{1em}
	\begin{tabular}{@{} p{0.45\textwidth} p{0.45\textwidth} @{}}
		\textit{Erick Jesse Angeles López} & \textit{Dr. Genaro Juárez Martínez} \\[5em]
		Firma: \rule{6cm}{0.4pt} & Firma: \rule{6cm}{0.4pt} \\
	\end{tabular}

\end{document}
"
%      pdflatex -jobname=2ReporteSeguimiento "\def\corte{2}% ============================================================================
%  Reporte de seguimiento - Trabajo Terminal II
%
%  Un solo archivo genera los reportes de cada mes. El mes se elige con \corte:
%      1 = Agosto 2026, 2 = Septiembre 2026
%
%  Compilación (desde esta carpeta):
%      pdflatex -jobname=1ReporteSeguimiento "\def\corte{1}% ============================================================================
%  Reporte de seguimiento - Trabajo Terminal II
%
%  Un solo archivo genera los reportes de cada mes. El mes se elige con \corte:
%      1 = Agosto 2026, 2 = Septiembre 2026
%
%  Compilación (desde esta carpeta):
%      pdflatex -jobname=1ReporteSeguimiento "\def\corte{1}% ============================================================================
%  Reporte de seguimiento - Trabajo Terminal II
%
%  Un solo archivo genera los reportes de cada mes. El mes se elige con \corte:
%      1 = Agosto 2026, 2 = Septiembre 2026
%
%  Compilación (desde esta carpeta):
%      pdflatex -jobname=1ReporteSeguimiento "\def\corte{1}\input{reporteSeguimiento}"
%      pdflatex -jobname=2ReporteSeguimiento "\def\corte{2}\input{reporteSeguimiento}"
%  o simplemente:  ./compilar.sh
%
%  Las páginas y secciones citadas corresponden a Documentation/main.pdf
%  (reporte técnico de TT), que es el PDF que se adjunta como evidencia.
% ============================================================================
\documentclass[letterpaper,10pt,oneside]{article}
\usepackage[margin=2.54cm]{geometry}
\usepackage[spanish]{babel}
\newcommand\authormark[1]{\textsuperscript{#1}}
\usepackage[colorlinks=true,bookmarks=true,linkcolor=black,citecolor=black,urlcolor=blue]{hyperref}
\usepackage{times}
\usepackage[table]{xcolor}
\usepackage{booktabs}
\usepackage{graphicx}
\usepackage{titlesec}
\usepackage{tabularx}
\usepackage{array}
\usepackage{enumitem}
\usepackage{float}

\titleformat{\section}
{\normalfont\fontsize{12pt}{14pt}\bfseries}
{\thesection}{1em}{}

% ----------------------------------------------------------------------------
%  Mes reportado
% ----------------------------------------------------------------------------
\providecommand\corte{2}
\ifnum\corte=1
	\def\numReporte{1}
	\def\mes{Agosto}
	\def\fechaCorte{31 de agosto de 2026}
\else
	\def\numReporte{2}
	\def\mes{Septiembre}
	\def\fechaCorte{30 de septiembre de 2026}
\fi
% \segun{texto agosto}{texto septiembre}
\newcommand\segun[2]{\ifnum\corte=1 #1\else #2\fi}

% ----------------------------------------------------------------------------
%  Porcentajes (editar aquí). Formato: \segun{agosto}{septiembre}
%  "esp" = avance esperado según cronograma al corte, "real" = avance logrado.
% ----------------------------------------------------------------------------
\def\espEntrenamiento{\segun{33}{67}}   \def\realEntrenamiento{\segun{33}{67}}
\def\espExperimentos{\segun{25}{50}}    \def\realExperimentos{\segun{25}{50}}
\def\espAnalisis{\segun{25}{50}}        \def\realAnalisis{\segun{25}{50}}

% ----------------------------------------------------------------------------
%  Enlaces a video (dejar vacío si no aplica)
% ----------------------------------------------------------------------------
\def\videoSimulacion{}
\def\videoEntrenamiento{}

% ----------------------------------------------------------------------------
%  Bloque de evidencia por actividad
%  \actividad{nombre}{fase}{periodo}{tipo de evidencia}{ubicación en reporte TT}
%            {descripción}{porcentaje}{enlace de video}
% ----------------------------------------------------------------------------
\definecolor{encabezado}{gray}{0.85}
\newcommand\video[1]{\edef\tmp{#1}\ifx\tmp\empty No aplica.\else\url{#1}\fi}
\newcommand\actividad[8]{%
	\par\noindent
	\begin{minipage}{\linewidth}
		\renewcommand{\arraystretch}{1.25}
		\begin{tabularx}{\linewidth}{|>{\raggedright\bfseries\arraybackslash}p{0.24\linewidth}|X|}
			\hline
			\rowcolor{encabezado}
			Nombre de la actividad & \textbf{#1} \\ \hline
			Fase y periodo & #2 \,(#3) \\ \hline
			Evidencia de trabajo & #4 \\ \hline
			Ubicación en el reporte de TT & #5 \\ \hline
			Descripción del avance & #6 \\ \hline
			Porcentaje de cumplimiento & \textbf{#7} \\ \hline
			Enlace de video & \video{#8} \\ \hline
		\end{tabularx}
	\end{minipage}
	\par\vspace{1.2em}%
}

% Porcentaje de cumplimiento: real respecto a lo esperado al corte
\newcommand\cumplimiento[2]{\fpeval{min(100, round(100*(#2)/(#1), 0))}\%}

% Columna del mes reportado en el cronograma
\newcommand\mesActual[2]{\ifnum\corte=#1\cellcolor{black!70}\textcolor{white}{\textbf{#2}}\else\textbf{#2}\fi}

\begin{document}

	\begin{center}
		\fontsize{14}{16}\selectfont\bfseries
		\numReporte° Reporte de seguimiento (\mes{} 2026): \\
		Simulación de sistemas complejos mediante autómatas celulares neuronales
	\end{center}

	\begin{center}
		\fontsize{14}{16}\selectfont\bfseries
		\textit{Trabajo Terminal No. 2026-A114}
	\end{center}

	\begin{center}
		\fontsize{12}{16}\selectfont
		\textit{Alumno: \authormark{*}Angeles López Erick Jesse\authormark{1}} \\
		\textit{Director: Juárez Martínez Genaro\authormark{1}} \\
		\textit{\authormark{*}e-mail: eangelesl1700@alumno.ipn.mx}
	\end{center}

	\section{Información general}

	\begin{itemize}[leftmargin=*]
		\item \textbf{Periodo reportado:} \mes{} 2026 (corte al \fechaCorte).
		\item \textbf{Documento de evidencia adjunto:} reporte técnico de Trabajo Terminal
		\textit{``Simulación de sistemas complejos mediante autómatas celulares neuronales''} (PDF, 204 páginas).
		El avance es incremental: cada evidencia listada en este reporte se encuentra integrada en el capítulo
		y sección correspondiente de dicho documento, cuyas páginas se indican en cada actividad.
		\item \textbf{Actividades programadas en el mes:} \segun{Módulos de entrenamiento (cierre), Entrenamiento y ajuste,
		Experimentación y análisis, Análisis de resultados (global).}{Entrenamiento y ajuste, Experimentación y análisis,
		Análisis de resultados (global).}
	\end{itemize}

	\section{Resumen de cumplimiento}

	\begin{table}[H]
		\centering
		\renewcommand{\arraystretch}{1.2}
		\begin{tabular}{l c c c c}
			\toprule
			\textbf{Actividad} & \textbf{Periodo} & \textbf{Esperado al corte} & \textbf{Avance real} & \textbf{Cumplimiento} \\
			\midrule
			\multicolumn{5}{l}{\textit{Fase 1: Planificación y Revisión Bibliográfica}} \\
			Estado del arte                  & Ago/25 -- Oct/25 & 100\% & 100\% & 100\% \\
			Marco teórico y metodológico     & Sep/25 -- Oct/25 & 100\% & 100\% & 100\% \\
			\multicolumn{5}{l}{\textit{Fase 2: Diseño e implementación del simulador}} \\
			Arquitectura del simulador       & Oct/25 -- Nov/25 & 100\% & 100\% & 100\% \\
			Módulos de simulación            & Nov/25 -- Dic/25 & 100\% & 100\% & 100\% \\
			Módulos entrenamiento            & Dic/25 -- Ago/26 & 100\% & 100\% & 100\% \\
			\multicolumn{5}{l}{\textit{Fase 3: Entrenamiento y Validación}} \\
			Entrenamiento y ajuste           & Ago/26 -- Oct/26 & \espEntrenamiento\% & \realEntrenamiento\% & \cumplimiento{\espEntrenamiento}{\realEntrenamiento} \\
			Experimentación y análisis       & Ago/26 -- Nov/26 & \espExperimentos\%  & \realExperimentos\%  & \cumplimiento{\espExperimentos}{\realExperimentos} \\
			Análisis de resultados (global)  & Ago/26 -- Nov/26 & \espAnalisis\%      & \realAnalisis\%      & \cumplimiento{\espAnalisis}{\realAnalisis} \\
			\multicolumn{5}{l}{\textit{Fase 4: Documentación y Difusión}} \\
			Documentación técnica            & Oct/26 -- Dic/26 & 0\% & -- & No iniciada \\
			Artículo académico               & Oct/26 -- Dic/26 & 0\% & -- & No iniciada \\
			\bottomrule
		\end{tabular}
		\caption{Avance acumulado por actividad al \fechaCorte. ``Esperado al corte'' es la proporción del periodo
		programado en el cronograma que ha transcurrido; ``Cumplimiento'' compara el avance real con dicho valor.}
	\end{table}

	\section{Evidencias por actividad}

	\subsection*{Fase 1: Planificación y Revisión Bibliográfica}

	\actividad{Estado del arte}
	{Fase 1}{Ago/25 -- Oct/25}
	{Documento escrito.}
	{Capítulo 5, sección 5.1 \textit{Estado del arte, aplicaciones y líneas actuales de investigación} (pp.~61--77);
	Capítulo 1 \textit{Introducción} (pp.~1--6).}
	{Revisión de trabajos sobre autómatas celulares neuronales: morfogénesis, NCA isotrópicos, autoclasificación,
	representaciones tridimensionales, NCA en malla, modelos generativos y aplicaciones adicionales.
	Actividad concluida en Trabajo Terminal I.}
	{100\%}{}

	\actividad{Marco teórico y metodológico}
	{Fase 1}{Sep/25 -- Oct/25}
	{Documento escrito.}
	{Capítulos 2 a 5 (marco teórico, pp.~7--77) y Capítulo 6 \textit{Marco metodológico} (pp.~78--90).}
	{Marco teórico de sistemas complejos y teoría del caos, autómatas celulares, inteligencia artificial y redes
	neuronales, y autómatas celulares neuronales. Marco metodológico con diseño cuasiexperimental, población,
	muestra, técnicas de recolección y procesamiento de datos, y contextualización tecnológica.
	Actividad concluida en Trabajo Terminal I.}
	{100\%}{}

	\subsection*{Fase 2: Diseño e implementación del simulador}

	\actividad{Arquitectura del simulador}
	{Fase 2}{Oct/25 -- Nov/25}
	{Documento escrito (análisis y diseño de software, diagramas UML).}
	{Capítulo 7 \textit{Análisis de software} (pp.~91--135); Capítulo 8, secciones 8.1 \textit{Arquitectura},
	8.2 \textit{Diagrama de clases}, 8.3 \textit{Diseño de datos} y 8.5 \textit{Interfaces de usuarios} (pp.~136--143, 153--156).}
	{Metodología de desarrollo, reglas de negocio, requerimientos, análisis de riesgos, estudio de factibilidad,
	casos de uso con sus especificaciones, diagramas de secuencia, arquitectura MVC, diagrama de clases,
	diseño de datos e interfaces gráficas.}
	{100\%}{}

	\actividad{Módulos de simulación}
	{Fase 2}{Nov/25 -- Dic/25}
	{Documento escrito, tablas de resultados y pruebas de operatividad.}
	{Sección 8.4 \textit{Diseño de algoritmos} (8.4.1--8.4.4 y 8.4.6, pp.~144--152); secciones 9.1 \textit{Entorno de desarrollo}
	y 9.2 \textit{Estructura del proyecto} (pp.~157--162); subsecciones 9.4.1, 9.4.3, 9.4.4 y 9.4.5 de
	\textit{Validación y pruebas} (pp.~166--169), tablas 9.6 y 9.7.}
	{Núcleo de simulación en C++/CUDA con enlaces a Python e interfaz en Qt: indexación de la regla semitotalística,
	paso de evolución, condiciones de frontera, \textit{forward pass} del NCA y \textit{prefetch} asíncrono.
	Validación funcional contra reglas conocidas y pruebas de rendimiento (hasta 345\,000 pasos/s en
	$100\times100$ y más de 100 fps de renderizado en $2000\times2000$).}
	{100\%}{\videoSimulacion}

	\actividad{Módulos entrenamiento}
	{Fase 2}{Dic/25 -- Ago/26}
	{Documento escrito, pruebas de operatividad.}
	{Subsección 7.8.4 \textit{Especificaciones de casos de uso: Entrenamiento} (pp.~125--130); 8.3.5
	\textit{Configuración del NCA} (p.~143); 8.4.5 \textit{Entrenamiento del autómata celular neuronal} (pp.~150--151);
	sección 9.3 \textit{Entrenamiento del modelo} (pp.~163--166, tabla 9.5); subsecciones 9.4.2 y 9.4.6 (pp.~167, 169).}
	{Módulo de entrenamiento en PyTorch con construcción dinámica del perceptrón multicapa, ciclo de optimización
	con Adam y detención temprana, persistencia de pesos en \texttt{.npz} compatible con el \textit{kernel} CUDA e
	integración con la pestaña de entrenamiento de la GUI (subproceso con \textit{log} en tiempo real).
	Se agregó el parámetro de semilla aleatoria para reproducir experimentos.
	\segun{Actividad concluida en el mes reportado.}{Actividad concluida en agosto de 2026.}}
	{100\%}{\videoEntrenamiento}

	\subsection*{Fase 3: Entrenamiento y Validación}

	\actividad{Entrenamiento y ajuste}
	{Fase 3}{Ago/26 -- Oct/26}
	{Documento escrito y tablas de resultados.}
	{Capítulo 10, sección 10.1 \textit{Duplicación de reglas del juego de la vida}:
	\segun{tablas 10.1 y 10.2 y figura 10.3 (pp.~173--174).}
	{tablas 10.1 y 10.2 y figura 10.3 (pp.~173--174); figuras 10.5 a 10.9 (pp.~176--178).}}
	{\segun{Barrido de hiperparámetros para la regla B3/S23: 10 combinaciones de funciones de activación con
	una capa oculta de 4 neuronas (una sola converge) y con 4 capas ocultas de 8 neuronas (cinco convergen a una
	pérdida menor a $10^{-4}$). Se identifica el estancamiento en las primeras 12 épocas.}
	{Además del barrido de hiperparámetros de agosto, se entrenaron las reglas B3/S23, B3/S238, B38/S23, B3/S236
	y B38/S238 con distintas pérdidas objetivo y semillas, midiendo la distancia Manhattan entre el vector de regla
	aprendido y el original. Se determinó que la pérdida objetivo debe ser $\leq 0.1$ para B3/S23.}}
	{\realEntrenamiento\% (esperado: \espEntrenamiento\%)}{}

	\actividad{Experimentación y análisis}
	{Fase 3}{Ago/26 -- Nov/26}
	{Documento escrito, tablas y gráficas de resultados.}
	{Capítulo 10, sección 10.1:
	\segun{figuras 10.1 y 10.2 (pp.~170--172).}
	{figuras 10.1 a 10.13 (pp.~170--181).}}
	{\segun{Generación del conjunto de datos (100 condiciones iniciales en $100\times100$ durante 100 pasos) y
	análisis de su balance: el estado 1 aparece un $16.5\% \pm 4.2\%$ y los patrones $C0N8$ y $C1N8$ representan
	menos del $0.02\%$ de las entradas.}
	{Primer experimento (duplicación de reglas del juego de la vida) completo. Se diseñó un espacio con las
	512 permutaciones de la vecindad para extraer el vector de regla aprendido y verificar si el NCA conserva la
	propiedad semitotalística. Se identificaron reglas aprendidas alternativas (B3/S23 en lugar de B3/S238,
	B3/S23678 en lugar de B3/S236) y reglas no semitotalísticas.}}
	{\realExperimentos\% (esperado: \espExperimentos\%)}{}

	\actividad{Análisis de resultados (global)}
	{Fase 3}{Ago/26 -- Nov/26}
	{Documento escrito.}
	{Capítulo 10, sección 10.1\segun{ (pp.~170--172).}{, conclusiones del experimento (p.~181).}}
	{\segun{Análisis estadístico de la distribución de estados y de patrones de vecindad para la regla B3/S23,
	identificando el sesgo de muestra y proponiendo sobremuestreo o ponderación de la función de costo.}
	{Conclusiones del primer experimento: la generalización depende del sesgo de muestra de patrones poco frecuentes,
	el entrenamiento de reglas semitotalísticas no garantiza una regla aprendida semitotalística, y no se requieren
	arquitecturas complejas para replicar estas dinámicas. Se valida la metodología de extracción del vector de
	regla para los siguientes experimentos.}}
	{\realAnalisis\% (esperado: \espAnalisis\%)}{}

	\subsection*{Fase 4: Documentación y Difusión}

	\actividad{Documentación técnica}
	{Fase 4}{Oct/26 -- Dic/26}
	{No aplica en el periodo.}
	{Reporte técnico de TT completo (204 páginas).}
	{Actividad programada a partir de octubre de 2026. El reporte técnico se ha redactado de forma incremental
	conforme avanza cada actividad, por lo que su integración final partirá del documento adjunto.}
	{No iniciada (0\% esperado al corte)}{}

	\actividad{Artículo académico}
	{Fase 4}{Oct/26 -- Dic/26}
	{No aplica en el periodo.}
	{--}
	{Actividad programada a partir de octubre de 2026.}
	{No iniciada (0\% esperado al corte)}{}

	\section{Cronograma}

	\resizebox{\linewidth}{!}{%
		\begin{tabular}{p{5cm}*{10}{>{\centering\arraybackslash}p{1cm}}}
			\toprule
			\textbf{Actividad} & \textbf{Ago/25} & \textbf{Sep/25} & \textbf{Oct/25} & \textbf{Nov/25} & \textbf{Dic/25} & \mesActual{1}{Ago/26} & \mesActual{2}{Sep/26} & \textbf{Oct/26} & \textbf{Nov/26} & \textbf{Dic/26} \\
			\midrule
			\multicolumn{11}{l}{\textbf{Fase 1: Planificación y Revisión Bibliográfica}} \\

			Estado del arte       & \cellcolor{gray!30} & \cellcolor{gray!30} & \cellcolor{gray!30} &  &  &  &  &  &  &   \\

			Marco teórico y metodológico &  & \cellcolor{gray!30} & \cellcolor{gray!30} &  &  &  &  &  &  &   \\

			\midrule

			\multicolumn{11}{l}{\textbf{Fase 2: Diseño e implementación del simulador}} \\
			Arquitectura del simulador &  &  & \cellcolor{gray!30} & \cellcolor{gray!30} &  &  &  &  &  &   \\

			Módulos de simulación  &  &  &  & \cellcolor{gray!30} & \cellcolor{gray!30} &  &  &  &  &   \\

			Módulos entrenamiento   &  &  &  &  & \cellcolor{gray!30} & \cellcolor{gray!30} &  &  &  &   \\

			\midrule

			\multicolumn{11}{l}{\textbf{Fase 3: Entrenamiento y Validación}} \\
			Entrenamiento y ajuste   &  &  &  &  &  & \cellcolor{gray!30} & \cellcolor{gray!30} & \cellcolor{gray!30} &   &   \\

			Experimentación y análisis  &  &  &  &  &  & \cellcolor{gray!30}  & \cellcolor{gray!30} & \cellcolor{gray!30} & \cellcolor{gray!30} &   \\

			Análisis de resultados (global) &  &  &  &  &  & \cellcolor{gray!30} & \cellcolor{gray!30} & \cellcolor{gray!30} & \cellcolor{gray!30} &   \\

			\midrule

			\multicolumn{11}{l}{\textbf{Fase 4: Documentación y Difusión}} \\
			Documentación técnica  &  &  &  &  &  &  &  & \cellcolor{gray!30}  & \cellcolor{gray!30} & \cellcolor{gray!30}  \\

			Artículo académico  &  &  &  &  &  &  &  & \cellcolor{gray!30}  & \cellcolor{gray!30}  & \cellcolor{gray!30}  \\

			\bottomrule
		\end{tabular}
	}

	\vspace{0.5em}
	{\small La columna resaltada corresponde al mes reportado.}

	\section{Alumnos y directores}

	\vspace{1em}
	\begin{tabular}{@{} p{0.45\textwidth} p{0.45\textwidth} @{}}
		\textit{Erick Jesse Angeles López} & \textit{Dr. Genaro Juárez Martínez} \\[5em]
		Firma: \rule{6cm}{0.4pt} & Firma: \rule{6cm}{0.4pt} \\
	\end{tabular}

\end{document}
"
%      pdflatex -jobname=2ReporteSeguimiento "\def\corte{2}% ============================================================================
%  Reporte de seguimiento - Trabajo Terminal II
%
%  Un solo archivo genera los reportes de cada mes. El mes se elige con \corte:
%      1 = Agosto 2026, 2 = Septiembre 2026
%
%  Compilación (desde esta carpeta):
%      pdflatex -jobname=1ReporteSeguimiento "\def\corte{1}\input{reporteSeguimiento}"
%      pdflatex -jobname=2ReporteSeguimiento "\def\corte{2}\input{reporteSeguimiento}"
%  o simplemente:  ./compilar.sh
%
%  Las páginas y secciones citadas corresponden a Documentation/main.pdf
%  (reporte técnico de TT), que es el PDF que se adjunta como evidencia.
% ============================================================================
\documentclass[letterpaper,10pt,oneside]{article}
\usepackage[margin=2.54cm]{geometry}
\usepackage[spanish]{babel}
\newcommand\authormark[1]{\textsuperscript{#1}}
\usepackage[colorlinks=true,bookmarks=true,linkcolor=black,citecolor=black,urlcolor=blue]{hyperref}
\usepackage{times}
\usepackage[table]{xcolor}
\usepackage{booktabs}
\usepackage{graphicx}
\usepackage{titlesec}
\usepackage{tabularx}
\usepackage{array}
\usepackage{enumitem}
\usepackage{float}

\titleformat{\section}
{\normalfont\fontsize{12pt}{14pt}\bfseries}
{\thesection}{1em}{}

% ----------------------------------------------------------------------------
%  Mes reportado
% ----------------------------------------------------------------------------
\providecommand\corte{2}
\ifnum\corte=1
	\def\numReporte{1}
	\def\mes{Agosto}
	\def\fechaCorte{31 de agosto de 2026}
\else
	\def\numReporte{2}
	\def\mes{Septiembre}
	\def\fechaCorte{30 de septiembre de 2026}
\fi
% \segun{texto agosto}{texto septiembre}
\newcommand\segun[2]{\ifnum\corte=1 #1\else #2\fi}

% ----------------------------------------------------------------------------
%  Porcentajes (editar aquí). Formato: \segun{agosto}{septiembre}
%  "esp" = avance esperado según cronograma al corte, "real" = avance logrado.
% ----------------------------------------------------------------------------
\def\espEntrenamiento{\segun{33}{67}}   \def\realEntrenamiento{\segun{33}{67}}
\def\espExperimentos{\segun{25}{50}}    \def\realExperimentos{\segun{25}{50}}
\def\espAnalisis{\segun{25}{50}}        \def\realAnalisis{\segun{25}{50}}

% ----------------------------------------------------------------------------
%  Enlaces a video (dejar vacío si no aplica)
% ----------------------------------------------------------------------------
\def\videoSimulacion{}
\def\videoEntrenamiento{}

% ----------------------------------------------------------------------------
%  Bloque de evidencia por actividad
%  \actividad{nombre}{fase}{periodo}{tipo de evidencia}{ubicación en reporte TT}
%            {descripción}{porcentaje}{enlace de video}
% ----------------------------------------------------------------------------
\definecolor{encabezado}{gray}{0.85}
\newcommand\video[1]{\edef\tmp{#1}\ifx\tmp\empty No aplica.\else\url{#1}\fi}
\newcommand\actividad[8]{%
	\par\noindent
	\begin{minipage}{\linewidth}
		\renewcommand{\arraystretch}{1.25}
		\begin{tabularx}{\linewidth}{|>{\raggedright\bfseries\arraybackslash}p{0.24\linewidth}|X|}
			\hline
			\rowcolor{encabezado}
			Nombre de la actividad & \textbf{#1} \\ \hline
			Fase y periodo & #2 \,(#3) \\ \hline
			Evidencia de trabajo & #4 \\ \hline
			Ubicación en el reporte de TT & #5 \\ \hline
			Descripción del avance & #6 \\ \hline
			Porcentaje de cumplimiento & \textbf{#7} \\ \hline
			Enlace de video & \video{#8} \\ \hline
		\end{tabularx}
	\end{minipage}
	\par\vspace{1.2em}%
}

% Porcentaje de cumplimiento: real respecto a lo esperado al corte
\newcommand\cumplimiento[2]{\fpeval{min(100, round(100*(#2)/(#1), 0))}\%}

% Columna del mes reportado en el cronograma
\newcommand\mesActual[2]{\ifnum\corte=#1\cellcolor{black!70}\textcolor{white}{\textbf{#2}}\else\textbf{#2}\fi}

\begin{document}

	\begin{center}
		\fontsize{14}{16}\selectfont\bfseries
		\numReporte° Reporte de seguimiento (\mes{} 2026): \\
		Simulación de sistemas complejos mediante autómatas celulares neuronales
	\end{center}

	\begin{center}
		\fontsize{14}{16}\selectfont\bfseries
		\textit{Trabajo Terminal No. 2026-A114}
	\end{center}

	\begin{center}
		\fontsize{12}{16}\selectfont
		\textit{Alumno: \authormark{*}Angeles López Erick Jesse\authormark{1}} \\
		\textit{Director: Juárez Martínez Genaro\authormark{1}} \\
		\textit{\authormark{*}e-mail: eangelesl1700@alumno.ipn.mx}
	\end{center}

	\section{Información general}

	\begin{itemize}[leftmargin=*]
		\item \textbf{Periodo reportado:} \mes{} 2026 (corte al \fechaCorte).
		\item \textbf{Documento de evidencia adjunto:} reporte técnico de Trabajo Terminal
		\textit{``Simulación de sistemas complejos mediante autómatas celulares neuronales''} (PDF, 204 páginas).
		El avance es incremental: cada evidencia listada en este reporte se encuentra integrada en el capítulo
		y sección correspondiente de dicho documento, cuyas páginas se indican en cada actividad.
		\item \textbf{Actividades programadas en el mes:} \segun{Módulos de entrenamiento (cierre), Entrenamiento y ajuste,
		Experimentación y análisis, Análisis de resultados (global).}{Entrenamiento y ajuste, Experimentación y análisis,
		Análisis de resultados (global).}
	\end{itemize}

	\section{Resumen de cumplimiento}

	\begin{table}[H]
		\centering
		\renewcommand{\arraystretch}{1.2}
		\begin{tabular}{l c c c c}
			\toprule
			\textbf{Actividad} & \textbf{Periodo} & \textbf{Esperado al corte} & \textbf{Avance real} & \textbf{Cumplimiento} \\
			\midrule
			\multicolumn{5}{l}{\textit{Fase 1: Planificación y Revisión Bibliográfica}} \\
			Estado del arte                  & Ago/25 -- Oct/25 & 100\% & 100\% & 100\% \\
			Marco teórico y metodológico     & Sep/25 -- Oct/25 & 100\% & 100\% & 100\% \\
			\multicolumn{5}{l}{\textit{Fase 2: Diseño e implementación del simulador}} \\
			Arquitectura del simulador       & Oct/25 -- Nov/25 & 100\% & 100\% & 100\% \\
			Módulos de simulación            & Nov/25 -- Dic/25 & 100\% & 100\% & 100\% \\
			Módulos entrenamiento            & Dic/25 -- Ago/26 & 100\% & 100\% & 100\% \\
			\multicolumn{5}{l}{\textit{Fase 3: Entrenamiento y Validación}} \\
			Entrenamiento y ajuste           & Ago/26 -- Oct/26 & \espEntrenamiento\% & \realEntrenamiento\% & \cumplimiento{\espEntrenamiento}{\realEntrenamiento} \\
			Experimentación y análisis       & Ago/26 -- Nov/26 & \espExperimentos\%  & \realExperimentos\%  & \cumplimiento{\espExperimentos}{\realExperimentos} \\
			Análisis de resultados (global)  & Ago/26 -- Nov/26 & \espAnalisis\%      & \realAnalisis\%      & \cumplimiento{\espAnalisis}{\realAnalisis} \\
			\multicolumn{5}{l}{\textit{Fase 4: Documentación y Difusión}} \\
			Documentación técnica            & Oct/26 -- Dic/26 & 0\% & -- & No iniciada \\
			Artículo académico               & Oct/26 -- Dic/26 & 0\% & -- & No iniciada \\
			\bottomrule
		\end{tabular}
		\caption{Avance acumulado por actividad al \fechaCorte. ``Esperado al corte'' es la proporción del periodo
		programado en el cronograma que ha transcurrido; ``Cumplimiento'' compara el avance real con dicho valor.}
	\end{table}

	\section{Evidencias por actividad}

	\subsection*{Fase 1: Planificación y Revisión Bibliográfica}

	\actividad{Estado del arte}
	{Fase 1}{Ago/25 -- Oct/25}
	{Documento escrito.}
	{Capítulo 5, sección 5.1 \textit{Estado del arte, aplicaciones y líneas actuales de investigación} (pp.~61--77);
	Capítulo 1 \textit{Introducción} (pp.~1--6).}
	{Revisión de trabajos sobre autómatas celulares neuronales: morfogénesis, NCA isotrópicos, autoclasificación,
	representaciones tridimensionales, NCA en malla, modelos generativos y aplicaciones adicionales.
	Actividad concluida en Trabajo Terminal I.}
	{100\%}{}

	\actividad{Marco teórico y metodológico}
	{Fase 1}{Sep/25 -- Oct/25}
	{Documento escrito.}
	{Capítulos 2 a 5 (marco teórico, pp.~7--77) y Capítulo 6 \textit{Marco metodológico} (pp.~78--90).}
	{Marco teórico de sistemas complejos y teoría del caos, autómatas celulares, inteligencia artificial y redes
	neuronales, y autómatas celulares neuronales. Marco metodológico con diseño cuasiexperimental, población,
	muestra, técnicas de recolección y procesamiento de datos, y contextualización tecnológica.
	Actividad concluida en Trabajo Terminal I.}
	{100\%}{}

	\subsection*{Fase 2: Diseño e implementación del simulador}

	\actividad{Arquitectura del simulador}
	{Fase 2}{Oct/25 -- Nov/25}
	{Documento escrito (análisis y diseño de software, diagramas UML).}
	{Capítulo 7 \textit{Análisis de software} (pp.~91--135); Capítulo 8, secciones 8.1 \textit{Arquitectura},
	8.2 \textit{Diagrama de clases}, 8.3 \textit{Diseño de datos} y 8.5 \textit{Interfaces de usuarios} (pp.~136--143, 153--156).}
	{Metodología de desarrollo, reglas de negocio, requerimientos, análisis de riesgos, estudio de factibilidad,
	casos de uso con sus especificaciones, diagramas de secuencia, arquitectura MVC, diagrama de clases,
	diseño de datos e interfaces gráficas.}
	{100\%}{}

	\actividad{Módulos de simulación}
	{Fase 2}{Nov/25 -- Dic/25}
	{Documento escrito, tablas de resultados y pruebas de operatividad.}
	{Sección 8.4 \textit{Diseño de algoritmos} (8.4.1--8.4.4 y 8.4.6, pp.~144--152); secciones 9.1 \textit{Entorno de desarrollo}
	y 9.2 \textit{Estructura del proyecto} (pp.~157--162); subsecciones 9.4.1, 9.4.3, 9.4.4 y 9.4.5 de
	\textit{Validación y pruebas} (pp.~166--169), tablas 9.6 y 9.7.}
	{Núcleo de simulación en C++/CUDA con enlaces a Python e interfaz en Qt: indexación de la regla semitotalística,
	paso de evolución, condiciones de frontera, \textit{forward pass} del NCA y \textit{prefetch} asíncrono.
	Validación funcional contra reglas conocidas y pruebas de rendimiento (hasta 345\,000 pasos/s en
	$100\times100$ y más de 100 fps de renderizado en $2000\times2000$).}
	{100\%}{\videoSimulacion}

	\actividad{Módulos entrenamiento}
	{Fase 2}{Dic/25 -- Ago/26}
	{Documento escrito, pruebas de operatividad.}
	{Subsección 7.8.4 \textit{Especificaciones de casos de uso: Entrenamiento} (pp.~125--130); 8.3.5
	\textit{Configuración del NCA} (p.~143); 8.4.5 \textit{Entrenamiento del autómata celular neuronal} (pp.~150--151);
	sección 9.3 \textit{Entrenamiento del modelo} (pp.~163--166, tabla 9.5); subsecciones 9.4.2 y 9.4.6 (pp.~167, 169).}
	{Módulo de entrenamiento en PyTorch con construcción dinámica del perceptrón multicapa, ciclo de optimización
	con Adam y detención temprana, persistencia de pesos en \texttt{.npz} compatible con el \textit{kernel} CUDA e
	integración con la pestaña de entrenamiento de la GUI (subproceso con \textit{log} en tiempo real).
	Se agregó el parámetro de semilla aleatoria para reproducir experimentos.
	\segun{Actividad concluida en el mes reportado.}{Actividad concluida en agosto de 2026.}}
	{100\%}{\videoEntrenamiento}

	\subsection*{Fase 3: Entrenamiento y Validación}

	\actividad{Entrenamiento y ajuste}
	{Fase 3}{Ago/26 -- Oct/26}
	{Documento escrito y tablas de resultados.}
	{Capítulo 10, sección 10.1 \textit{Duplicación de reglas del juego de la vida}:
	\segun{tablas 10.1 y 10.2 y figura 10.3 (pp.~173--174).}
	{tablas 10.1 y 10.2 y figura 10.3 (pp.~173--174); figuras 10.5 a 10.9 (pp.~176--178).}}
	{\segun{Barrido de hiperparámetros para la regla B3/S23: 10 combinaciones de funciones de activación con
	una capa oculta de 4 neuronas (una sola converge) y con 4 capas ocultas de 8 neuronas (cinco convergen a una
	pérdida menor a $10^{-4}$). Se identifica el estancamiento en las primeras 12 épocas.}
	{Además del barrido de hiperparámetros de agosto, se entrenaron las reglas B3/S23, B3/S238, B38/S23, B3/S236
	y B38/S238 con distintas pérdidas objetivo y semillas, midiendo la distancia Manhattan entre el vector de regla
	aprendido y el original. Se determinó que la pérdida objetivo debe ser $\leq 0.1$ para B3/S23.}}
	{\realEntrenamiento\% (esperado: \espEntrenamiento\%)}{}

	\actividad{Experimentación y análisis}
	{Fase 3}{Ago/26 -- Nov/26}
	{Documento escrito, tablas y gráficas de resultados.}
	{Capítulo 10, sección 10.1:
	\segun{figuras 10.1 y 10.2 (pp.~170--172).}
	{figuras 10.1 a 10.13 (pp.~170--181).}}
	{\segun{Generación del conjunto de datos (100 condiciones iniciales en $100\times100$ durante 100 pasos) y
	análisis de su balance: el estado 1 aparece un $16.5\% \pm 4.2\%$ y los patrones $C0N8$ y $C1N8$ representan
	menos del $0.02\%$ de las entradas.}
	{Primer experimento (duplicación de reglas del juego de la vida) completo. Se diseñó un espacio con las
	512 permutaciones de la vecindad para extraer el vector de regla aprendido y verificar si el NCA conserva la
	propiedad semitotalística. Se identificaron reglas aprendidas alternativas (B3/S23 en lugar de B3/S238,
	B3/S23678 en lugar de B3/S236) y reglas no semitotalísticas.}}
	{\realExperimentos\% (esperado: \espExperimentos\%)}{}

	\actividad{Análisis de resultados (global)}
	{Fase 3}{Ago/26 -- Nov/26}
	{Documento escrito.}
	{Capítulo 10, sección 10.1\segun{ (pp.~170--172).}{, conclusiones del experimento (p.~181).}}
	{\segun{Análisis estadístico de la distribución de estados y de patrones de vecindad para la regla B3/S23,
	identificando el sesgo de muestra y proponiendo sobremuestreo o ponderación de la función de costo.}
	{Conclusiones del primer experimento: la generalización depende del sesgo de muestra de patrones poco frecuentes,
	el entrenamiento de reglas semitotalísticas no garantiza una regla aprendida semitotalística, y no se requieren
	arquitecturas complejas para replicar estas dinámicas. Se valida la metodología de extracción del vector de
	regla para los siguientes experimentos.}}
	{\realAnalisis\% (esperado: \espAnalisis\%)}{}

	\subsection*{Fase 4: Documentación y Difusión}

	\actividad{Documentación técnica}
	{Fase 4}{Oct/26 -- Dic/26}
	{No aplica en el periodo.}
	{Reporte técnico de TT completo (204 páginas).}
	{Actividad programada a partir de octubre de 2026. El reporte técnico se ha redactado de forma incremental
	conforme avanza cada actividad, por lo que su integración final partirá del documento adjunto.}
	{No iniciada (0\% esperado al corte)}{}

	\actividad{Artículo académico}
	{Fase 4}{Oct/26 -- Dic/26}
	{No aplica en el periodo.}
	{--}
	{Actividad programada a partir de octubre de 2026.}
	{No iniciada (0\% esperado al corte)}{}

	\section{Cronograma}

	\resizebox{\linewidth}{!}{%
		\begin{tabular}{p{5cm}*{10}{>{\centering\arraybackslash}p{1cm}}}
			\toprule
			\textbf{Actividad} & \textbf{Ago/25} & \textbf{Sep/25} & \textbf{Oct/25} & \textbf{Nov/25} & \textbf{Dic/25} & \mesActual{1}{Ago/26} & \mesActual{2}{Sep/26} & \textbf{Oct/26} & \textbf{Nov/26} & \textbf{Dic/26} \\
			\midrule
			\multicolumn{11}{l}{\textbf{Fase 1: Planificación y Revisión Bibliográfica}} \\

			Estado del arte       & \cellcolor{gray!30} & \cellcolor{gray!30} & \cellcolor{gray!30} &  &  &  &  &  &  &   \\

			Marco teórico y metodológico &  & \cellcolor{gray!30} & \cellcolor{gray!30} &  &  &  &  &  &  &   \\

			\midrule

			\multicolumn{11}{l}{\textbf{Fase 2: Diseño e implementación del simulador}} \\
			Arquitectura del simulador &  &  & \cellcolor{gray!30} & \cellcolor{gray!30} &  &  &  &  &  &   \\

			Módulos de simulación  &  &  &  & \cellcolor{gray!30} & \cellcolor{gray!30} &  &  &  &  &   \\

			Módulos entrenamiento   &  &  &  &  & \cellcolor{gray!30} & \cellcolor{gray!30} &  &  &  &   \\

			\midrule

			\multicolumn{11}{l}{\textbf{Fase 3: Entrenamiento y Validación}} \\
			Entrenamiento y ajuste   &  &  &  &  &  & \cellcolor{gray!30} & \cellcolor{gray!30} & \cellcolor{gray!30} &   &   \\

			Experimentación y análisis  &  &  &  &  &  & \cellcolor{gray!30}  & \cellcolor{gray!30} & \cellcolor{gray!30} & \cellcolor{gray!30} &   \\

			Análisis de resultados (global) &  &  &  &  &  & \cellcolor{gray!30} & \cellcolor{gray!30} & \cellcolor{gray!30} & \cellcolor{gray!30} &   \\

			\midrule

			\multicolumn{11}{l}{\textbf{Fase 4: Documentación y Difusión}} \\
			Documentación técnica  &  &  &  &  &  &  &  & \cellcolor{gray!30}  & \cellcolor{gray!30} & \cellcolor{gray!30}  \\

			Artículo académico  &  &  &  &  &  &  &  & \cellcolor{gray!30}  & \cellcolor{gray!30}  & \cellcolor{gray!30}  \\

			\bottomrule
		\end{tabular}
	}

	\vspace{0.5em}
	{\small La columna resaltada corresponde al mes reportado.}

	\section{Alumnos y directores}

	\vspace{1em}
	\begin{tabular}{@{} p{0.45\textwidth} p{0.45\textwidth} @{}}
		\textit{Erick Jesse Angeles López} & \textit{Dr. Genaro Juárez Martínez} \\[5em]
		Firma: \rule{6cm}{0.4pt} & Firma: \rule{6cm}{0.4pt} \\
	\end{tabular}

\end{document}
"
%  o simplemente:  ./compilar.sh
%
%  Las páginas y secciones citadas corresponden a Documentation/main.pdf
%  (reporte técnico de TT), que es el PDF que se adjunta como evidencia.
% ============================================================================
\documentclass[letterpaper,10pt,oneside]{article}
\usepackage[margin=2.54cm]{geometry}
\usepackage[spanish]{babel}
\newcommand\authormark[1]{\textsuperscript{#1}}
\usepackage[colorlinks=true,bookmarks=true,linkcolor=black,citecolor=black,urlcolor=blue]{hyperref}
\usepackage{times}
\usepackage[table]{xcolor}
\usepackage{booktabs}
\usepackage{graphicx}
\usepackage{titlesec}
\usepackage{tabularx}
\usepackage{array}
\usepackage{enumitem}
\usepackage{float}

\titleformat{\section}
{\normalfont\fontsize{12pt}{14pt}\bfseries}
{\thesection}{1em}{}

% ----------------------------------------------------------------------------
%  Mes reportado
% ----------------------------------------------------------------------------
\providecommand\corte{2}
\ifnum\corte=1
	\def\numReporte{1}
	\def\mes{Agosto}
	\def\fechaCorte{31 de agosto de 2026}
\else
	\def\numReporte{2}
	\def\mes{Septiembre}
	\def\fechaCorte{30 de septiembre de 2026}
\fi
% \segun{texto agosto}{texto septiembre}
\newcommand\segun[2]{\ifnum\corte=1 #1\else #2\fi}

% ----------------------------------------------------------------------------
%  Porcentajes (editar aquí). Formato: \segun{agosto}{septiembre}
%  "esp" = avance esperado según cronograma al corte, "real" = avance logrado.
% ----------------------------------------------------------------------------
\def\espEntrenamiento{\segun{33}{67}}   \def\realEntrenamiento{\segun{33}{67}}
\def\espExperimentos{\segun{25}{50}}    \def\realExperimentos{\segun{25}{50}}
\def\espAnalisis{\segun{25}{50}}        \def\realAnalisis{\segun{25}{50}}

% ----------------------------------------------------------------------------
%  Enlaces a video (dejar vacío si no aplica)
% ----------------------------------------------------------------------------
\def\videoSimulacion{}
\def\videoEntrenamiento{}

% ----------------------------------------------------------------------------
%  Bloque de evidencia por actividad
%  \actividad{nombre}{fase}{periodo}{tipo de evidencia}{ubicación en reporte TT}
%            {descripción}{porcentaje}{enlace de video}
% ----------------------------------------------------------------------------
\definecolor{encabezado}{gray}{0.85}
\newcommand\video[1]{\edef\tmp{#1}\ifx\tmp\empty No aplica.\else\url{#1}\fi}
\newcommand\actividad[8]{%
	\par\noindent
	\begin{minipage}{\linewidth}
		\renewcommand{\arraystretch}{1.25}
		\begin{tabularx}{\linewidth}{|>{\raggedright\bfseries\arraybackslash}p{0.24\linewidth}|X|}
			\hline
			\rowcolor{encabezado}
			Nombre de la actividad & \textbf{#1} \\ \hline
			Fase y periodo & #2 \,(#3) \\ \hline
			Evidencia de trabajo & #4 \\ \hline
			Ubicación en el reporte de TT & #5 \\ \hline
			Descripción del avance & #6 \\ \hline
			Porcentaje de cumplimiento & \textbf{#7} \\ \hline
			Enlace de video & \video{#8} \\ \hline
		\end{tabularx}
	\end{minipage}
	\par\vspace{1.2em}%
}

% Porcentaje de cumplimiento: real respecto a lo esperado al corte
\newcommand\cumplimiento[2]{\fpeval{min(100, round(100*(#2)/(#1), 0))}\%}

% Columna del mes reportado en el cronograma
\newcommand\mesActual[2]{\ifnum\corte=#1\cellcolor{black!70}\textcolor{white}{\textbf{#2}}\else\textbf{#2}\fi}

\begin{document}

	\begin{center}
		\fontsize{14}{16}\selectfont\bfseries
		\numReporte° Reporte de seguimiento (\mes{} 2026): \\
		Simulación de sistemas complejos mediante autómatas celulares neuronales
	\end{center}

	\begin{center}
		\fontsize{14}{16}\selectfont\bfseries
		\textit{Trabajo Terminal No. 2026-A114}
	\end{center}

	\begin{center}
		\fontsize{12}{16}\selectfont
		\textit{Alumno: \authormark{*}Angeles López Erick Jesse\authormark{1}} \\
		\textit{Director: Juárez Martínez Genaro\authormark{1}} \\
		\textit{\authormark{*}e-mail: eangelesl1700@alumno.ipn.mx}
	\end{center}

	\section{Información general}

	\begin{itemize}[leftmargin=*]
		\item \textbf{Periodo reportado:} \mes{} 2026 (corte al \fechaCorte).
		\item \textbf{Documento de evidencia adjunto:} reporte técnico de Trabajo Terminal
		\textit{``Simulación de sistemas complejos mediante autómatas celulares neuronales''} (PDF, 204 páginas).
		El avance es incremental: cada evidencia listada en este reporte se encuentra integrada en el capítulo
		y sección correspondiente de dicho documento, cuyas páginas se indican en cada actividad.
		\item \textbf{Actividades programadas en el mes:} \segun{Módulos de entrenamiento (cierre), Entrenamiento y ajuste,
		Experimentación y análisis, Análisis de resultados (global).}{Entrenamiento y ajuste, Experimentación y análisis,
		Análisis de resultados (global).}
	\end{itemize}

	\section{Resumen de cumplimiento}

	\begin{table}[H]
		\centering
		\renewcommand{\arraystretch}{1.2}
		\begin{tabular}{l c c c c}
			\toprule
			\textbf{Actividad} & \textbf{Periodo} & \textbf{Esperado al corte} & \textbf{Avance real} & \textbf{Cumplimiento} \\
			\midrule
			\multicolumn{5}{l}{\textit{Fase 1: Planificación y Revisión Bibliográfica}} \\
			Estado del arte                  & Ago/25 -- Oct/25 & 100\% & 100\% & 100\% \\
			Marco teórico y metodológico     & Sep/25 -- Oct/25 & 100\% & 100\% & 100\% \\
			\multicolumn{5}{l}{\textit{Fase 2: Diseño e implementación del simulador}} \\
			Arquitectura del simulador       & Oct/25 -- Nov/25 & 100\% & 100\% & 100\% \\
			Módulos de simulación            & Nov/25 -- Dic/25 & 100\% & 100\% & 100\% \\
			Módulos entrenamiento            & Dic/25 -- Ago/26 & 100\% & 100\% & 100\% \\
			\multicolumn{5}{l}{\textit{Fase 3: Entrenamiento y Validación}} \\
			Entrenamiento y ajuste           & Ago/26 -- Oct/26 & \espEntrenamiento\% & \realEntrenamiento\% & \cumplimiento{\espEntrenamiento}{\realEntrenamiento} \\
			Experimentación y análisis       & Ago/26 -- Nov/26 & \espExperimentos\%  & \realExperimentos\%  & \cumplimiento{\espExperimentos}{\realExperimentos} \\
			Análisis de resultados (global)  & Ago/26 -- Nov/26 & \espAnalisis\%      & \realAnalisis\%      & \cumplimiento{\espAnalisis}{\realAnalisis} \\
			\multicolumn{5}{l}{\textit{Fase 4: Documentación y Difusión}} \\
			Documentación técnica            & Oct/26 -- Dic/26 & 0\% & -- & No iniciada \\
			Artículo académico               & Oct/26 -- Dic/26 & 0\% & -- & No iniciada \\
			\bottomrule
		\end{tabular}
		\caption{Avance acumulado por actividad al \fechaCorte. ``Esperado al corte'' es la proporción del periodo
		programado en el cronograma que ha transcurrido; ``Cumplimiento'' compara el avance real con dicho valor.}
	\end{table}

	\section{Evidencias por actividad}

	\subsection*{Fase 1: Planificación y Revisión Bibliográfica}

	\actividad{Estado del arte}
	{Fase 1}{Ago/25 -- Oct/25}
	{Documento escrito.}
	{Capítulo 5, sección 5.1 \textit{Estado del arte, aplicaciones y líneas actuales de investigación} (pp.~61--77);
	Capítulo 1 \textit{Introducción} (pp.~1--6).}
	{Revisión de trabajos sobre autómatas celulares neuronales: morfogénesis, NCA isotrópicos, autoclasificación,
	representaciones tridimensionales, NCA en malla, modelos generativos y aplicaciones adicionales.
	Actividad concluida en Trabajo Terminal I.}
	{100\%}{}

	\actividad{Marco teórico y metodológico}
	{Fase 1}{Sep/25 -- Oct/25}
	{Documento escrito.}
	{Capítulos 2 a 5 (marco teórico, pp.~7--77) y Capítulo 6 \textit{Marco metodológico} (pp.~78--90).}
	{Marco teórico de sistemas complejos y teoría del caos, autómatas celulares, inteligencia artificial y redes
	neuronales, y autómatas celulares neuronales. Marco metodológico con diseño cuasiexperimental, población,
	muestra, técnicas de recolección y procesamiento de datos, y contextualización tecnológica.
	Actividad concluida en Trabajo Terminal I.}
	{100\%}{}

	\subsection*{Fase 2: Diseño e implementación del simulador}

	\actividad{Arquitectura del simulador}
	{Fase 2}{Oct/25 -- Nov/25}
	{Documento escrito (análisis y diseño de software, diagramas UML).}
	{Capítulo 7 \textit{Análisis de software} (pp.~91--135); Capítulo 8, secciones 8.1 \textit{Arquitectura},
	8.2 \textit{Diagrama de clases}, 8.3 \textit{Diseño de datos} y 8.5 \textit{Interfaces de usuarios} (pp.~136--143, 153--156).}
	{Metodología de desarrollo, reglas de negocio, requerimientos, análisis de riesgos, estudio de factibilidad,
	casos de uso con sus especificaciones, diagramas de secuencia, arquitectura MVC, diagrama de clases,
	diseño de datos e interfaces gráficas.}
	{100\%}{}

	\actividad{Módulos de simulación}
	{Fase 2}{Nov/25 -- Dic/25}
	{Documento escrito, tablas de resultados y pruebas de operatividad.}
	{Sección 8.4 \textit{Diseño de algoritmos} (8.4.1--8.4.4 y 8.4.6, pp.~144--152); secciones 9.1 \textit{Entorno de desarrollo}
	y 9.2 \textit{Estructura del proyecto} (pp.~157--162); subsecciones 9.4.1, 9.4.3, 9.4.4 y 9.4.5 de
	\textit{Validación y pruebas} (pp.~166--169), tablas 9.6 y 9.7.}
	{Núcleo de simulación en C++/CUDA con enlaces a Python e interfaz en Qt: indexación de la regla semitotalística,
	paso de evolución, condiciones de frontera, \textit{forward pass} del NCA y \textit{prefetch} asíncrono.
	Validación funcional contra reglas conocidas y pruebas de rendimiento (hasta 345\,000 pasos/s en
	$100\times100$ y más de 100 fps de renderizado en $2000\times2000$).}
	{100\%}{\videoSimulacion}

	\actividad{Módulos entrenamiento}
	{Fase 2}{Dic/25 -- Ago/26}
	{Documento escrito, pruebas de operatividad.}
	{Subsección 7.8.4 \textit{Especificaciones de casos de uso: Entrenamiento} (pp.~125--130); 8.3.5
	\textit{Configuración del NCA} (p.~143); 8.4.5 \textit{Entrenamiento del autómata celular neuronal} (pp.~150--151);
	sección 9.3 \textit{Entrenamiento del modelo} (pp.~163--166, tabla 9.5); subsecciones 9.4.2 y 9.4.6 (pp.~167, 169).}
	{Módulo de entrenamiento en PyTorch con construcción dinámica del perceptrón multicapa, ciclo de optimización
	con Adam y detención temprana, persistencia de pesos en \texttt{.npz} compatible con el \textit{kernel} CUDA e
	integración con la pestaña de entrenamiento de la GUI (subproceso con \textit{log} en tiempo real).
	Se agregó el parámetro de semilla aleatoria para reproducir experimentos.
	\segun{Actividad concluida en el mes reportado.}{Actividad concluida en agosto de 2026.}}
	{100\%}{\videoEntrenamiento}

	\subsection*{Fase 3: Entrenamiento y Validación}

	\actividad{Entrenamiento y ajuste}
	{Fase 3}{Ago/26 -- Oct/26}
	{Documento escrito y tablas de resultados.}
	{Capítulo 10, sección 10.1 \textit{Duplicación de reglas del juego de la vida}:
	\segun{tablas 10.1 y 10.2 y figura 10.3 (pp.~173--174).}
	{tablas 10.1 y 10.2 y figura 10.3 (pp.~173--174); figuras 10.5 a 10.9 (pp.~176--178).}}
	{\segun{Barrido de hiperparámetros para la regla B3/S23: 10 combinaciones de funciones de activación con
	una capa oculta de 4 neuronas (una sola converge) y con 4 capas ocultas de 8 neuronas (cinco convergen a una
	pérdida menor a $10^{-4}$). Se identifica el estancamiento en las primeras 12 épocas.}
	{Además del barrido de hiperparámetros de agosto, se entrenaron las reglas B3/S23, B3/S238, B38/S23, B3/S236
	y B38/S238 con distintas pérdidas objetivo y semillas, midiendo la distancia Manhattan entre el vector de regla
	aprendido y el original. Se determinó que la pérdida objetivo debe ser $\leq 0.1$ para B3/S23.}}
	{\realEntrenamiento\% (esperado: \espEntrenamiento\%)}{}

	\actividad{Experimentación y análisis}
	{Fase 3}{Ago/26 -- Nov/26}
	{Documento escrito, tablas y gráficas de resultados.}
	{Capítulo 10, sección 10.1:
	\segun{figuras 10.1 y 10.2 (pp.~170--172).}
	{figuras 10.1 a 10.13 (pp.~170--181).}}
	{\segun{Generación del conjunto de datos (100 condiciones iniciales en $100\times100$ durante 100 pasos) y
	análisis de su balance: el estado 1 aparece un $16.5\% \pm 4.2\%$ y los patrones $C0N8$ y $C1N8$ representan
	menos del $0.02\%$ de las entradas.}
	{Primer experimento (duplicación de reglas del juego de la vida) completo. Se diseñó un espacio con las
	512 permutaciones de la vecindad para extraer el vector de regla aprendido y verificar si el NCA conserva la
	propiedad semitotalística. Se identificaron reglas aprendidas alternativas (B3/S23 en lugar de B3/S238,
	B3/S23678 en lugar de B3/S236) y reglas no semitotalísticas.}}
	{\realExperimentos\% (esperado: \espExperimentos\%)}{}

	\actividad{Análisis de resultados (global)}
	{Fase 3}{Ago/26 -- Nov/26}
	{Documento escrito.}
	{Capítulo 10, sección 10.1\segun{ (pp.~170--172).}{, conclusiones del experimento (p.~181).}}
	{\segun{Análisis estadístico de la distribución de estados y de patrones de vecindad para la regla B3/S23,
	identificando el sesgo de muestra y proponiendo sobremuestreo o ponderación de la función de costo.}
	{Conclusiones del primer experimento: la generalización depende del sesgo de muestra de patrones poco frecuentes,
	el entrenamiento de reglas semitotalísticas no garantiza una regla aprendida semitotalística, y no se requieren
	arquitecturas complejas para replicar estas dinámicas. Se valida la metodología de extracción del vector de
	regla para los siguientes experimentos.}}
	{\realAnalisis\% (esperado: \espAnalisis\%)}{}

	\subsection*{Fase 4: Documentación y Difusión}

	\actividad{Documentación técnica}
	{Fase 4}{Oct/26 -- Dic/26}
	{No aplica en el periodo.}
	{Reporte técnico de TT completo (204 páginas).}
	{Actividad programada a partir de octubre de 2026. El reporte técnico se ha redactado de forma incremental
	conforme avanza cada actividad, por lo que su integración final partirá del documento adjunto.}
	{No iniciada (0\% esperado al corte)}{}

	\actividad{Artículo académico}
	{Fase 4}{Oct/26 -- Dic/26}
	{No aplica en el periodo.}
	{--}
	{Actividad programada a partir de octubre de 2026.}
	{No iniciada (0\% esperado al corte)}{}

	\section{Cronograma}

	\resizebox{\linewidth}{!}{%
		\begin{tabular}{p{5cm}*{10}{>{\centering\arraybackslash}p{1cm}}}
			\toprule
			\textbf{Actividad} & \textbf{Ago/25} & \textbf{Sep/25} & \textbf{Oct/25} & \textbf{Nov/25} & \textbf{Dic/25} & \mesActual{1}{Ago/26} & \mesActual{2}{Sep/26} & \textbf{Oct/26} & \textbf{Nov/26} & \textbf{Dic/26} \\
			\midrule
			\multicolumn{11}{l}{\textbf{Fase 1: Planificación y Revisión Bibliográfica}} \\

			Estado del arte       & \cellcolor{gray!30} & \cellcolor{gray!30} & \cellcolor{gray!30} &  &  &  &  &  &  &   \\

			Marco teórico y metodológico &  & \cellcolor{gray!30} & \cellcolor{gray!30} &  &  &  &  &  &  &   \\

			\midrule

			\multicolumn{11}{l}{\textbf{Fase 2: Diseño e implementación del simulador}} \\
			Arquitectura del simulador &  &  & \cellcolor{gray!30} & \cellcolor{gray!30} &  &  &  &  &  &   \\

			Módulos de simulación  &  &  &  & \cellcolor{gray!30} & \cellcolor{gray!30} &  &  &  &  &   \\

			Módulos entrenamiento   &  &  &  &  & \cellcolor{gray!30} & \cellcolor{gray!30} &  &  &  &   \\

			\midrule

			\multicolumn{11}{l}{\textbf{Fase 3: Entrenamiento y Validación}} \\
			Entrenamiento y ajuste   &  &  &  &  &  & \cellcolor{gray!30} & \cellcolor{gray!30} & \cellcolor{gray!30} &   &   \\

			Experimentación y análisis  &  &  &  &  &  & \cellcolor{gray!30}  & \cellcolor{gray!30} & \cellcolor{gray!30} & \cellcolor{gray!30} &   \\

			Análisis de resultados (global) &  &  &  &  &  & \cellcolor{gray!30} & \cellcolor{gray!30} & \cellcolor{gray!30} & \cellcolor{gray!30} &   \\

			\midrule

			\multicolumn{11}{l}{\textbf{Fase 4: Documentación y Difusión}} \\
			Documentación técnica  &  &  &  &  &  &  &  & \cellcolor{gray!30}  & \cellcolor{gray!30} & \cellcolor{gray!30}  \\

			Artículo académico  &  &  &  &  &  &  &  & \cellcolor{gray!30}  & \cellcolor{gray!30}  & \cellcolor{gray!30}  \\

			\bottomrule
		\end{tabular}
	}

	\vspace{0.5em}
	{\small La columna resaltada corresponde al mes reportado.}

	\section{Alumnos y directores}

	\vspace{1em}
	\begin{tabular}{@{} p{0.45\textwidth} p{0.45\textwidth} @{}}
		\textit{Erick Jesse Angeles López} & \textit{Dr. Genaro Juárez Martínez} \\[5em]
		Firma: \rule{6cm}{0.4pt} & Firma: \rule{6cm}{0.4pt} \\
	\end{tabular}

\end{document}
"
%      pdflatex -jobname=2ReporteSeguimiento "\def\corte{2}% ============================================================================
%  Reporte de seguimiento - Trabajo Terminal II
%
%  Un solo archivo genera los reportes de cada mes. El mes se elige con \corte:
%      1 = Agosto 2026, 2 = Septiembre 2026
%
%  Compilación (desde esta carpeta):
%      pdflatex -jobname=1ReporteSeguimiento "\def\corte{1}% ============================================================================
%  Reporte de seguimiento - Trabajo Terminal II
%
%  Un solo archivo genera los reportes de cada mes. El mes se elige con \corte:
%      1 = Agosto 2026, 2 = Septiembre 2026
%
%  Compilación (desde esta carpeta):
%      pdflatex -jobname=1ReporteSeguimiento "\def\corte{1}\input{reporteSeguimiento}"
%      pdflatex -jobname=2ReporteSeguimiento "\def\corte{2}\input{reporteSeguimiento}"
%  o simplemente:  ./compilar.sh
%
%  Las páginas y secciones citadas corresponden a Documentation/main.pdf
%  (reporte técnico de TT), que es el PDF que se adjunta como evidencia.
% ============================================================================
\documentclass[letterpaper,10pt,oneside]{article}
\usepackage[margin=2.54cm]{geometry}
\usepackage[spanish]{babel}
\newcommand\authormark[1]{\textsuperscript{#1}}
\usepackage[colorlinks=true,bookmarks=true,linkcolor=black,citecolor=black,urlcolor=blue]{hyperref}
\usepackage{times}
\usepackage[table]{xcolor}
\usepackage{booktabs}
\usepackage{graphicx}
\usepackage{titlesec}
\usepackage{tabularx}
\usepackage{array}
\usepackage{enumitem}
\usepackage{float}

\titleformat{\section}
{\normalfont\fontsize{12pt}{14pt}\bfseries}
{\thesection}{1em}{}

% ----------------------------------------------------------------------------
%  Mes reportado
% ----------------------------------------------------------------------------
\providecommand\corte{2}
\ifnum\corte=1
	\def\numReporte{1}
	\def\mes{Agosto}
	\def\fechaCorte{31 de agosto de 2026}
\else
	\def\numReporte{2}
	\def\mes{Septiembre}
	\def\fechaCorte{30 de septiembre de 2026}
\fi
% \segun{texto agosto}{texto septiembre}
\newcommand\segun[2]{\ifnum\corte=1 #1\else #2\fi}

% ----------------------------------------------------------------------------
%  Porcentajes (editar aquí). Formato: \segun{agosto}{septiembre}
%  "esp" = avance esperado según cronograma al corte, "real" = avance logrado.
% ----------------------------------------------------------------------------
\def\espEntrenamiento{\segun{33}{67}}   \def\realEntrenamiento{\segun{33}{67}}
\def\espExperimentos{\segun{25}{50}}    \def\realExperimentos{\segun{25}{50}}
\def\espAnalisis{\segun{25}{50}}        \def\realAnalisis{\segun{25}{50}}

% ----------------------------------------------------------------------------
%  Enlaces a video (dejar vacío si no aplica)
% ----------------------------------------------------------------------------
\def\videoSimulacion{}
\def\videoEntrenamiento{}

% ----------------------------------------------------------------------------
%  Bloque de evidencia por actividad
%  \actividad{nombre}{fase}{periodo}{tipo de evidencia}{ubicación en reporte TT}
%            {descripción}{porcentaje}{enlace de video}
% ----------------------------------------------------------------------------
\definecolor{encabezado}{gray}{0.85}
\newcommand\video[1]{\edef\tmp{#1}\ifx\tmp\empty No aplica.\else\url{#1}\fi}
\newcommand\actividad[8]{%
	\par\noindent
	\begin{minipage}{\linewidth}
		\renewcommand{\arraystretch}{1.25}
		\begin{tabularx}{\linewidth}{|>{\raggedright\bfseries\arraybackslash}p{0.24\linewidth}|X|}
			\hline
			\rowcolor{encabezado}
			Nombre de la actividad & \textbf{#1} \\ \hline
			Fase y periodo & #2 \,(#3) \\ \hline
			Evidencia de trabajo & #4 \\ \hline
			Ubicación en el reporte de TT & #5 \\ \hline
			Descripción del avance & #6 \\ \hline
			Porcentaje de cumplimiento & \textbf{#7} \\ \hline
			Enlace de video & \video{#8} \\ \hline
		\end{tabularx}
	\end{minipage}
	\par\vspace{1.2em}%
}

% Porcentaje de cumplimiento: real respecto a lo esperado al corte
\newcommand\cumplimiento[2]{\fpeval{min(100, round(100*(#2)/(#1), 0))}\%}

% Columna del mes reportado en el cronograma
\newcommand\mesActual[2]{\ifnum\corte=#1\cellcolor{black!70}\textcolor{white}{\textbf{#2}}\else\textbf{#2}\fi}

\begin{document}

	\begin{center}
		\fontsize{14}{16}\selectfont\bfseries
		\numReporte° Reporte de seguimiento (\mes{} 2026): \\
		Simulación de sistemas complejos mediante autómatas celulares neuronales
	\end{center}

	\begin{center}
		\fontsize{14}{16}\selectfont\bfseries
		\textit{Trabajo Terminal No. 2026-A114}
	\end{center}

	\begin{center}
		\fontsize{12}{16}\selectfont
		\textit{Alumno: \authormark{*}Angeles López Erick Jesse\authormark{1}} \\
		\textit{Director: Juárez Martínez Genaro\authormark{1}} \\
		\textit{\authormark{*}e-mail: eangelesl1700@alumno.ipn.mx}
	\end{center}

	\section{Información general}

	\begin{itemize}[leftmargin=*]
		\item \textbf{Periodo reportado:} \mes{} 2026 (corte al \fechaCorte).
		\item \textbf{Documento de evidencia adjunto:} reporte técnico de Trabajo Terminal
		\textit{``Simulación de sistemas complejos mediante autómatas celulares neuronales''} (PDF, 204 páginas).
		El avance es incremental: cada evidencia listada en este reporte se encuentra integrada en el capítulo
		y sección correspondiente de dicho documento, cuyas páginas se indican en cada actividad.
		\item \textbf{Actividades programadas en el mes:} \segun{Módulos de entrenamiento (cierre), Entrenamiento y ajuste,
		Experimentación y análisis, Análisis de resultados (global).}{Entrenamiento y ajuste, Experimentación y análisis,
		Análisis de resultados (global).}
	\end{itemize}

	\section{Resumen de cumplimiento}

	\begin{table}[H]
		\centering
		\renewcommand{\arraystretch}{1.2}
		\begin{tabular}{l c c c c}
			\toprule
			\textbf{Actividad} & \textbf{Periodo} & \textbf{Esperado al corte} & \textbf{Avance real} & \textbf{Cumplimiento} \\
			\midrule
			\multicolumn{5}{l}{\textit{Fase 1: Planificación y Revisión Bibliográfica}} \\
			Estado del arte                  & Ago/25 -- Oct/25 & 100\% & 100\% & 100\% \\
			Marco teórico y metodológico     & Sep/25 -- Oct/25 & 100\% & 100\% & 100\% \\
			\multicolumn{5}{l}{\textit{Fase 2: Diseño e implementación del simulador}} \\
			Arquitectura del simulador       & Oct/25 -- Nov/25 & 100\% & 100\% & 100\% \\
			Módulos de simulación            & Nov/25 -- Dic/25 & 100\% & 100\% & 100\% \\
			Módulos entrenamiento            & Dic/25 -- Ago/26 & 100\% & 100\% & 100\% \\
			\multicolumn{5}{l}{\textit{Fase 3: Entrenamiento y Validación}} \\
			Entrenamiento y ajuste           & Ago/26 -- Oct/26 & \espEntrenamiento\% & \realEntrenamiento\% & \cumplimiento{\espEntrenamiento}{\realEntrenamiento} \\
			Experimentación y análisis       & Ago/26 -- Nov/26 & \espExperimentos\%  & \realExperimentos\%  & \cumplimiento{\espExperimentos}{\realExperimentos} \\
			Análisis de resultados (global)  & Ago/26 -- Nov/26 & \espAnalisis\%      & \realAnalisis\%      & \cumplimiento{\espAnalisis}{\realAnalisis} \\
			\multicolumn{5}{l}{\textit{Fase 4: Documentación y Difusión}} \\
			Documentación técnica            & Oct/26 -- Dic/26 & 0\% & -- & No iniciada \\
			Artículo académico               & Oct/26 -- Dic/26 & 0\% & -- & No iniciada \\
			\bottomrule
		\end{tabular}
		\caption{Avance acumulado por actividad al \fechaCorte. ``Esperado al corte'' es la proporción del periodo
		programado en el cronograma que ha transcurrido; ``Cumplimiento'' compara el avance real con dicho valor.}
	\end{table}

	\section{Evidencias por actividad}

	\subsection*{Fase 1: Planificación y Revisión Bibliográfica}

	\actividad{Estado del arte}
	{Fase 1}{Ago/25 -- Oct/25}
	{Documento escrito.}
	{Capítulo 5, sección 5.1 \textit{Estado del arte, aplicaciones y líneas actuales de investigación} (pp.~61--77);
	Capítulo 1 \textit{Introducción} (pp.~1--6).}
	{Revisión de trabajos sobre autómatas celulares neuronales: morfogénesis, NCA isotrópicos, autoclasificación,
	representaciones tridimensionales, NCA en malla, modelos generativos y aplicaciones adicionales.
	Actividad concluida en Trabajo Terminal I.}
	{100\%}{}

	\actividad{Marco teórico y metodológico}
	{Fase 1}{Sep/25 -- Oct/25}
	{Documento escrito.}
	{Capítulos 2 a 5 (marco teórico, pp.~7--77) y Capítulo 6 \textit{Marco metodológico} (pp.~78--90).}
	{Marco teórico de sistemas complejos y teoría del caos, autómatas celulares, inteligencia artificial y redes
	neuronales, y autómatas celulares neuronales. Marco metodológico con diseño cuasiexperimental, población,
	muestra, técnicas de recolección y procesamiento de datos, y contextualización tecnológica.
	Actividad concluida en Trabajo Terminal I.}
	{100\%}{}

	\subsection*{Fase 2: Diseño e implementación del simulador}

	\actividad{Arquitectura del simulador}
	{Fase 2}{Oct/25 -- Nov/25}
	{Documento escrito (análisis y diseño de software, diagramas UML).}
	{Capítulo 7 \textit{Análisis de software} (pp.~91--135); Capítulo 8, secciones 8.1 \textit{Arquitectura},
	8.2 \textit{Diagrama de clases}, 8.3 \textit{Diseño de datos} y 8.5 \textit{Interfaces de usuarios} (pp.~136--143, 153--156).}
	{Metodología de desarrollo, reglas de negocio, requerimientos, análisis de riesgos, estudio de factibilidad,
	casos de uso con sus especificaciones, diagramas de secuencia, arquitectura MVC, diagrama de clases,
	diseño de datos e interfaces gráficas.}
	{100\%}{}

	\actividad{Módulos de simulación}
	{Fase 2}{Nov/25 -- Dic/25}
	{Documento escrito, tablas de resultados y pruebas de operatividad.}
	{Sección 8.4 \textit{Diseño de algoritmos} (8.4.1--8.4.4 y 8.4.6, pp.~144--152); secciones 9.1 \textit{Entorno de desarrollo}
	y 9.2 \textit{Estructura del proyecto} (pp.~157--162); subsecciones 9.4.1, 9.4.3, 9.4.4 y 9.4.5 de
	\textit{Validación y pruebas} (pp.~166--169), tablas 9.6 y 9.7.}
	{Núcleo de simulación en C++/CUDA con enlaces a Python e interfaz en Qt: indexación de la regla semitotalística,
	paso de evolución, condiciones de frontera, \textit{forward pass} del NCA y \textit{prefetch} asíncrono.
	Validación funcional contra reglas conocidas y pruebas de rendimiento (hasta 345\,000 pasos/s en
	$100\times100$ y más de 100 fps de renderizado en $2000\times2000$).}
	{100\%}{\videoSimulacion}

	\actividad{Módulos entrenamiento}
	{Fase 2}{Dic/25 -- Ago/26}
	{Documento escrito, pruebas de operatividad.}
	{Subsección 7.8.4 \textit{Especificaciones de casos de uso: Entrenamiento} (pp.~125--130); 8.3.5
	\textit{Configuración del NCA} (p.~143); 8.4.5 \textit{Entrenamiento del autómata celular neuronal} (pp.~150--151);
	sección 9.3 \textit{Entrenamiento del modelo} (pp.~163--166, tabla 9.5); subsecciones 9.4.2 y 9.4.6 (pp.~167, 169).}
	{Módulo de entrenamiento en PyTorch con construcción dinámica del perceptrón multicapa, ciclo de optimización
	con Adam y detención temprana, persistencia de pesos en \texttt{.npz} compatible con el \textit{kernel} CUDA e
	integración con la pestaña de entrenamiento de la GUI (subproceso con \textit{log} en tiempo real).
	Se agregó el parámetro de semilla aleatoria para reproducir experimentos.
	\segun{Actividad concluida en el mes reportado.}{Actividad concluida en agosto de 2026.}}
	{100\%}{\videoEntrenamiento}

	\subsection*{Fase 3: Entrenamiento y Validación}

	\actividad{Entrenamiento y ajuste}
	{Fase 3}{Ago/26 -- Oct/26}
	{Documento escrito y tablas de resultados.}
	{Capítulo 10, sección 10.1 \textit{Duplicación de reglas del juego de la vida}:
	\segun{tablas 10.1 y 10.2 y figura 10.3 (pp.~173--174).}
	{tablas 10.1 y 10.2 y figura 10.3 (pp.~173--174); figuras 10.5 a 10.9 (pp.~176--178).}}
	{\segun{Barrido de hiperparámetros para la regla B3/S23: 10 combinaciones de funciones de activación con
	una capa oculta de 4 neuronas (una sola converge) y con 4 capas ocultas de 8 neuronas (cinco convergen a una
	pérdida menor a $10^{-4}$). Se identifica el estancamiento en las primeras 12 épocas.}
	{Además del barrido de hiperparámetros de agosto, se entrenaron las reglas B3/S23, B3/S238, B38/S23, B3/S236
	y B38/S238 con distintas pérdidas objetivo y semillas, midiendo la distancia Manhattan entre el vector de regla
	aprendido y el original. Se determinó que la pérdida objetivo debe ser $\leq 0.1$ para B3/S23.}}
	{\realEntrenamiento\% (esperado: \espEntrenamiento\%)}{}

	\actividad{Experimentación y análisis}
	{Fase 3}{Ago/26 -- Nov/26}
	{Documento escrito, tablas y gráficas de resultados.}
	{Capítulo 10, sección 10.1:
	\segun{figuras 10.1 y 10.2 (pp.~170--172).}
	{figuras 10.1 a 10.13 (pp.~170--181).}}
	{\segun{Generación del conjunto de datos (100 condiciones iniciales en $100\times100$ durante 100 pasos) y
	análisis de su balance: el estado 1 aparece un $16.5\% \pm 4.2\%$ y los patrones $C0N8$ y $C1N8$ representan
	menos del $0.02\%$ de las entradas.}
	{Primer experimento (duplicación de reglas del juego de la vida) completo. Se diseñó un espacio con las
	512 permutaciones de la vecindad para extraer el vector de regla aprendido y verificar si el NCA conserva la
	propiedad semitotalística. Se identificaron reglas aprendidas alternativas (B3/S23 en lugar de B3/S238,
	B3/S23678 en lugar de B3/S236) y reglas no semitotalísticas.}}
	{\realExperimentos\% (esperado: \espExperimentos\%)}{}

	\actividad{Análisis de resultados (global)}
	{Fase 3}{Ago/26 -- Nov/26}
	{Documento escrito.}
	{Capítulo 10, sección 10.1\segun{ (pp.~170--172).}{, conclusiones del experimento (p.~181).}}
	{\segun{Análisis estadístico de la distribución de estados y de patrones de vecindad para la regla B3/S23,
	identificando el sesgo de muestra y proponiendo sobremuestreo o ponderación de la función de costo.}
	{Conclusiones del primer experimento: la generalización depende del sesgo de muestra de patrones poco frecuentes,
	el entrenamiento de reglas semitotalísticas no garantiza una regla aprendida semitotalística, y no se requieren
	arquitecturas complejas para replicar estas dinámicas. Se valida la metodología de extracción del vector de
	regla para los siguientes experimentos.}}
	{\realAnalisis\% (esperado: \espAnalisis\%)}{}

	\subsection*{Fase 4: Documentación y Difusión}

	\actividad{Documentación técnica}
	{Fase 4}{Oct/26 -- Dic/26}
	{No aplica en el periodo.}
	{Reporte técnico de TT completo (204 páginas).}
	{Actividad programada a partir de octubre de 2026. El reporte técnico se ha redactado de forma incremental
	conforme avanza cada actividad, por lo que su integración final partirá del documento adjunto.}
	{No iniciada (0\% esperado al corte)}{}

	\actividad{Artículo académico}
	{Fase 4}{Oct/26 -- Dic/26}
	{No aplica en el periodo.}
	{--}
	{Actividad programada a partir de octubre de 2026.}
	{No iniciada (0\% esperado al corte)}{}

	\section{Cronograma}

	\resizebox{\linewidth}{!}{%
		\begin{tabular}{p{5cm}*{10}{>{\centering\arraybackslash}p{1cm}}}
			\toprule
			\textbf{Actividad} & \textbf{Ago/25} & \textbf{Sep/25} & \textbf{Oct/25} & \textbf{Nov/25} & \textbf{Dic/25} & \mesActual{1}{Ago/26} & \mesActual{2}{Sep/26} & \textbf{Oct/26} & \textbf{Nov/26} & \textbf{Dic/26} \\
			\midrule
			\multicolumn{11}{l}{\textbf{Fase 1: Planificación y Revisión Bibliográfica}} \\

			Estado del arte       & \cellcolor{gray!30} & \cellcolor{gray!30} & \cellcolor{gray!30} &  &  &  &  &  &  &   \\

			Marco teórico y metodológico &  & \cellcolor{gray!30} & \cellcolor{gray!30} &  &  &  &  &  &  &   \\

			\midrule

			\multicolumn{11}{l}{\textbf{Fase 2: Diseño e implementación del simulador}} \\
			Arquitectura del simulador &  &  & \cellcolor{gray!30} & \cellcolor{gray!30} &  &  &  &  &  &   \\

			Módulos de simulación  &  &  &  & \cellcolor{gray!30} & \cellcolor{gray!30} &  &  &  &  &   \\

			Módulos entrenamiento   &  &  &  &  & \cellcolor{gray!30} & \cellcolor{gray!30} &  &  &  &   \\

			\midrule

			\multicolumn{11}{l}{\textbf{Fase 3: Entrenamiento y Validación}} \\
			Entrenamiento y ajuste   &  &  &  &  &  & \cellcolor{gray!30} & \cellcolor{gray!30} & \cellcolor{gray!30} &   &   \\

			Experimentación y análisis  &  &  &  &  &  & \cellcolor{gray!30}  & \cellcolor{gray!30} & \cellcolor{gray!30} & \cellcolor{gray!30} &   \\

			Análisis de resultados (global) &  &  &  &  &  & \cellcolor{gray!30} & \cellcolor{gray!30} & \cellcolor{gray!30} & \cellcolor{gray!30} &   \\

			\midrule

			\multicolumn{11}{l}{\textbf{Fase 4: Documentación y Difusión}} \\
			Documentación técnica  &  &  &  &  &  &  &  & \cellcolor{gray!30}  & \cellcolor{gray!30} & \cellcolor{gray!30}  \\

			Artículo académico  &  &  &  &  &  &  &  & \cellcolor{gray!30}  & \cellcolor{gray!30}  & \cellcolor{gray!30}  \\

			\bottomrule
		\end{tabular}
	}

	\vspace{0.5em}
	{\small La columna resaltada corresponde al mes reportado.}

	\section{Alumnos y directores}

	\vspace{1em}
	\begin{tabular}{@{} p{0.45\textwidth} p{0.45\textwidth} @{}}
		\textit{Erick Jesse Angeles López} & \textit{Dr. Genaro Juárez Martínez} \\[5em]
		Firma: \rule{6cm}{0.4pt} & Firma: \rule{6cm}{0.4pt} \\
	\end{tabular}

\end{document}
"
%      pdflatex -jobname=2ReporteSeguimiento "\def\corte{2}% ============================================================================
%  Reporte de seguimiento - Trabajo Terminal II
%
%  Un solo archivo genera los reportes de cada mes. El mes se elige con \corte:
%      1 = Agosto 2026, 2 = Septiembre 2026
%
%  Compilación (desde esta carpeta):
%      pdflatex -jobname=1ReporteSeguimiento "\def\corte{1}\input{reporteSeguimiento}"
%      pdflatex -jobname=2ReporteSeguimiento "\def\corte{2}\input{reporteSeguimiento}"
%  o simplemente:  ./compilar.sh
%
%  Las páginas y secciones citadas corresponden a Documentation/main.pdf
%  (reporte técnico de TT), que es el PDF que se adjunta como evidencia.
% ============================================================================
\documentclass[letterpaper,10pt,oneside]{article}
\usepackage[margin=2.54cm]{geometry}
\usepackage[spanish]{babel}
\newcommand\authormark[1]{\textsuperscript{#1}}
\usepackage[colorlinks=true,bookmarks=true,linkcolor=black,citecolor=black,urlcolor=blue]{hyperref}
\usepackage{times}
\usepackage[table]{xcolor}
\usepackage{booktabs}
\usepackage{graphicx}
\usepackage{titlesec}
\usepackage{tabularx}
\usepackage{array}
\usepackage{enumitem}
\usepackage{float}

\titleformat{\section}
{\normalfont\fontsize{12pt}{14pt}\bfseries}
{\thesection}{1em}{}

% ----------------------------------------------------------------------------
%  Mes reportado
% ----------------------------------------------------------------------------
\providecommand\corte{2}
\ifnum\corte=1
	\def\numReporte{1}
	\def\mes{Agosto}
	\def\fechaCorte{31 de agosto de 2026}
\else
	\def\numReporte{2}
	\def\mes{Septiembre}
	\def\fechaCorte{30 de septiembre de 2026}
\fi
% \segun{texto agosto}{texto septiembre}
\newcommand\segun[2]{\ifnum\corte=1 #1\else #2\fi}

% ----------------------------------------------------------------------------
%  Porcentajes (editar aquí). Formato: \segun{agosto}{septiembre}
%  "esp" = avance esperado según cronograma al corte, "real" = avance logrado.
% ----------------------------------------------------------------------------
\def\espEntrenamiento{\segun{33}{67}}   \def\realEntrenamiento{\segun{33}{67}}
\def\espExperimentos{\segun{25}{50}}    \def\realExperimentos{\segun{25}{50}}
\def\espAnalisis{\segun{25}{50}}        \def\realAnalisis{\segun{25}{50}}

% ----------------------------------------------------------------------------
%  Enlaces a video (dejar vacío si no aplica)
% ----------------------------------------------------------------------------
\def\videoSimulacion{}
\def\videoEntrenamiento{}

% ----------------------------------------------------------------------------
%  Bloque de evidencia por actividad
%  \actividad{nombre}{fase}{periodo}{tipo de evidencia}{ubicación en reporte TT}
%            {descripción}{porcentaje}{enlace de video}
% ----------------------------------------------------------------------------
\definecolor{encabezado}{gray}{0.85}
\newcommand\video[1]{\edef\tmp{#1}\ifx\tmp\empty No aplica.\else\url{#1}\fi}
\newcommand\actividad[8]{%
	\par\noindent
	\begin{minipage}{\linewidth}
		\renewcommand{\arraystretch}{1.25}
		\begin{tabularx}{\linewidth}{|>{\raggedright\bfseries\arraybackslash}p{0.24\linewidth}|X|}
			\hline
			\rowcolor{encabezado}
			Nombre de la actividad & \textbf{#1} \\ \hline
			Fase y periodo & #2 \,(#3) \\ \hline
			Evidencia de trabajo & #4 \\ \hline
			Ubicación en el reporte de TT & #5 \\ \hline
			Descripción del avance & #6 \\ \hline
			Porcentaje de cumplimiento & \textbf{#7} \\ \hline
			Enlace de video & \video{#8} \\ \hline
		\end{tabularx}
	\end{minipage}
	\par\vspace{1.2em}%
}

% Porcentaje de cumplimiento: real respecto a lo esperado al corte
\newcommand\cumplimiento[2]{\fpeval{min(100, round(100*(#2)/(#1), 0))}\%}

% Columna del mes reportado en el cronograma
\newcommand\mesActual[2]{\ifnum\corte=#1\cellcolor{black!70}\textcolor{white}{\textbf{#2}}\else\textbf{#2}\fi}

\begin{document}

	\begin{center}
		\fontsize{14}{16}\selectfont\bfseries
		\numReporte° Reporte de seguimiento (\mes{} 2026): \\
		Simulación de sistemas complejos mediante autómatas celulares neuronales
	\end{center}

	\begin{center}
		\fontsize{14}{16}\selectfont\bfseries
		\textit{Trabajo Terminal No. 2026-A114}
	\end{center}

	\begin{center}
		\fontsize{12}{16}\selectfont
		\textit{Alumno: \authormark{*}Angeles López Erick Jesse\authormark{1}} \\
		\textit{Director: Juárez Martínez Genaro\authormark{1}} \\
		\textit{\authormark{*}e-mail: eangelesl1700@alumno.ipn.mx}
	\end{center}

	\section{Información general}

	\begin{itemize}[leftmargin=*]
		\item \textbf{Periodo reportado:} \mes{} 2026 (corte al \fechaCorte).
		\item \textbf{Documento de evidencia adjunto:} reporte técnico de Trabajo Terminal
		\textit{``Simulación de sistemas complejos mediante autómatas celulares neuronales''} (PDF, 204 páginas).
		El avance es incremental: cada evidencia listada en este reporte se encuentra integrada en el capítulo
		y sección correspondiente de dicho documento, cuyas páginas se indican en cada actividad.
		\item \textbf{Actividades programadas en el mes:} \segun{Módulos de entrenamiento (cierre), Entrenamiento y ajuste,
		Experimentación y análisis, Análisis de resultados (global).}{Entrenamiento y ajuste, Experimentación y análisis,
		Análisis de resultados (global).}
	\end{itemize}

	\section{Resumen de cumplimiento}

	\begin{table}[H]
		\centering
		\renewcommand{\arraystretch}{1.2}
		\begin{tabular}{l c c c c}
			\toprule
			\textbf{Actividad} & \textbf{Periodo} & \textbf{Esperado al corte} & \textbf{Avance real} & \textbf{Cumplimiento} \\
			\midrule
			\multicolumn{5}{l}{\textit{Fase 1: Planificación y Revisión Bibliográfica}} \\
			Estado del arte                  & Ago/25 -- Oct/25 & 100\% & 100\% & 100\% \\
			Marco teórico y metodológico     & Sep/25 -- Oct/25 & 100\% & 100\% & 100\% \\
			\multicolumn{5}{l}{\textit{Fase 2: Diseño e implementación del simulador}} \\
			Arquitectura del simulador       & Oct/25 -- Nov/25 & 100\% & 100\% & 100\% \\
			Módulos de simulación            & Nov/25 -- Dic/25 & 100\% & 100\% & 100\% \\
			Módulos entrenamiento            & Dic/25 -- Ago/26 & 100\% & 100\% & 100\% \\
			\multicolumn{5}{l}{\textit{Fase 3: Entrenamiento y Validación}} \\
			Entrenamiento y ajuste           & Ago/26 -- Oct/26 & \espEntrenamiento\% & \realEntrenamiento\% & \cumplimiento{\espEntrenamiento}{\realEntrenamiento} \\
			Experimentación y análisis       & Ago/26 -- Nov/26 & \espExperimentos\%  & \realExperimentos\%  & \cumplimiento{\espExperimentos}{\realExperimentos} \\
			Análisis de resultados (global)  & Ago/26 -- Nov/26 & \espAnalisis\%      & \realAnalisis\%      & \cumplimiento{\espAnalisis}{\realAnalisis} \\
			\multicolumn{5}{l}{\textit{Fase 4: Documentación y Difusión}} \\
			Documentación técnica            & Oct/26 -- Dic/26 & 0\% & -- & No iniciada \\
			Artículo académico               & Oct/26 -- Dic/26 & 0\% & -- & No iniciada \\
			\bottomrule
		\end{tabular}
		\caption{Avance acumulado por actividad al \fechaCorte. ``Esperado al corte'' es la proporción del periodo
		programado en el cronograma que ha transcurrido; ``Cumplimiento'' compara el avance real con dicho valor.}
	\end{table}

	\section{Evidencias por actividad}

	\subsection*{Fase 1: Planificación y Revisión Bibliográfica}

	\actividad{Estado del arte}
	{Fase 1}{Ago/25 -- Oct/25}
	{Documento escrito.}
	{Capítulo 5, sección 5.1 \textit{Estado del arte, aplicaciones y líneas actuales de investigación} (pp.~61--77);
	Capítulo 1 \textit{Introducción} (pp.~1--6).}
	{Revisión de trabajos sobre autómatas celulares neuronales: morfogénesis, NCA isotrópicos, autoclasificación,
	representaciones tridimensionales, NCA en malla, modelos generativos y aplicaciones adicionales.
	Actividad concluida en Trabajo Terminal I.}
	{100\%}{}

	\actividad{Marco teórico y metodológico}
	{Fase 1}{Sep/25 -- Oct/25}
	{Documento escrito.}
	{Capítulos 2 a 5 (marco teórico, pp.~7--77) y Capítulo 6 \textit{Marco metodológico} (pp.~78--90).}
	{Marco teórico de sistemas complejos y teoría del caos, autómatas celulares, inteligencia artificial y redes
	neuronales, y autómatas celulares neuronales. Marco metodológico con diseño cuasiexperimental, población,
	muestra, técnicas de recolección y procesamiento de datos, y contextualización tecnológica.
	Actividad concluida en Trabajo Terminal I.}
	{100\%}{}

	\subsection*{Fase 2: Diseño e implementación del simulador}

	\actividad{Arquitectura del simulador}
	{Fase 2}{Oct/25 -- Nov/25}
	{Documento escrito (análisis y diseño de software, diagramas UML).}
	{Capítulo 7 \textit{Análisis de software} (pp.~91--135); Capítulo 8, secciones 8.1 \textit{Arquitectura},
	8.2 \textit{Diagrama de clases}, 8.3 \textit{Diseño de datos} y 8.5 \textit{Interfaces de usuarios} (pp.~136--143, 153--156).}
	{Metodología de desarrollo, reglas de negocio, requerimientos, análisis de riesgos, estudio de factibilidad,
	casos de uso con sus especificaciones, diagramas de secuencia, arquitectura MVC, diagrama de clases,
	diseño de datos e interfaces gráficas.}
	{100\%}{}

	\actividad{Módulos de simulación}
	{Fase 2}{Nov/25 -- Dic/25}
	{Documento escrito, tablas de resultados y pruebas de operatividad.}
	{Sección 8.4 \textit{Diseño de algoritmos} (8.4.1--8.4.4 y 8.4.6, pp.~144--152); secciones 9.1 \textit{Entorno de desarrollo}
	y 9.2 \textit{Estructura del proyecto} (pp.~157--162); subsecciones 9.4.1, 9.4.3, 9.4.4 y 9.4.5 de
	\textit{Validación y pruebas} (pp.~166--169), tablas 9.6 y 9.7.}
	{Núcleo de simulación en C++/CUDA con enlaces a Python e interfaz en Qt: indexación de la regla semitotalística,
	paso de evolución, condiciones de frontera, \textit{forward pass} del NCA y \textit{prefetch} asíncrono.
	Validación funcional contra reglas conocidas y pruebas de rendimiento (hasta 345\,000 pasos/s en
	$100\times100$ y más de 100 fps de renderizado en $2000\times2000$).}
	{100\%}{\videoSimulacion}

	\actividad{Módulos entrenamiento}
	{Fase 2}{Dic/25 -- Ago/26}
	{Documento escrito, pruebas de operatividad.}
	{Subsección 7.8.4 \textit{Especificaciones de casos de uso: Entrenamiento} (pp.~125--130); 8.3.5
	\textit{Configuración del NCA} (p.~143); 8.4.5 \textit{Entrenamiento del autómata celular neuronal} (pp.~150--151);
	sección 9.3 \textit{Entrenamiento del modelo} (pp.~163--166, tabla 9.5); subsecciones 9.4.2 y 9.4.6 (pp.~167, 169).}
	{Módulo de entrenamiento en PyTorch con construcción dinámica del perceptrón multicapa, ciclo de optimización
	con Adam y detención temprana, persistencia de pesos en \texttt{.npz} compatible con el \textit{kernel} CUDA e
	integración con la pestaña de entrenamiento de la GUI (subproceso con \textit{log} en tiempo real).
	Se agregó el parámetro de semilla aleatoria para reproducir experimentos.
	\segun{Actividad concluida en el mes reportado.}{Actividad concluida en agosto de 2026.}}
	{100\%}{\videoEntrenamiento}

	\subsection*{Fase 3: Entrenamiento y Validación}

	\actividad{Entrenamiento y ajuste}
	{Fase 3}{Ago/26 -- Oct/26}
	{Documento escrito y tablas de resultados.}
	{Capítulo 10, sección 10.1 \textit{Duplicación de reglas del juego de la vida}:
	\segun{tablas 10.1 y 10.2 y figura 10.3 (pp.~173--174).}
	{tablas 10.1 y 10.2 y figura 10.3 (pp.~173--174); figuras 10.5 a 10.9 (pp.~176--178).}}
	{\segun{Barrido de hiperparámetros para la regla B3/S23: 10 combinaciones de funciones de activación con
	una capa oculta de 4 neuronas (una sola converge) y con 4 capas ocultas de 8 neuronas (cinco convergen a una
	pérdida menor a $10^{-4}$). Se identifica el estancamiento en las primeras 12 épocas.}
	{Además del barrido de hiperparámetros de agosto, se entrenaron las reglas B3/S23, B3/S238, B38/S23, B3/S236
	y B38/S238 con distintas pérdidas objetivo y semillas, midiendo la distancia Manhattan entre el vector de regla
	aprendido y el original. Se determinó que la pérdida objetivo debe ser $\leq 0.1$ para B3/S23.}}
	{\realEntrenamiento\% (esperado: \espEntrenamiento\%)}{}

	\actividad{Experimentación y análisis}
	{Fase 3}{Ago/26 -- Nov/26}
	{Documento escrito, tablas y gráficas de resultados.}
	{Capítulo 10, sección 10.1:
	\segun{figuras 10.1 y 10.2 (pp.~170--172).}
	{figuras 10.1 a 10.13 (pp.~170--181).}}
	{\segun{Generación del conjunto de datos (100 condiciones iniciales en $100\times100$ durante 100 pasos) y
	análisis de su balance: el estado 1 aparece un $16.5\% \pm 4.2\%$ y los patrones $C0N8$ y $C1N8$ representan
	menos del $0.02\%$ de las entradas.}
	{Primer experimento (duplicación de reglas del juego de la vida) completo. Se diseñó un espacio con las
	512 permutaciones de la vecindad para extraer el vector de regla aprendido y verificar si el NCA conserva la
	propiedad semitotalística. Se identificaron reglas aprendidas alternativas (B3/S23 en lugar de B3/S238,
	B3/S23678 en lugar de B3/S236) y reglas no semitotalísticas.}}
	{\realExperimentos\% (esperado: \espExperimentos\%)}{}

	\actividad{Análisis de resultados (global)}
	{Fase 3}{Ago/26 -- Nov/26}
	{Documento escrito.}
	{Capítulo 10, sección 10.1\segun{ (pp.~170--172).}{, conclusiones del experimento (p.~181).}}
	{\segun{Análisis estadístico de la distribución de estados y de patrones de vecindad para la regla B3/S23,
	identificando el sesgo de muestra y proponiendo sobremuestreo o ponderación de la función de costo.}
	{Conclusiones del primer experimento: la generalización depende del sesgo de muestra de patrones poco frecuentes,
	el entrenamiento de reglas semitotalísticas no garantiza una regla aprendida semitotalística, y no se requieren
	arquitecturas complejas para replicar estas dinámicas. Se valida la metodología de extracción del vector de
	regla para los siguientes experimentos.}}
	{\realAnalisis\% (esperado: \espAnalisis\%)}{}

	\subsection*{Fase 4: Documentación y Difusión}

	\actividad{Documentación técnica}
	{Fase 4}{Oct/26 -- Dic/26}
	{No aplica en el periodo.}
	{Reporte técnico de TT completo (204 páginas).}
	{Actividad programada a partir de octubre de 2026. El reporte técnico se ha redactado de forma incremental
	conforme avanza cada actividad, por lo que su integración final partirá del documento adjunto.}
	{No iniciada (0\% esperado al corte)}{}

	\actividad{Artículo académico}
	{Fase 4}{Oct/26 -- Dic/26}
	{No aplica en el periodo.}
	{--}
	{Actividad programada a partir de octubre de 2026.}
	{No iniciada (0\% esperado al corte)}{}

	\section{Cronograma}

	\resizebox{\linewidth}{!}{%
		\begin{tabular}{p{5cm}*{10}{>{\centering\arraybackslash}p{1cm}}}
			\toprule
			\textbf{Actividad} & \textbf{Ago/25} & \textbf{Sep/25} & \textbf{Oct/25} & \textbf{Nov/25} & \textbf{Dic/25} & \mesActual{1}{Ago/26} & \mesActual{2}{Sep/26} & \textbf{Oct/26} & \textbf{Nov/26} & \textbf{Dic/26} \\
			\midrule
			\multicolumn{11}{l}{\textbf{Fase 1: Planificación y Revisión Bibliográfica}} \\

			Estado del arte       & \cellcolor{gray!30} & \cellcolor{gray!30} & \cellcolor{gray!30} &  &  &  &  &  &  &   \\

			Marco teórico y metodológico &  & \cellcolor{gray!30} & \cellcolor{gray!30} &  &  &  &  &  &  &   \\

			\midrule

			\multicolumn{11}{l}{\textbf{Fase 2: Diseño e implementación del simulador}} \\
			Arquitectura del simulador &  &  & \cellcolor{gray!30} & \cellcolor{gray!30} &  &  &  &  &  &   \\

			Módulos de simulación  &  &  &  & \cellcolor{gray!30} & \cellcolor{gray!30} &  &  &  &  &   \\

			Módulos entrenamiento   &  &  &  &  & \cellcolor{gray!30} & \cellcolor{gray!30} &  &  &  &   \\

			\midrule

			\multicolumn{11}{l}{\textbf{Fase 3: Entrenamiento y Validación}} \\
			Entrenamiento y ajuste   &  &  &  &  &  & \cellcolor{gray!30} & \cellcolor{gray!30} & \cellcolor{gray!30} &   &   \\

			Experimentación y análisis  &  &  &  &  &  & \cellcolor{gray!30}  & \cellcolor{gray!30} & \cellcolor{gray!30} & \cellcolor{gray!30} &   \\

			Análisis de resultados (global) &  &  &  &  &  & \cellcolor{gray!30} & \cellcolor{gray!30} & \cellcolor{gray!30} & \cellcolor{gray!30} &   \\

			\midrule

			\multicolumn{11}{l}{\textbf{Fase 4: Documentación y Difusión}} \\
			Documentación técnica  &  &  &  &  &  &  &  & \cellcolor{gray!30}  & \cellcolor{gray!30} & \cellcolor{gray!30}  \\

			Artículo académico  &  &  &  &  &  &  &  & \cellcolor{gray!30}  & \cellcolor{gray!30}  & \cellcolor{gray!30}  \\

			\bottomrule
		\end{tabular}
	}

	\vspace{0.5em}
	{\small La columna resaltada corresponde al mes reportado.}

	\section{Alumnos y directores}

	\vspace{1em}
	\begin{tabular}{@{} p{0.45\textwidth} p{0.45\textwidth} @{}}
		\textit{Erick Jesse Angeles López} & \textit{Dr. Genaro Juárez Martínez} \\[5em]
		Firma: \rule{6cm}{0.4pt} & Firma: \rule{6cm}{0.4pt} \\
	\end{tabular}

\end{document}
"
%  o simplemente:  ./compilar.sh
%
%  Las páginas y secciones citadas corresponden a Documentation/main.pdf
%  (reporte técnico de TT), que es el PDF que se adjunta como evidencia.
% ============================================================================
\documentclass[letterpaper,10pt,oneside]{article}
\usepackage[margin=2.54cm]{geometry}
\usepackage[spanish]{babel}
\newcommand\authormark[1]{\textsuperscript{#1}}
\usepackage[colorlinks=true,bookmarks=true,linkcolor=black,citecolor=black,urlcolor=blue]{hyperref}
\usepackage{times}
\usepackage[table]{xcolor}
\usepackage{booktabs}
\usepackage{graphicx}
\usepackage{titlesec}
\usepackage{tabularx}
\usepackage{array}
\usepackage{enumitem}
\usepackage{float}

\titleformat{\section}
{\normalfont\fontsize{12pt}{14pt}\bfseries}
{\thesection}{1em}{}

% ----------------------------------------------------------------------------
%  Mes reportado
% ----------------------------------------------------------------------------
\providecommand\corte{2}
\ifnum\corte=1
	\def\numReporte{1}
	\def\mes{Agosto}
	\def\fechaCorte{31 de agosto de 2026}
\else
	\def\numReporte{2}
	\def\mes{Septiembre}
	\def\fechaCorte{30 de septiembre de 2026}
\fi
% \segun{texto agosto}{texto septiembre}
\newcommand\segun[2]{\ifnum\corte=1 #1\else #2\fi}

% ----------------------------------------------------------------------------
%  Porcentajes (editar aquí). Formato: \segun{agosto}{septiembre}
%  "esp" = avance esperado según cronograma al corte, "real" = avance logrado.
% ----------------------------------------------------------------------------
\def\espEntrenamiento{\segun{33}{67}}   \def\realEntrenamiento{\segun{33}{67}}
\def\espExperimentos{\segun{25}{50}}    \def\realExperimentos{\segun{25}{50}}
\def\espAnalisis{\segun{25}{50}}        \def\realAnalisis{\segun{25}{50}}

% ----------------------------------------------------------------------------
%  Enlaces a video (dejar vacío si no aplica)
% ----------------------------------------------------------------------------
\def\videoSimulacion{}
\def\videoEntrenamiento{}

% ----------------------------------------------------------------------------
%  Bloque de evidencia por actividad
%  \actividad{nombre}{fase}{periodo}{tipo de evidencia}{ubicación en reporte TT}
%            {descripción}{porcentaje}{enlace de video}
% ----------------------------------------------------------------------------
\definecolor{encabezado}{gray}{0.85}
\newcommand\video[1]{\edef\tmp{#1}\ifx\tmp\empty No aplica.\else\url{#1}\fi}
\newcommand\actividad[8]{%
	\par\noindent
	\begin{minipage}{\linewidth}
		\renewcommand{\arraystretch}{1.25}
		\begin{tabularx}{\linewidth}{|>{\raggedright\bfseries\arraybackslash}p{0.24\linewidth}|X|}
			\hline
			\rowcolor{encabezado}
			Nombre de la actividad & \textbf{#1} \\ \hline
			Fase y periodo & #2 \,(#3) \\ \hline
			Evidencia de trabajo & #4 \\ \hline
			Ubicación en el reporte de TT & #5 \\ \hline
			Descripción del avance & #6 \\ \hline
			Porcentaje de cumplimiento & \textbf{#7} \\ \hline
			Enlace de video & \video{#8} \\ \hline
		\end{tabularx}
	\end{minipage}
	\par\vspace{1.2em}%
}

% Porcentaje de cumplimiento: real respecto a lo esperado al corte
\newcommand\cumplimiento[2]{\fpeval{min(100, round(100*(#2)/(#1), 0))}\%}

% Columna del mes reportado en el cronograma
\newcommand\mesActual[2]{\ifnum\corte=#1\cellcolor{black!70}\textcolor{white}{\textbf{#2}}\else\textbf{#2}\fi}

\begin{document}

	\begin{center}
		\fontsize{14}{16}\selectfont\bfseries
		\numReporte° Reporte de seguimiento (\mes{} 2026): \\
		Simulación de sistemas complejos mediante autómatas celulares neuronales
	\end{center}

	\begin{center}
		\fontsize{14}{16}\selectfont\bfseries
		\textit{Trabajo Terminal No. 2026-A114}
	\end{center}

	\begin{center}
		\fontsize{12}{16}\selectfont
		\textit{Alumno: \authormark{*}Angeles López Erick Jesse\authormark{1}} \\
		\textit{Director: Juárez Martínez Genaro\authormark{1}} \\
		\textit{\authormark{*}e-mail: eangelesl1700@alumno.ipn.mx}
	\end{center}

	\section{Información general}

	\begin{itemize}[leftmargin=*]
		\item \textbf{Periodo reportado:} \mes{} 2026 (corte al \fechaCorte).
		\item \textbf{Documento de evidencia adjunto:} reporte técnico de Trabajo Terminal
		\textit{``Simulación de sistemas complejos mediante autómatas celulares neuronales''} (PDF, 204 páginas).
		El avance es incremental: cada evidencia listada en este reporte se encuentra integrada en el capítulo
		y sección correspondiente de dicho documento, cuyas páginas se indican en cada actividad.
		\item \textbf{Actividades programadas en el mes:} \segun{Módulos de entrenamiento (cierre), Entrenamiento y ajuste,
		Experimentación y análisis, Análisis de resultados (global).}{Entrenamiento y ajuste, Experimentación y análisis,
		Análisis de resultados (global).}
	\end{itemize}

	\section{Resumen de cumplimiento}

	\begin{table}[H]
		\centering
		\renewcommand{\arraystretch}{1.2}
		\begin{tabular}{l c c c c}
			\toprule
			\textbf{Actividad} & \textbf{Periodo} & \textbf{Esperado al corte} & \textbf{Avance real} & \textbf{Cumplimiento} \\
			\midrule
			\multicolumn{5}{l}{\textit{Fase 1: Planificación y Revisión Bibliográfica}} \\
			Estado del arte                  & Ago/25 -- Oct/25 & 100\% & 100\% & 100\% \\
			Marco teórico y metodológico     & Sep/25 -- Oct/25 & 100\% & 100\% & 100\% \\
			\multicolumn{5}{l}{\textit{Fase 2: Diseño e implementación del simulador}} \\
			Arquitectura del simulador       & Oct/25 -- Nov/25 & 100\% & 100\% & 100\% \\
			Módulos de simulación            & Nov/25 -- Dic/25 & 100\% & 100\% & 100\% \\
			Módulos entrenamiento            & Dic/25 -- Ago/26 & 100\% & 100\% & 100\% \\
			\multicolumn{5}{l}{\textit{Fase 3: Entrenamiento y Validación}} \\
			Entrenamiento y ajuste           & Ago/26 -- Oct/26 & \espEntrenamiento\% & \realEntrenamiento\% & \cumplimiento{\espEntrenamiento}{\realEntrenamiento} \\
			Experimentación y análisis       & Ago/26 -- Nov/26 & \espExperimentos\%  & \realExperimentos\%  & \cumplimiento{\espExperimentos}{\realExperimentos} \\
			Análisis de resultados (global)  & Ago/26 -- Nov/26 & \espAnalisis\%      & \realAnalisis\%      & \cumplimiento{\espAnalisis}{\realAnalisis} \\
			\multicolumn{5}{l}{\textit{Fase 4: Documentación y Difusión}} \\
			Documentación técnica            & Oct/26 -- Dic/26 & 0\% & -- & No iniciada \\
			Artículo académico               & Oct/26 -- Dic/26 & 0\% & -- & No iniciada \\
			\bottomrule
		\end{tabular}
		\caption{Avance acumulado por actividad al \fechaCorte. ``Esperado al corte'' es la proporción del periodo
		programado en el cronograma que ha transcurrido; ``Cumplimiento'' compara el avance real con dicho valor.}
	\end{table}

	\section{Evidencias por actividad}

	\subsection*{Fase 1: Planificación y Revisión Bibliográfica}

	\actividad{Estado del arte}
	{Fase 1}{Ago/25 -- Oct/25}
	{Documento escrito.}
	{Capítulo 5, sección 5.1 \textit{Estado del arte, aplicaciones y líneas actuales de investigación} (pp.~61--77);
	Capítulo 1 \textit{Introducción} (pp.~1--6).}
	{Revisión de trabajos sobre autómatas celulares neuronales: morfogénesis, NCA isotrópicos, autoclasificación,
	representaciones tridimensionales, NCA en malla, modelos generativos y aplicaciones adicionales.
	Actividad concluida en Trabajo Terminal I.}
	{100\%}{}

	\actividad{Marco teórico y metodológico}
	{Fase 1}{Sep/25 -- Oct/25}
	{Documento escrito.}
	{Capítulos 2 a 5 (marco teórico, pp.~7--77) y Capítulo 6 \textit{Marco metodológico} (pp.~78--90).}
	{Marco teórico de sistemas complejos y teoría del caos, autómatas celulares, inteligencia artificial y redes
	neuronales, y autómatas celulares neuronales. Marco metodológico con diseño cuasiexperimental, población,
	muestra, técnicas de recolección y procesamiento de datos, y contextualización tecnológica.
	Actividad concluida en Trabajo Terminal I.}
	{100\%}{}

	\subsection*{Fase 2: Diseño e implementación del simulador}

	\actividad{Arquitectura del simulador}
	{Fase 2}{Oct/25 -- Nov/25}
	{Documento escrito (análisis y diseño de software, diagramas UML).}
	{Capítulo 7 \textit{Análisis de software} (pp.~91--135); Capítulo 8, secciones 8.1 \textit{Arquitectura},
	8.2 \textit{Diagrama de clases}, 8.3 \textit{Diseño de datos} y 8.5 \textit{Interfaces de usuarios} (pp.~136--143, 153--156).}
	{Metodología de desarrollo, reglas de negocio, requerimientos, análisis de riesgos, estudio de factibilidad,
	casos de uso con sus especificaciones, diagramas de secuencia, arquitectura MVC, diagrama de clases,
	diseño de datos e interfaces gráficas.}
	{100\%}{}

	\actividad{Módulos de simulación}
	{Fase 2}{Nov/25 -- Dic/25}
	{Documento escrito, tablas de resultados y pruebas de operatividad.}
	{Sección 8.4 \textit{Diseño de algoritmos} (8.4.1--8.4.4 y 8.4.6, pp.~144--152); secciones 9.1 \textit{Entorno de desarrollo}
	y 9.2 \textit{Estructura del proyecto} (pp.~157--162); subsecciones 9.4.1, 9.4.3, 9.4.4 y 9.4.5 de
	\textit{Validación y pruebas} (pp.~166--169), tablas 9.6 y 9.7.}
	{Núcleo de simulación en C++/CUDA con enlaces a Python e interfaz en Qt: indexación de la regla semitotalística,
	paso de evolución, condiciones de frontera, \textit{forward pass} del NCA y \textit{prefetch} asíncrono.
	Validación funcional contra reglas conocidas y pruebas de rendimiento (hasta 345\,000 pasos/s en
	$100\times100$ y más de 100 fps de renderizado en $2000\times2000$).}
	{100\%}{\videoSimulacion}

	\actividad{Módulos entrenamiento}
	{Fase 2}{Dic/25 -- Ago/26}
	{Documento escrito, pruebas de operatividad.}
	{Subsección 7.8.4 \textit{Especificaciones de casos de uso: Entrenamiento} (pp.~125--130); 8.3.5
	\textit{Configuración del NCA} (p.~143); 8.4.5 \textit{Entrenamiento del autómata celular neuronal} (pp.~150--151);
	sección 9.3 \textit{Entrenamiento del modelo} (pp.~163--166, tabla 9.5); subsecciones 9.4.2 y 9.4.6 (pp.~167, 169).}
	{Módulo de entrenamiento en PyTorch con construcción dinámica del perceptrón multicapa, ciclo de optimización
	con Adam y detención temprana, persistencia de pesos en \texttt{.npz} compatible con el \textit{kernel} CUDA e
	integración con la pestaña de entrenamiento de la GUI (subproceso con \textit{log} en tiempo real).
	Se agregó el parámetro de semilla aleatoria para reproducir experimentos.
	\segun{Actividad concluida en el mes reportado.}{Actividad concluida en agosto de 2026.}}
	{100\%}{\videoEntrenamiento}

	\subsection*{Fase 3: Entrenamiento y Validación}

	\actividad{Entrenamiento y ajuste}
	{Fase 3}{Ago/26 -- Oct/26}
	{Documento escrito y tablas de resultados.}
	{Capítulo 10, sección 10.1 \textit{Duplicación de reglas del juego de la vida}:
	\segun{tablas 10.1 y 10.2 y figura 10.3 (pp.~173--174).}
	{tablas 10.1 y 10.2 y figura 10.3 (pp.~173--174); figuras 10.5 a 10.9 (pp.~176--178).}}
	{\segun{Barrido de hiperparámetros para la regla B3/S23: 10 combinaciones de funciones de activación con
	una capa oculta de 4 neuronas (una sola converge) y con 4 capas ocultas de 8 neuronas (cinco convergen a una
	pérdida menor a $10^{-4}$). Se identifica el estancamiento en las primeras 12 épocas.}
	{Además del barrido de hiperparámetros de agosto, se entrenaron las reglas B3/S23, B3/S238, B38/S23, B3/S236
	y B38/S238 con distintas pérdidas objetivo y semillas, midiendo la distancia Manhattan entre el vector de regla
	aprendido y el original. Se determinó que la pérdida objetivo debe ser $\leq 0.1$ para B3/S23.}}
	{\realEntrenamiento\% (esperado: \espEntrenamiento\%)}{}

	\actividad{Experimentación y análisis}
	{Fase 3}{Ago/26 -- Nov/26}
	{Documento escrito, tablas y gráficas de resultados.}
	{Capítulo 10, sección 10.1:
	\segun{figuras 10.1 y 10.2 (pp.~170--172).}
	{figuras 10.1 a 10.13 (pp.~170--181).}}
	{\segun{Generación del conjunto de datos (100 condiciones iniciales en $100\times100$ durante 100 pasos) y
	análisis de su balance: el estado 1 aparece un $16.5\% \pm 4.2\%$ y los patrones $C0N8$ y $C1N8$ representan
	menos del $0.02\%$ de las entradas.}
	{Primer experimento (duplicación de reglas del juego de la vida) completo. Se diseñó un espacio con las
	512 permutaciones de la vecindad para extraer el vector de regla aprendido y verificar si el NCA conserva la
	propiedad semitotalística. Se identificaron reglas aprendidas alternativas (B3/S23 en lugar de B3/S238,
	B3/S23678 en lugar de B3/S236) y reglas no semitotalísticas.}}
	{\realExperimentos\% (esperado: \espExperimentos\%)}{}

	\actividad{Análisis de resultados (global)}
	{Fase 3}{Ago/26 -- Nov/26}
	{Documento escrito.}
	{Capítulo 10, sección 10.1\segun{ (pp.~170--172).}{, conclusiones del experimento (p.~181).}}
	{\segun{Análisis estadístico de la distribución de estados y de patrones de vecindad para la regla B3/S23,
	identificando el sesgo de muestra y proponiendo sobremuestreo o ponderación de la función de costo.}
	{Conclusiones del primer experimento: la generalización depende del sesgo de muestra de patrones poco frecuentes,
	el entrenamiento de reglas semitotalísticas no garantiza una regla aprendida semitotalística, y no se requieren
	arquitecturas complejas para replicar estas dinámicas. Se valida la metodología de extracción del vector de
	regla para los siguientes experimentos.}}
	{\realAnalisis\% (esperado: \espAnalisis\%)}{}

	\subsection*{Fase 4: Documentación y Difusión}

	\actividad{Documentación técnica}
	{Fase 4}{Oct/26 -- Dic/26}
	{No aplica en el periodo.}
	{Reporte técnico de TT completo (204 páginas).}
	{Actividad programada a partir de octubre de 2026. El reporte técnico se ha redactado de forma incremental
	conforme avanza cada actividad, por lo que su integración final partirá del documento adjunto.}
	{No iniciada (0\% esperado al corte)}{}

	\actividad{Artículo académico}
	{Fase 4}{Oct/26 -- Dic/26}
	{No aplica en el periodo.}
	{--}
	{Actividad programada a partir de octubre de 2026.}
	{No iniciada (0\% esperado al corte)}{}

	\section{Cronograma}

	\resizebox{\linewidth}{!}{%
		\begin{tabular}{p{5cm}*{10}{>{\centering\arraybackslash}p{1cm}}}
			\toprule
			\textbf{Actividad} & \textbf{Ago/25} & \textbf{Sep/25} & \textbf{Oct/25} & \textbf{Nov/25} & \textbf{Dic/25} & \mesActual{1}{Ago/26} & \mesActual{2}{Sep/26} & \textbf{Oct/26} & \textbf{Nov/26} & \textbf{Dic/26} \\
			\midrule
			\multicolumn{11}{l}{\textbf{Fase 1: Planificación y Revisión Bibliográfica}} \\

			Estado del arte       & \cellcolor{gray!30} & \cellcolor{gray!30} & \cellcolor{gray!30} &  &  &  &  &  &  &   \\

			Marco teórico y metodológico &  & \cellcolor{gray!30} & \cellcolor{gray!30} &  &  &  &  &  &  &   \\

			\midrule

			\multicolumn{11}{l}{\textbf{Fase 2: Diseño e implementación del simulador}} \\
			Arquitectura del simulador &  &  & \cellcolor{gray!30} & \cellcolor{gray!30} &  &  &  &  &  &   \\

			Módulos de simulación  &  &  &  & \cellcolor{gray!30} & \cellcolor{gray!30} &  &  &  &  &   \\

			Módulos entrenamiento   &  &  &  &  & \cellcolor{gray!30} & \cellcolor{gray!30} &  &  &  &   \\

			\midrule

			\multicolumn{11}{l}{\textbf{Fase 3: Entrenamiento y Validación}} \\
			Entrenamiento y ajuste   &  &  &  &  &  & \cellcolor{gray!30} & \cellcolor{gray!30} & \cellcolor{gray!30} &   &   \\

			Experimentación y análisis  &  &  &  &  &  & \cellcolor{gray!30}  & \cellcolor{gray!30} & \cellcolor{gray!30} & \cellcolor{gray!30} &   \\

			Análisis de resultados (global) &  &  &  &  &  & \cellcolor{gray!30} & \cellcolor{gray!30} & \cellcolor{gray!30} & \cellcolor{gray!30} &   \\

			\midrule

			\multicolumn{11}{l}{\textbf{Fase 4: Documentación y Difusión}} \\
			Documentación técnica  &  &  &  &  &  &  &  & \cellcolor{gray!30}  & \cellcolor{gray!30} & \cellcolor{gray!30}  \\

			Artículo académico  &  &  &  &  &  &  &  & \cellcolor{gray!30}  & \cellcolor{gray!30}  & \cellcolor{gray!30}  \\

			\bottomrule
		\end{tabular}
	}

	\vspace{0.5em}
	{\small La columna resaltada corresponde al mes reportado.}

	\section{Alumnos y directores}

	\vspace{1em}
	\begin{tabular}{@{} p{0.45\textwidth} p{0.45\textwidth} @{}}
		\textit{Erick Jesse Angeles López} & \textit{Dr. Genaro Juárez Martínez} \\[5em]
		Firma: \rule{6cm}{0.4pt} & Firma: \rule{6cm}{0.4pt} \\
	\end{tabular}

\end{document}
"
%  o simplemente:  ./compilar.sh
%
%  Las páginas y secciones citadas corresponden a Documentation/main.pdf
%  (reporte técnico de TT), que es el PDF que se adjunta como evidencia.
% ============================================================================
\documentclass[letterpaper,10pt,oneside]{article}
\usepackage[margin=2.54cm]{geometry}
\usepackage[spanish]{babel}
\newcommand\authormark[1]{\textsuperscript{#1}}
\usepackage[colorlinks=true,bookmarks=true,linkcolor=black,citecolor=black,urlcolor=blue]{hyperref}
\usepackage{times}
\usepackage[table]{xcolor}
\usepackage{booktabs}
\usepackage{graphicx}
\usepackage{titlesec}
\usepackage{tabularx}
\usepackage{array}
\usepackage{enumitem}
\usepackage{float}

\titleformat{\section}
{\normalfont\fontsize{12pt}{14pt}\bfseries}
{\thesection}{1em}{}

% ----------------------------------------------------------------------------
%  Mes reportado
% ----------------------------------------------------------------------------
\providecommand\corte{2}
\ifnum\corte=1
	\def\numReporte{1}
	\def\mes{Agosto}
	\def\fechaCorte{31 de agosto de 2026}
\else
	\def\numReporte{2}
	\def\mes{Septiembre}
	\def\fechaCorte{30 de septiembre de 2026}
\fi
% \segun{texto agosto}{texto septiembre}
\newcommand\segun[2]{\ifnum\corte=1 #1\else #2\fi}

% ----------------------------------------------------------------------------
%  Porcentajes (editar aquí). Formato: \segun{agosto}{septiembre}
%  "esp" = avance esperado según cronograma al corte, "real" = avance logrado.
% ----------------------------------------------------------------------------
\def\espEntrenamiento{\segun{33}{67}}   \def\realEntrenamiento{\segun{33}{67}}
\def\espExperimentos{\segun{25}{50}}    \def\realExperimentos{\segun{25}{50}}
\def\espAnalisis{\segun{25}{50}}        \def\realAnalisis{\segun{25}{50}}

% ----------------------------------------------------------------------------
%  Enlaces a video (dejar vacío si no aplica)
% ----------------------------------------------------------------------------
\def\videoSimulacion{}
\def\videoEntrenamiento{}

% ----------------------------------------------------------------------------
%  Bloque de evidencia por actividad
%  \actividad{nombre}{fase}{periodo}{tipo de evidencia}{ubicación en reporte TT}
%            {descripción}{porcentaje}{enlace de video}
% ----------------------------------------------------------------------------
\definecolor{encabezado}{gray}{0.85}
\newcommand\video[1]{\edef\tmp{#1}\ifx\tmp\empty No aplica.\else\url{#1}\fi}
\newcommand\actividad[8]{%
	\par\noindent
	\begin{minipage}{\linewidth}
		\renewcommand{\arraystretch}{1.25}
		\begin{tabularx}{\linewidth}{|>{\raggedright\bfseries\arraybackslash}p{0.24\linewidth}|X|}
			\hline
			\rowcolor{encabezado}
			Nombre de la actividad & \textbf{#1} \\ \hline
			Fase y periodo & #2 \,(#3) \\ \hline
			Evidencia de trabajo & #4 \\ \hline
			Ubicación en el reporte de TT & #5 \\ \hline
			Descripción del avance & #6 \\ \hline
			Porcentaje de cumplimiento & \textbf{#7} \\ \hline
			Enlace de video & \video{#8} \\ \hline
		\end{tabularx}
	\end{minipage}
	\par\vspace{1.2em}%
}

% Porcentaje de cumplimiento: real respecto a lo esperado al corte
\newcommand\cumplimiento[2]{\fpeval{min(100, round(100*(#2)/(#1), 0))}\%}

% Columna del mes reportado en el cronograma
\newcommand\mesActual[2]{\ifnum\corte=#1\cellcolor{black!70}\textcolor{white}{\textbf{#2}}\else\textbf{#2}\fi}

\begin{document}

	\begin{center}
		\fontsize{14}{16}\selectfont\bfseries
		\numReporte° Reporte de seguimiento (\mes{} 2026): \\
		Simulación de sistemas complejos mediante autómatas celulares neuronales
	\end{center}

	\begin{center}
		\fontsize{14}{16}\selectfont\bfseries
		\textit{Trabajo Terminal No. 2026-A114}
	\end{center}

	\begin{center}
		\fontsize{12}{16}\selectfont
		\textit{Alumno: \authormark{*}Angeles López Erick Jesse\authormark{1}} \\
		\textit{Director: Juárez Martínez Genaro\authormark{1}} \\
		\textit{\authormark{*}e-mail: eangelesl1700@alumno.ipn.mx}
	\end{center}

	\section{Información general}

	\begin{itemize}[leftmargin=*]
		\item \textbf{Periodo reportado:} \mes{} 2026 (corte al \fechaCorte).
		\item \textbf{Documento de evidencia adjunto:} reporte técnico de Trabajo Terminal
		\textit{``Simulación de sistemas complejos mediante autómatas celulares neuronales''} (PDF, 204 páginas).
		El avance es incremental: cada evidencia listada en este reporte se encuentra integrada en el capítulo
		y sección correspondiente de dicho documento, cuyas páginas se indican en cada actividad.
		\item \textbf{Actividades programadas en el mes:} \segun{Módulos de entrenamiento (cierre), Entrenamiento y ajuste,
		Experimentación y análisis, Análisis de resultados (global).}{Entrenamiento y ajuste, Experimentación y análisis,
		Análisis de resultados (global).}
	\end{itemize}

	\section{Resumen de cumplimiento}

	\begin{table}[H]
		\centering
		\renewcommand{\arraystretch}{1.2}
		\begin{tabular}{l c c c c}
			\toprule
			\textbf{Actividad} & \textbf{Periodo} & \textbf{Esperado al corte} & \textbf{Avance real} & \textbf{Cumplimiento} \\
			\midrule
			\multicolumn{5}{l}{\textit{Fase 1: Planificación y Revisión Bibliográfica}} \\
			Estado del arte                  & Ago/25 -- Oct/25 & 100\% & 100\% & 100\% \\
			Marco teórico y metodológico     & Sep/25 -- Oct/25 & 100\% & 100\% & 100\% \\
			\multicolumn{5}{l}{\textit{Fase 2: Diseño e implementación del simulador}} \\
			Arquitectura del simulador       & Oct/25 -- Nov/25 & 100\% & 100\% & 100\% \\
			Módulos de simulación            & Nov/25 -- Dic/25 & 100\% & 100\% & 100\% \\
			Módulos entrenamiento            & Dic/25 -- Ago/26 & 100\% & 100\% & 100\% \\
			\multicolumn{5}{l}{\textit{Fase 3: Entrenamiento y Validación}} \\
			Entrenamiento y ajuste           & Ago/26 -- Oct/26 & \espEntrenamiento\% & \realEntrenamiento\% & \cumplimiento{\espEntrenamiento}{\realEntrenamiento} \\
			Experimentación y análisis       & Ago/26 -- Nov/26 & \espExperimentos\%  & \realExperimentos\%  & \cumplimiento{\espExperimentos}{\realExperimentos} \\
			Análisis de resultados (global)  & Ago/26 -- Nov/26 & \espAnalisis\%      & \realAnalisis\%      & \cumplimiento{\espAnalisis}{\realAnalisis} \\
			\multicolumn{5}{l}{\textit{Fase 4: Documentación y Difusión}} \\
			Documentación técnica            & Oct/26 -- Dic/26 & 0\% & -- & No iniciada \\
			Artículo académico               & Oct/26 -- Dic/26 & 0\% & -- & No iniciada \\
			\bottomrule
		\end{tabular}
		\caption{Avance acumulado por actividad al \fechaCorte. ``Esperado al corte'' es la proporción del periodo
		programado en el cronograma que ha transcurrido; ``Cumplimiento'' compara el avance real con dicho valor.}
	\end{table}

	\section{Evidencias por actividad}

	\subsection*{Fase 1: Planificación y Revisión Bibliográfica}

	\actividad{Estado del arte}
	{Fase 1}{Ago/25 -- Oct/25}
	{Documento escrito.}
	{Capítulo 5, sección 5.1 \textit{Estado del arte, aplicaciones y líneas actuales de investigación} (pp.~61--77);
	Capítulo 1 \textit{Introducción} (pp.~1--6).}
	{Revisión de trabajos sobre autómatas celulares neuronales: morfogénesis, NCA isotrópicos, autoclasificación,
	representaciones tridimensionales, NCA en malla, modelos generativos y aplicaciones adicionales.
	Actividad concluida en Trabajo Terminal I.}
	{100\%}{}

	\actividad{Marco teórico y metodológico}
	{Fase 1}{Sep/25 -- Oct/25}
	{Documento escrito.}
	{Capítulos 2 a 5 (marco teórico, pp.~7--77) y Capítulo 6 \textit{Marco metodológico} (pp.~78--90).}
	{Marco teórico de sistemas complejos y teoría del caos, autómatas celulares, inteligencia artificial y redes
	neuronales, y autómatas celulares neuronales. Marco metodológico con diseño cuasiexperimental, población,
	muestra, técnicas de recolección y procesamiento de datos, y contextualización tecnológica.
	Actividad concluida en Trabajo Terminal I.}
	{100\%}{}

	\subsection*{Fase 2: Diseño e implementación del simulador}

	\actividad{Arquitectura del simulador}
	{Fase 2}{Oct/25 -- Nov/25}
	{Documento escrito (análisis y diseño de software, diagramas UML).}
	{Capítulo 7 \textit{Análisis de software} (pp.~91--135); Capítulo 8, secciones 8.1 \textit{Arquitectura},
	8.2 \textit{Diagrama de clases}, 8.3 \textit{Diseño de datos} y 8.5 \textit{Interfaces de usuarios} (pp.~136--143, 153--156).}
	{Metodología de desarrollo, reglas de negocio, requerimientos, análisis de riesgos, estudio de factibilidad,
	casos de uso con sus especificaciones, diagramas de secuencia, arquitectura MVC, diagrama de clases,
	diseño de datos e interfaces gráficas.}
	{100\%}{}

	\actividad{Módulos de simulación}
	{Fase 2}{Nov/25 -- Dic/25}
	{Documento escrito, tablas de resultados y pruebas de operatividad.}
	{Sección 8.4 \textit{Diseño de algoritmos} (8.4.1--8.4.4 y 8.4.6, pp.~144--152); secciones 9.1 \textit{Entorno de desarrollo}
	y 9.2 \textit{Estructura del proyecto} (pp.~157--162); subsecciones 9.4.1, 9.4.3, 9.4.4 y 9.4.5 de
	\textit{Validación y pruebas} (pp.~166--169), tablas 9.6 y 9.7.}
	{Núcleo de simulación en C++/CUDA con enlaces a Python e interfaz en Qt: indexación de la regla semitotalística,
	paso de evolución, condiciones de frontera, \textit{forward pass} del NCA y \textit{prefetch} asíncrono.
	Validación funcional contra reglas conocidas y pruebas de rendimiento (hasta 345\,000 pasos/s en
	$100\times100$ y más de 100 fps de renderizado en $2000\times2000$).}
	{100\%}{\videoSimulacion}

	\actividad{Módulos entrenamiento}
	{Fase 2}{Dic/25 -- Ago/26}
	{Documento escrito, pruebas de operatividad.}
	{Subsección 7.8.4 \textit{Especificaciones de casos de uso: Entrenamiento} (pp.~125--130); 8.3.5
	\textit{Configuración del NCA} (p.~143); 8.4.5 \textit{Entrenamiento del autómata celular neuronal} (pp.~150--151);
	sección 9.3 \textit{Entrenamiento del modelo} (pp.~163--166, tabla 9.5); subsecciones 9.4.2 y 9.4.6 (pp.~167, 169).}
	{Módulo de entrenamiento en PyTorch con construcción dinámica del perceptrón multicapa, ciclo de optimización
	con Adam y detención temprana, persistencia de pesos en \texttt{.npz} compatible con el \textit{kernel} CUDA e
	integración con la pestaña de entrenamiento de la GUI (subproceso con \textit{log} en tiempo real).
	Se agregó el parámetro de semilla aleatoria para reproducir experimentos.
	\segun{Actividad concluida en el mes reportado.}{Actividad concluida en agosto de 2026.}}
	{100\%}{\videoEntrenamiento}

	\subsection*{Fase 3: Entrenamiento y Validación}

	\actividad{Entrenamiento y ajuste}
	{Fase 3}{Ago/26 -- Oct/26}
	{Documento escrito y tablas de resultados.}
	{Capítulo 10, sección 10.1 \textit{Duplicación de reglas del juego de la vida}:
	\segun{tablas 10.1 y 10.2 y figura 10.3 (pp.~173--174).}
	{tablas 10.1 y 10.2 y figura 10.3 (pp.~173--174); figuras 10.5 a 10.9 (pp.~176--178).}}
	{\segun{Barrido de hiperparámetros para la regla B3/S23: 10 combinaciones de funciones de activación con
	una capa oculta de 4 neuronas (una sola converge) y con 4 capas ocultas de 8 neuronas (cinco convergen a una
	pérdida menor a $10^{-4}$). Se identifica el estancamiento en las primeras 12 épocas.}
	{Además del barrido de hiperparámetros de agosto, se entrenaron las reglas B3/S23, B3/S238, B38/S23, B3/S236
	y B38/S238 con distintas pérdidas objetivo y semillas, midiendo la distancia Manhattan entre el vector de regla
	aprendido y el original. Se determinó que la pérdida objetivo debe ser $\leq 0.1$ para B3/S23.}}
	{\realEntrenamiento\% (esperado: \espEntrenamiento\%)}{}

	\actividad{Experimentación y análisis}
	{Fase 3}{Ago/26 -- Nov/26}
	{Documento escrito, tablas y gráficas de resultados.}
	{Capítulo 10, sección 10.1:
	\segun{figuras 10.1 y 10.2 (pp.~170--172).}
	{figuras 10.1 a 10.13 (pp.~170--181).}}
	{\segun{Generación del conjunto de datos (100 condiciones iniciales en $100\times100$ durante 100 pasos) y
	análisis de su balance: el estado 1 aparece un $16.5\% \pm 4.2\%$ y los patrones $C0N8$ y $C1N8$ representan
	menos del $0.02\%$ de las entradas.}
	{Primer experimento (duplicación de reglas del juego de la vida) completo. Se diseñó un espacio con las
	512 permutaciones de la vecindad para extraer el vector de regla aprendido y verificar si el NCA conserva la
	propiedad semitotalística. Se identificaron reglas aprendidas alternativas (B3/S23 en lugar de B3/S238,
	B3/S23678 en lugar de B3/S236) y reglas no semitotalísticas.}}
	{\realExperimentos\% (esperado: \espExperimentos\%)}{}

	\actividad{Análisis de resultados (global)}
	{Fase 3}{Ago/26 -- Nov/26}
	{Documento escrito.}
	{Capítulo 10, sección 10.1\segun{ (pp.~170--172).}{, conclusiones del experimento (p.~181).}}
	{\segun{Análisis estadístico de la distribución de estados y de patrones de vecindad para la regla B3/S23,
	identificando el sesgo de muestra y proponiendo sobremuestreo o ponderación de la función de costo.}
	{Conclusiones del primer experimento: la generalización depende del sesgo de muestra de patrones poco frecuentes,
	el entrenamiento de reglas semitotalísticas no garantiza una regla aprendida semitotalística, y no se requieren
	arquitecturas complejas para replicar estas dinámicas. Se valida la metodología de extracción del vector de
	regla para los siguientes experimentos.}}
	{\realAnalisis\% (esperado: \espAnalisis\%)}{}

	\subsection*{Fase 4: Documentación y Difusión}

	\actividad{Documentación técnica}
	{Fase 4}{Oct/26 -- Dic/26}
	{No aplica en el periodo.}
	{Reporte técnico de TT completo (204 páginas).}
	{Actividad programada a partir de octubre de 2026. El reporte técnico se ha redactado de forma incremental
	conforme avanza cada actividad, por lo que su integración final partirá del documento adjunto.}
	{No iniciada (0\% esperado al corte)}{}

	\actividad{Artículo académico}
	{Fase 4}{Oct/26 -- Dic/26}
	{No aplica en el periodo.}
	{--}
	{Actividad programada a partir de octubre de 2026.}
	{No iniciada (0\% esperado al corte)}{}

	\section{Cronograma}

	\resizebox{\linewidth}{!}{%
		\begin{tabular}{p{5cm}*{10}{>{\centering\arraybackslash}p{1cm}}}
			\toprule
			\textbf{Actividad} & \textbf{Ago/25} & \textbf{Sep/25} & \textbf{Oct/25} & \textbf{Nov/25} & \textbf{Dic/25} & \mesActual{1}{Ago/26} & \mesActual{2}{Sep/26} & \textbf{Oct/26} & \textbf{Nov/26} & \textbf{Dic/26} \\
			\midrule
			\multicolumn{11}{l}{\textbf{Fase 1: Planificación y Revisión Bibliográfica}} \\

			Estado del arte       & \cellcolor{gray!30} & \cellcolor{gray!30} & \cellcolor{gray!30} &  &  &  &  &  &  &   \\

			Marco teórico y metodológico &  & \cellcolor{gray!30} & \cellcolor{gray!30} &  &  &  &  &  &  &   \\

			\midrule

			\multicolumn{11}{l}{\textbf{Fase 2: Diseño e implementación del simulador}} \\
			Arquitectura del simulador &  &  & \cellcolor{gray!30} & \cellcolor{gray!30} &  &  &  &  &  &   \\

			Módulos de simulación  &  &  &  & \cellcolor{gray!30} & \cellcolor{gray!30} &  &  &  &  &   \\

			Módulos entrenamiento   &  &  &  &  & \cellcolor{gray!30} & \cellcolor{gray!30} &  &  &  &   \\

			\midrule

			\multicolumn{11}{l}{\textbf{Fase 3: Entrenamiento y Validación}} \\
			Entrenamiento y ajuste   &  &  &  &  &  & \cellcolor{gray!30} & \cellcolor{gray!30} & \cellcolor{gray!30} &   &   \\

			Experimentación y análisis  &  &  &  &  &  & \cellcolor{gray!30}  & \cellcolor{gray!30} & \cellcolor{gray!30} & \cellcolor{gray!30} &   \\

			Análisis de resultados (global) &  &  &  &  &  & \cellcolor{gray!30} & \cellcolor{gray!30} & \cellcolor{gray!30} & \cellcolor{gray!30} &   \\

			\midrule

			\multicolumn{11}{l}{\textbf{Fase 4: Documentación y Difusión}} \\
			Documentación técnica  &  &  &  &  &  &  &  & \cellcolor{gray!30}  & \cellcolor{gray!30} & \cellcolor{gray!30}  \\

			Artículo académico  &  &  &  &  &  &  &  & \cellcolor{gray!30}  & \cellcolor{gray!30}  & \cellcolor{gray!30}  \\

			\bottomrule
		\end{tabular}
	}

	\vspace{0.5em}
	{\small La columna resaltada corresponde al mes reportado.}

	\section{Alumnos y directores}

	\vspace{1em}
	\begin{tabular}{@{} p{0.45\textwidth} p{0.45\textwidth} @{}}
		\textit{Erick Jesse Angeles López} & \textit{Dr. Genaro Juárez Martínez} \\[5em]
		Firma: \rule{6cm}{0.4pt} & Firma: \rule{6cm}{0.4pt} \\
	\end{tabular}

\end{document}
"
%  o simplemente:  ./compilar.sh
%
%  Las páginas y secciones citadas corresponden a Documentation/main.pdf
%  (reporte técnico de TT), que es el PDF que se adjunta como evidencia.
% ============================================================================
\documentclass[letterpaper,10pt,oneside]{article}
\usepackage[margin=2.54cm]{geometry}
\usepackage[spanish]{babel}
\newcommand\authormark[1]{\textsuperscript{#1}}
\usepackage[colorlinks=true,bookmarks=true,linkcolor=black,citecolor=black,urlcolor=blue]{hyperref}
\usepackage{times}
\usepackage[table]{xcolor}
\usepackage{booktabs}
\usepackage{graphicx}
\usepackage{titlesec}
\usepackage{tabularx}
\usepackage{array}
\usepackage{enumitem}
\usepackage{float}

\titleformat{\section}
{\normalfont\fontsize{12pt}{14pt}\bfseries}
{\thesection}{1em}{}

% ----------------------------------------------------------------------------
%  Mes reportado
% ----------------------------------------------------------------------------
\providecommand\corte{2}
\ifnum\corte=1
	\def\numReporte{1}
	\def\mes{Agosto}
	\def\fechaCorte{31 de agosto de 2026}
\else
	\def\numReporte{2}
	\def\mes{Septiembre}
	\def\fechaCorte{30 de septiembre de 2026}
\fi
% \segun{texto agosto}{texto septiembre}
\newcommand\segun[2]{\ifnum\corte=1 #1\else #2\fi}

% ----------------------------------------------------------------------------
%  Porcentajes (editar aquí). Formato: \segun{agosto}{septiembre}
%  "esp" = avance esperado según cronograma al corte, "real" = avance logrado.
% ----------------------------------------------------------------------------
\def\espEntrenamiento{\segun{33}{67}}   \def\realEntrenamiento{\segun{33}{67}}
\def\espExperimentos{\segun{25}{50}}    \def\realExperimentos{\segun{25}{50}}
\def\espAnalisis{\segun{25}{50}}        \def\realAnalisis{\segun{25}{50}}

% ----------------------------------------------------------------------------
%  Enlaces a video (dejar vacío si no aplica)
% ----------------------------------------------------------------------------
\def\videoSimulacion{}
\def\videoEntrenamiento{}

% ----------------------------------------------------------------------------
%  Bloque de evidencia por actividad
%  \actividad{nombre}{fase}{periodo}{tipo de evidencia}{ubicación en reporte TT}
%            {descripción}{porcentaje}{enlace de video}
% ----------------------------------------------------------------------------
\definecolor{encabezado}{gray}{0.85}
\newcommand\video[1]{\edef\tmp{#1}\ifx\tmp\empty No aplica.\else\url{#1}\fi}
\newcommand\actividad[8]{%
	\par\noindent
	\begin{minipage}{\linewidth}
		\renewcommand{\arraystretch}{1.25}
		\begin{tabularx}{\linewidth}{|>{\raggedright\bfseries\arraybackslash}p{0.24\linewidth}|X|}
			\hline
			\rowcolor{encabezado}
			Nombre de la actividad & \textbf{#1} \\ \hline
			Fase y periodo & #2 \,(#3) \\ \hline
			Evidencia de trabajo & #4 \\ \hline
			Ubicación en el reporte de TT & #5 \\ \hline
			Descripción del avance & #6 \\ \hline
			Porcentaje de cumplimiento & \textbf{#7} \\ \hline
			Enlace de video & \video{#8} \\ \hline
		\end{tabularx}
	\end{minipage}
	\par\vspace{1.2em}%
}

% Porcentaje de cumplimiento: real respecto a lo esperado al corte
\newcommand\cumplimiento[2]{\fpeval{min(100, round(100*(#2)/(#1), 0))}\%}

% Columna del mes reportado en el cronograma
\newcommand\mesActual[2]{\ifnum\corte=#1\cellcolor{black!70}\textcolor{white}{\textbf{#2}}\else\textbf{#2}\fi}

\begin{document}

	\begin{center}
		\fontsize{14}{16}\selectfont\bfseries
		\numReporte° Reporte de seguimiento (\mes{} 2026): \\
		Simulación de sistemas complejos mediante autómatas celulares neuronales
	\end{center}

	\begin{center}
		\fontsize{14}{16}\selectfont\bfseries
		\textit{Trabajo Terminal No. 2026-A114}
	\end{center}

	\begin{center}
		\fontsize{12}{16}\selectfont
		\textit{Alumno: \authormark{*}Angeles López Erick Jesse\authormark{1}} \\
		\textit{Director: Juárez Martínez Genaro\authormark{1}} \\
		\textit{\authormark{*}e-mail: eangelesl1700@alumno.ipn.mx}
	\end{center}

	\section{Información general}

	\begin{itemize}[leftmargin=*]
		\item \textbf{Periodo reportado:} \mes{} 2026 (corte al \fechaCorte).
		\item \textbf{Documento de evidencia adjunto:} reporte técnico de Trabajo Terminal
		\textit{``Simulación de sistemas complejos mediante autómatas celulares neuronales''} (PDF, 204 páginas).
		El avance es incremental: cada evidencia listada en este reporte se encuentra integrada en el capítulo
		y sección correspondiente de dicho documento, cuyas páginas se indican en cada actividad.
		\item \textbf{Actividades programadas en el mes:} \segun{Módulos de entrenamiento (cierre), Entrenamiento y ajuste,
		Experimentación y análisis, Análisis de resultados (global).}{Entrenamiento y ajuste, Experimentación y análisis,
		Análisis de resultados (global).}
	\end{itemize}

	\section{Resumen de cumplimiento}

	\begin{table}[H]
		\centering
		\renewcommand{\arraystretch}{1.2}
		\begin{tabular}{l c c c c}
			\toprule
			\textbf{Actividad} & \textbf{Periodo} & \textbf{Esperado al corte} & \textbf{Avance real} & \textbf{Cumplimiento} \\
			\midrule
			\multicolumn{5}{l}{\textit{Fase 1: Planificación y Revisión Bibliográfica}} \\
			Estado del arte                  & Ago/25 -- Oct/25 & 100\% & 100\% & 100\% \\
			Marco teórico y metodológico     & Sep/25 -- Oct/25 & 100\% & 100\% & 100\% \\
			\multicolumn{5}{l}{\textit{Fase 2: Diseño e implementación del simulador}} \\
			Arquitectura del simulador       & Oct/25 -- Nov/25 & 100\% & 100\% & 100\% \\
			Módulos de simulación            & Nov/25 -- Dic/25 & 100\% & 100\% & 100\% \\
			Módulos entrenamiento            & Dic/25 -- Ago/26 & 100\% & 100\% & 100\% \\
			\multicolumn{5}{l}{\textit{Fase 3: Entrenamiento y Validación}} \\
			Entrenamiento y ajuste           & Ago/26 -- Oct/26 & \espEntrenamiento\% & \realEntrenamiento\% & \cumplimiento{\espEntrenamiento}{\realEntrenamiento} \\
			Experimentación y análisis       & Ago/26 -- Nov/26 & \espExperimentos\%  & \realExperimentos\%  & \cumplimiento{\espExperimentos}{\realExperimentos} \\
			Análisis de resultados (global)  & Ago/26 -- Nov/26 & \espAnalisis\%      & \realAnalisis\%      & \cumplimiento{\espAnalisis}{\realAnalisis} \\
			\multicolumn{5}{l}{\textit{Fase 4: Documentación y Difusión}} \\
			Documentación técnica            & Oct/26 -- Dic/26 & 0\% & -- & No iniciada \\
			Artículo académico               & Oct/26 -- Dic/26 & 0\% & -- & No iniciada \\
			\bottomrule
		\end{tabular}
		\caption{Avance acumulado por actividad al \fechaCorte. ``Esperado al corte'' es la proporción del periodo
		programado en el cronograma que ha transcurrido; ``Cumplimiento'' compara el avance real con dicho valor.}
	\end{table}

	\section{Evidencias por actividad}

	\subsection*{Fase 1: Planificación y Revisión Bibliográfica}

	\actividad{Estado del arte}
	{Fase 1}{Ago/25 -- Oct/25}
	{Documento escrito.}
	{Capítulo 5, sección 5.1 \textit{Estado del arte, aplicaciones y líneas actuales de investigación} (pp.~61--77);
	Capítulo 1 \textit{Introducción} (pp.~1--6).}
	{Revisión de trabajos sobre autómatas celulares neuronales: morfogénesis, NCA isotrópicos, autoclasificación,
	representaciones tridimensionales, NCA en malla, modelos generativos y aplicaciones adicionales.
	Actividad concluida en Trabajo Terminal I.}
	{100\%}{}

	\actividad{Marco teórico y metodológico}
	{Fase 1}{Sep/25 -- Oct/25}
	{Documento escrito.}
	{Capítulos 2 a 5 (marco teórico, pp.~7--77) y Capítulo 6 \textit{Marco metodológico} (pp.~78--90).}
	{Marco teórico de sistemas complejos y teoría del caos, autómatas celulares, inteligencia artificial y redes
	neuronales, y autómatas celulares neuronales. Marco metodológico con diseño cuasiexperimental, población,
	muestra, técnicas de recolección y procesamiento de datos, y contextualización tecnológica.
	Actividad concluida en Trabajo Terminal I.}
	{100\%}{}

	\subsection*{Fase 2: Diseño e implementación del simulador}

	\actividad{Arquitectura del simulador}
	{Fase 2}{Oct/25 -- Nov/25}
	{Documento escrito (análisis y diseño de software, diagramas UML).}
	{Capítulo 7 \textit{Análisis de software} (pp.~91--135); Capítulo 8, secciones 8.1 \textit{Arquitectura},
	8.2 \textit{Diagrama de clases}, 8.3 \textit{Diseño de datos} y 8.5 \textit{Interfaces de usuarios} (pp.~136--143, 153--156).}
	{Metodología de desarrollo, reglas de negocio, requerimientos, análisis de riesgos, estudio de factibilidad,
	casos de uso con sus especificaciones, diagramas de secuencia, arquitectura MVC, diagrama de clases,
	diseño de datos e interfaces gráficas.}
	{100\%}{}

	\actividad{Módulos de simulación}
	{Fase 2}{Nov/25 -- Dic/25}
	{Documento escrito, tablas de resultados y pruebas de operatividad.}
	{Sección 8.4 \textit{Diseño de algoritmos} (8.4.1--8.4.4 y 8.4.6, pp.~144--152); secciones 9.1 \textit{Entorno de desarrollo}
	y 9.2 \textit{Estructura del proyecto} (pp.~157--162); subsecciones 9.4.1, 9.4.3, 9.4.4 y 9.4.5 de
	\textit{Validación y pruebas} (pp.~166--169), tablas 9.6 y 9.7.}
	{Núcleo de simulación en C++/CUDA con enlaces a Python e interfaz en Qt: indexación de la regla semitotalística,
	paso de evolución, condiciones de frontera, \textit{forward pass} del NCA y \textit{prefetch} asíncrono.
	Validación funcional contra reglas conocidas y pruebas de rendimiento (hasta 345\,000 pasos/s en
	$100\times100$ y más de 100 fps de renderizado en $2000\times2000$).}
	{100\%}{\videoSimulacion}

	\actividad{Módulos entrenamiento}
	{Fase 2}{Dic/25 -- Ago/26}
	{Documento escrito, pruebas de operatividad.}
	{Subsección 7.8.4 \textit{Especificaciones de casos de uso: Entrenamiento} (pp.~125--130); 8.3.5
	\textit{Configuración del NCA} (p.~143); 8.4.5 \textit{Entrenamiento del autómata celular neuronal} (pp.~150--151);
	sección 9.3 \textit{Entrenamiento del modelo} (pp.~163--166, tabla 9.5); subsecciones 9.4.2 y 9.4.6 (pp.~167, 169).}
	{Módulo de entrenamiento en PyTorch con construcción dinámica del perceptrón multicapa, ciclo de optimización
	con Adam y detención temprana, persistencia de pesos en \texttt{.npz} compatible con el \textit{kernel} CUDA e
	integración con la pestaña de entrenamiento de la GUI (subproceso con \textit{log} en tiempo real).
	Se agregó el parámetro de semilla aleatoria para reproducir experimentos.
	\segun{Actividad concluida en el mes reportado.}{Actividad concluida en agosto de 2026.}}
	{100\%}{\videoEntrenamiento}

	\subsection*{Fase 3: Entrenamiento y Validación}

	\actividad{Entrenamiento y ajuste}
	{Fase 3}{Ago/26 -- Oct/26}
	{Documento escrito y tablas de resultados.}
	{Capítulo 10, sección 10.1 \textit{Duplicación de reglas del juego de la vida}:
	\segun{tablas 10.1 y 10.2 y figura 10.3 (pp.~173--174).}
	{tablas 10.1 y 10.2 y figura 10.3 (pp.~173--174); figuras 10.5 a 10.9 (pp.~176--178).}}
	{\segun{Barrido de hiperparámetros para la regla B3/S23: 10 combinaciones de funciones de activación con
	una capa oculta de 4 neuronas (una sola converge) y con 4 capas ocultas de 8 neuronas (cinco convergen a una
	pérdida menor a $10^{-4}$). Se identifica el estancamiento en las primeras 12 épocas.}
	{Además del barrido de hiperparámetros de agosto, se entrenaron las reglas B3/S23, B3/S238, B38/S23, B3/S236
	y B38/S238 con distintas pérdidas objetivo y semillas, midiendo la distancia Manhattan entre el vector de regla
	aprendido y el original. Se determinó que la pérdida objetivo debe ser $\leq 0.1$ para B3/S23.}}
	{\realEntrenamiento\% (esperado: \espEntrenamiento\%)}{}

	\actividad{Experimentación y análisis}
	{Fase 3}{Ago/26 -- Nov/26}
	{Documento escrito, tablas y gráficas de resultados.}
	{Capítulo 10, sección 10.1:
	\segun{figuras 10.1 y 10.2 (pp.~170--172).}
	{figuras 10.1 a 10.13 (pp.~170--181).}}
	{\segun{Generación del conjunto de datos (100 condiciones iniciales en $100\times100$ durante 100 pasos) y
	análisis de su balance: el estado 1 aparece un $16.5\% \pm 4.2\%$ y los patrones $C0N8$ y $C1N8$ representan
	menos del $0.02\%$ de las entradas.}
	{Primer experimento (duplicación de reglas del juego de la vida) completo. Se diseñó un espacio con las
	512 permutaciones de la vecindad para extraer el vector de regla aprendido y verificar si el NCA conserva la
	propiedad semitotalística. Se identificaron reglas aprendidas alternativas (B3/S23 en lugar de B3/S238,
	B3/S23678 en lugar de B3/S236) y reglas no semitotalísticas.}}
	{\realExperimentos\% (esperado: \espExperimentos\%)}{}

	\actividad{Análisis de resultados (global)}
	{Fase 3}{Ago/26 -- Nov/26}
	{Documento escrito.}
	{Capítulo 10, sección 10.1\segun{ (pp.~170--172).}{, conclusiones del experimento (p.~181).}}
	{\segun{Análisis estadístico de la distribución de estados y de patrones de vecindad para la regla B3/S23,
	identificando el sesgo de muestra y proponiendo sobremuestreo o ponderación de la función de costo.}
	{Conclusiones del primer experimento: la generalización depende del sesgo de muestra de patrones poco frecuentes,
	el entrenamiento de reglas semitotalísticas no garantiza una regla aprendida semitotalística, y no se requieren
	arquitecturas complejas para replicar estas dinámicas. Se valida la metodología de extracción del vector de
	regla para los siguientes experimentos.}}
	{\realAnalisis\% (esperado: \espAnalisis\%)}{}

	\subsection*{Fase 4: Documentación y Difusión}

	\actividad{Documentación técnica}
	{Fase 4}{Oct/26 -- Dic/26}
	{No aplica en el periodo.}
	{Reporte técnico de TT completo (204 páginas).}
	{Actividad programada a partir de octubre de 2026. El reporte técnico se ha redactado de forma incremental
	conforme avanza cada actividad, por lo que su integración final partirá del documento adjunto.}
	{No iniciada (0\% esperado al corte)}{}

	\actividad{Artículo académico}
	{Fase 4}{Oct/26 -- Dic/26}
	{No aplica en el periodo.}
	{--}
	{Actividad programada a partir de octubre de 2026.}
	{No iniciada (0\% esperado al corte)}{}

	\section{Cronograma}

	\resizebox{\linewidth}{!}{%
		\begin{tabular}{p{5cm}*{10}{>{\centering\arraybackslash}p{1cm}}}
			\toprule
			\textbf{Actividad} & \textbf{Ago/25} & \textbf{Sep/25} & \textbf{Oct/25} & \textbf{Nov/25} & \textbf{Dic/25} & \mesActual{1}{Ago/26} & \mesActual{2}{Sep/26} & \textbf{Oct/26} & \textbf{Nov/26} & \textbf{Dic/26} \\
			\midrule
			\multicolumn{11}{l}{\textbf{Fase 1: Planificación y Revisión Bibliográfica}} \\

			Estado del arte       & \cellcolor{gray!30} & \cellcolor{gray!30} & \cellcolor{gray!30} &  &  &  &  &  &  &   \\

			Marco teórico y metodológico &  & \cellcolor{gray!30} & \cellcolor{gray!30} &  &  &  &  &  &  &   \\

			\midrule

			\multicolumn{11}{l}{\textbf{Fase 2: Diseño e implementación del simulador}} \\
			Arquitectura del simulador &  &  & \cellcolor{gray!30} & \cellcolor{gray!30} &  &  &  &  &  &   \\

			Módulos de simulación  &  &  &  & \cellcolor{gray!30} & \cellcolor{gray!30} &  &  &  &  &   \\

			Módulos entrenamiento   &  &  &  &  & \cellcolor{gray!30} & \cellcolor{gray!30} &  &  &  &   \\

			\midrule

			\multicolumn{11}{l}{\textbf{Fase 3: Entrenamiento y Validación}} \\
			Entrenamiento y ajuste   &  &  &  &  &  & \cellcolor{gray!30} & \cellcolor{gray!30} & \cellcolor{gray!30} &   &   \\

			Experimentación y análisis  &  &  &  &  &  & \cellcolor{gray!30}  & \cellcolor{gray!30} & \cellcolor{gray!30} & \cellcolor{gray!30} &   \\

			Análisis de resultados (global) &  &  &  &  &  & \cellcolor{gray!30} & \cellcolor{gray!30} & \cellcolor{gray!30} & \cellcolor{gray!30} &   \\

			\midrule

			\multicolumn{11}{l}{\textbf{Fase 4: Documentación y Difusión}} \\
			Documentación técnica  &  &  &  &  &  &  &  & \cellcolor{gray!30}  & \cellcolor{gray!30} & \cellcolor{gray!30}  \\

			Artículo académico  &  &  &  &  &  &  &  & \cellcolor{gray!30}  & \cellcolor{gray!30}  & \cellcolor{gray!30}  \\

			\bottomrule
		\end{tabular}
	}

	\vspace{0.5em}
	{\small La columna resaltada corresponde al mes reportado.}

	\section{Alumnos y directores}

	\vspace{1em}
	\begin{tabular}{@{} p{0.45\textwidth} p{0.45\textwidth} @{}}
		\textit{Erick Jesse Angeles López} & \textit{Dr. Genaro Juárez Martínez} \\[5em]
		Firma: \rule{6cm}{0.4pt} & Firma: \rule{6cm}{0.4pt} \\
	\end{tabular}

\end{document}
